\section{Recovering zeros and characters from prime data}\label{sec:rec}
\providecommand{\recsinc}{\operatorname{sinc}}
\providecommand{\recgE}{\gamma_{\mathrm{E}}}
\providecommand{\recGen}{\operatorname{Gen}}
\providecommand{\recav}{\operatorname{av}}
\providecommand{\recA}{\mathcal{A}}
\providecommand{\recartanh}{\operatorname{artanh}}

This section collects inverse theorems. An observation is made on the prime side, or on the
completed $L$-function in a region where it is given by its Euler product, and the question
is what it determines: the number of zeros below a height, the zero multiset with
multiplicities, the character, the modulus, an auxiliary twist, or a set of omitted primes.
\Cref{src:thm:dgrh} enters in two ways. It turns zero contributions into pure phases
$x^{i\gamma}$, so that Bohr means, recurrence and sharp upper limits read off individual
spectral coefficients; and it makes certain spectral weights, such as $(1/4+\gamma^2)^{-1}$
or $h/(h^2+(t-\gamma)^2)$, positive, so that coincident zeros of different functions cannot
cancel. Results marked \stU{} use no information on the location of zeros. Simplicity of
zeros and linear independence of ordinates are never assumed.

Zeros are counted with multiplicity. For a Dirichlet character $\chi$ we write $\chi^*$ for
the primitive character inducing it; the Euler factors removed or added by induction have all
their zeros on $\Re s=0$, so $\chi$ and $\chi^*$ have the same zeros in the open critical
strip. The conductor-one character gives $\zeta$. The \emph{completed function} of a
primitive $\chi$ of conductor $q$ and parity $\kappa\in\set{0,1}$, $\chi(-1)=(-1)^\kappa$, is
\begin{equation}\label{rec:eq:completion}
 \Xi_\chi(s)=\bigl(s(s-1)\bigr)^{\mathbf 1_{q=1}}\,(q/\pi)^{(s+\kappa)/2}\,
 \Gamma\bigl((s+\kappa)/2\bigr)L(s,\chi).
\end{equation}
It is entire of order one, its zeros are exactly the nontrivial zeros of $L(s,\chi)$, and
$\Xi_\chi(s)=\epsilon_\chi\overline{\Xi_\chi(1-\bar s)}$ with $\abs{\epsilon_\chi}=1$
\cite{Davenport2000Multiplicative}. For $q=1$ some sources use $\xi=\tfrac12\Xi_1$; constant
factors play no role. Under \cref{src:thm:dgrh} the zeros are $\rho=\tfrac12+i\gamma$ with
$\gamma\in\R$, and both signs of $\gamma$ are listed for a complex character.

\subsection{Three tools}\label{rec:sec:tools}
\sources{dirichlet_li_rigidity; dirichlet_inverse_problems; one_gap_dirichlet_zeros; grh_zero_count_certificate; zeta_zero_moment_bounds; dirichlet_zeros_and_abelian_fields; local_euler_recovery; dirichlet_low_zero_identification}

The first tool turns a recovered zero multiset into a character, the second turns values of
the completed function in the half-plane of absolute convergence into a positive spectral
measure, and the third makes prime values of distinct characters almost orthogonal.

\begin{lemma}[Zeros determine the character; \stU]\label{rec:lem:char}
If two primitive Dirichlet $L$-functions have the same nontrivial zeros, with
multiplicities, outside a bounded set, their characters are equal.
\end{lemma}

\begin{proof}
Let $\chi,\psi$ be the characters and $a$ the parity of $\chi$. Hadamard factorization of the
order-one entire functions \eqref{rec:eq:completion} gives $\Xi_\chi/\Xi_\psi=e^{bs+c}R(s)$ for
a rational function $R$: the common canonical factors cancel and the finitely many leftover
factors have linear exponents. Put $H(s)=L(s,\chi)/L(s,\psi)$. Forming the quotient at $s+2$
and at $s$ makes the exponential factor constant, and both gamma shifts are rational by
$\Gamma(z+1)=z\Gamma(z)$. Hence $H(s+2)/H(s)$ is rational, including its constant prefactor.
For real $\sigma\to\infty$ the Dirichlet series give $H(\sigma)=1+O(2^{-\sigma})$, and a rational
function differing from one by $O(2^{-\sigma})$ on this ray is identically one, by its Laurent
expansion at infinity. So $H$ has period $2$, and for $\Re s>1$,
$H(s)=\lim_{j\to\infty}H(s+2j)=1$. Uniqueness of the coefficients of absolutely convergent
Dirichlet series gives $\chi(n)=\psi(n)$ for every $n\ge1$, and primitivity identifies the
characters and their conductors.
\end{proof}

\begin{remark}[Reading the character directly]\label{rec:rem:readchar}
When the whole multiset is known, the character can be computed from it. For real
$\sigma>1$ the Hadamard product gives the absolutely convergent sum
$J_\chi(\sigma)=-\sum_\rho(\sigma-\rho)^{-2}=(\log\Xi_\chi)''(\sigma)$, while the Euler product
gives, with $e=\mathbf 1_{q=1}$, $\psi_1$ the trigamma function and
$b_n=\Lambda(n)\chi(n)\log n$,
\[
 J_\chi(\sigma)=-e\bigl(\sigma^{-2}+(\sigma-1)^{-2}\bigr)+\tfrac14\psi_1\bigl((\sigma+\kappa)/2\bigr)
 +\sum_{n\ge2}b_nn^{-\sigma}.
\]
From $\psi_1(z)=z^{-1}+(2z^2)^{-1}+O(z^{-3})$,
$\lim_{\sigma\to\infty}\sigma^2(J_\chi(\sigma)-1/(2\sigma))=(1-\kappa)/2-2e$; the values
$-3/2,1/2,0$ distinguish $\zeta$, even nonprincipal and odd characters. Subtracting the known
rational and trigamma terms leaves $D(\sigma)=\sum b_nn^{-\sigma}$, and in increasing order of
$n$, $b_n=\lim_{\sigma\to\infty}n^\sigma\bigl(D(\sigma)-\sum_{2\le j<n}b_jj^{-\sigma}\bigr)$, by
domination at $\sigma=2$. At a prime $\ell$ division by $(\log\ell)^2$ gives $\chi(\ell)$.
\end{remark}

Primitivity is what is identified. If $\chi$ modulo $Q$ is induced by $\chi^*$, the two
prime-power sums $\sum_{n\le x}\Lambda(n)\chi(n)(1-n/x)$ differ only at powers of primes dividing
$Q$, by at most $\omega(Q)\log x$ for $x\ge2$. So the prime-side observations of this section
recover the inducing primitive character; for example all principal characters have the
same limiting measurements as $\zeta$.

\begin{lemma}[The squared coordinate; \stG]\label{rec:lem:square}
Let $\chi$ be primitive, $F=\Xi_\chi$, $t\in\R$ and $s_t=\tfrac12+it$. The function
$w\mapsto F(s_t+w)F(s_t-w)$ is even and entire, so it equals $H_t(w^2)$ for an entire function
$H_t$ of order at most $1/2$. Its zero multiset in $z=w^2$ is $\set{-(\gamma-t)^2}_\gamma$, a
zero of order $r$ of $F$ at $s_t$ giving order $r$ at $z=0$, and equal squared distances adding
their multiplicities. For $z\notin(-\infty,0]$,
\begin{equation}\label{rec:eq:stieltjes}
 \frac{H_t'(z)}{H_t(z)}=\sum_\gamma\frac1{z+(\gamma-t)^2},
 \qquad
 \frac{H_t'(h^2)}{H_t(h^2)}=\frac1h\Re\frac{F'}{F}(s_t+h)\quad(h>0),
\end{equation}
with locally uniform convergence. If $F(s_t)\ne0$, then
$H_t(z)/H_t(0)=\prod_\gamma\bigl(1+z/(\gamma-t)^2\bigr)$. In particular, writing $\theta_h(t)$ for
a continuous argument of $F(\tfrac12+h+it)$,
\begin{equation}\label{rec:eq:phase}
 h\,\theta_h'(t)=\sum_\gamma\frac{h^2}{h^2+(t-\gamma)^2}\qquad(h>0).
\end{equation}
For $t=0$ the function $H_0(w^2)$ is a constant multiple of
$\Xi_\chi(\tfrac12+w)\Xi_{\bar\chi}(\tfrac12+w)$.
\end{lemma}

\begin{proof}
The zero count in $\abs\gamma\le T$ is $O_q(T\log(2+T))$, so $H_t$ has order at most $1/2$ and
its Hadamard product has genus zero with constant exponential factor; this gives the product
and its logarithmic derivative, and the local uniform convergence. A zero of $F$ of order $r$ at
$s_t+i(\gamma-t)$ gives zeros of order $r$ of the product at $w=\pm i(\gamma-t)$, which is order
$r$ at $z=-(\gamma-t)^2$ when $\gamma\ne t$ and order $2r$ in $w$, hence $r$ in $z$, when
$\gamma=t$. Evaluate the derivative of $H_t(w^2)$ at $w=h$: the reflection equation gives
$(F'/F)(s_t-h)=-\overline{(F'/F)(s_t+h)}$, whence $2hH_t'(h^2)/H_t(h^2)=2\Re(F'/F)(s_t+h)$. The
last formula is $\Re(F'/F)(s_t+h)=\theta_h'(t)$ by the Cauchy--Riemann equations. For $t=0$,
$\Xi_\chi(\tfrac12-w)=\epsilon_\chi\overline{\Xi_\chi(\tfrac12-\bar w)}$ and
$\overline{\Xi_\chi(\bar s)}=\Xi_{\bar\chi}(s)$.
\end{proof}

Without \cref{src:thm:dgrh} one line of data already sees the real parts of all zeros.

\begin{corollary}[What a single line records; \stU]\label{rec:cor:radius}
Write $\rho=\beta+i\gamma$ for the nontrivial zeros of $L(s,\chi)$, $\Delta=\sup_\rho\abs{\beta-1/2}$,
and let $r_h(t)$ be the radius of convergence of the Taylor series of $H_t'/H_t$ about $z=h^2$,
with $h>1/2$. All its coefficients are real derivatives of $\Re(F'/F)(\tfrac12+\sqrt z+it)/\sqrt z$
at $z=h^2$. Then $\inf_{t\in\R}r_h(t)=h^2-\Delta^2$. Under \cref{src:thm:dgrh},
$r_h(t)=h^2+\min_\gamma(t-\gamma)^2$.
\end{corollary}

\begin{proof}
Without a hypothesis on $\beta$ the poles of $H_t'/H_t$ are the points $z=(\rho-s_t)^2$, with
positive integral residues that cannot cancel, so $r_h(t)$ is the distance from $h^2$ to the
nearest of them. For $u=\beta-\tfrac12$, $v=\gamma-t$,
$\abs{h^2-(u+iv)^2}^2=(h^2-u^2)^2+2(h^2+u^2)v^2+v^4$, minimized at $t=\gamma$. Since
$\abs u\le1/2<h$, taking infima gives the first formula, and $u=0$ gives the second.
\end{proof}

The third tool converts \cref{src:thm:dgrh} into almost-orthogonality of prime values, uniformly in the
conductors.

\begin{lemma}[Prime Gram estimate; \stG]\label{rec:lem:gram}
Fix a smooth compactly supported real weight $w$ on $(0,\infty)$. For primitive characters $\chi,\xi$ of
conductors $q,f$ and $X\ge X_w$,
\[
 \sum_p(\log p)w(p/X)\chi(p)\overline{\xi(p)}=\mathbf 1_{\chi=\xi}\,X\int_0^\infty w(y)\dd y+O_w\bigl(\sqrt X\log(2qf)\bigr).
\]
\end{lemma}

\begin{proof}
Let $\eta$ be the primitive character inducing $\chi\bar\xi$ on $(\Z/\lcm(q,f)\Z)^\times$; its conductor is at most
$qf$, and by uniqueness of primitive induction it is principal exactly when $\chi=\xi$. Mellin
inversion writes $\sum_n\Lambda(n)\eta(n)w(n/X)$ as the integral of $-(L'/L)(s,\eta)X^s\widetilde w(s)$ on $\Re s=2$,
$\widetilde w(s)=\int_0^\infty w(y)y^{s-1}\dd y$. Move the line to $\Re s=-\tfrac14$, with horizontal edges at
distance $\gg1/\log(f'(3+\abs t))$ from all ordinates, where $L'/L\ll\log^2(f'(3+\abs t))$ \cite{Davenport2000Multiplicative}.
The pole of $\zeta$ gives $X\widetilde w(1)$, a zero at $s=0$ costs $O(1)$, the new line costs
$O_w(X^{-1/4}\log(2qf))$ by the functional equation and Stirling, and by \cref{src:thm:dgrh} the
nontrivial zeros give $-\sum_\rho X^\rho\widetilde w(\rho)=O_w(\sqrt X\log(2qf))$, by the rapid decay of $\widetilde w$ on
the critical line and the unit-interval zero count. Prime powers of exponent at least two cost
$O_w(\sqrt X)$ by Chebyshev's bound, and passing from $\eta(p)$ to $\chi(p)\overline{\xi(p)}$ at primes dividing
$qf$ costs $O(\log(qf))$.
\end{proof}

\paragraph{Comparison.}
Several observations of this section determine a character, or its zero multiset, under different
data. \Cref{rec:thm:universal} uses one continuous autocorrelation for all characters at once;
\cref{rec:thm:factorial} one twisted sum, sampled only at factorials; \cref{rec:thm:threescale} and
\cref{rec:prop:oriented} sharp upper limits on the grid $6^n$; \cref{rec:thm:dyadic} one complex
sequence at the powers of two; \cref{rec:thm:mellin} sparse samples of split primes, with no hypothesis
on zeros; \cref{rec:thm:liseparation} filtered and subsampled Li coefficients with sparse corruption
and no hypothesis on zeros; \cref{rec:thm:oneinterval} the positions of zeros of high derivatives on
one interval; \cref{rec:thm:localrecovery} poles in one triple-logarithmic window;
\cref{rec:thm:localcoeff,rec:thm:localid} one window of complex zeros at a large height, uniformly in the
conductor; \cref{rec:thm:operatorinverse} one prime frequency of the zeros of a differentiated
completion; and \cref{rec:thm:lowzero,rec:thm:robust} the zeros below an explicit height
$O(\log^2Q/\log\log Q)$, respectively $O(\log^2Q\log\log Q)$ with explicit constants and a clipped error
metric. None of these data determines another, and each statement is kept. By contrast, \cref{rec:thm:jumps} and
\cref{rec:thm:recovery} recover the parameters of the source construction from signed data in which
the names of the characters have disappeared, and \cref{rec:thm:localclass} shows that differences of
pooled heights determine a packet only up to local rotations.

\subsection{Counting zeros from the half-plane of absolute convergence}\label{rec:sec:count}
\sources{grh_zero_count_certificate; dirichlet_inverse_problems}

By \cref{rec:lem:square}, real values of the logarithmic derivative of $\Xi_\chi$ to the right
of $\Re s=1$, or phases on vertical lines there, are Stieltjes transforms of the zero counting
measure. Two finite certificates follow: one counts the zeros with $\abs\gamma\le T$ from real
samples, the other counts the zeros in a height interval $[A,B]$ from phases. Neither needs a
zero list or an a priori lower bound on the distance of the boundary from the zeros.

\paragraph{Real samples.}
The certificate is stated for an abstract positive discrete measure. Let $(t_j)_{j\in J}$ be a
finite or countable list in $(0,1]$, repeated according to multiplicity, with
$M=\sum_jt_j<\infty$, and for $0<\tau<1$ put $N_>(\tau)=\#\set{j:t_j>\tau}$,
$N_\ge(\tau)=\#\set{j:t_j\ge\tau}$. Fix $n\ge1$ and distinct $d_1,\dots,d_n>0$, and put
\[
 D(t)=\prod_{i=1}^n(1+d_it),\qquad d=D(1),\qquad V_i=\sum_j\frac{t_j}{1+d_it_j},\qquad
 B_{n,k}(t)=\sum_{r=k}^n\binom nrt^r(1-t)^{n-r},
\]
with $B_{n,0}=1$; at negative arguments $B_{n,k}$ is evaluated as a polynomial. For $1\le k\le n$
set
\begin{equation}\label{rec:eq:binomcoef}
 c_i^{(k)}=\frac{B_{n,k}(-1/d_i)}{(-1/d_i)\prod_{h\ne i}(1-d_h/d_i)},\qquad
 S_k=\sum_{i=1}^nc_i^{(k)}V_i .
\end{equation}

\begin{theorem}[Rational-binomial count certificate; \stU]\label{rec:thm:count}
For integers $1\le k_-\le k_+\le n$ put
\[
 L=S_{k_+}-\frac n{k_+}B_{n-1,k_+-1}(\tau)\,M,\qquad U=\frac d{B_{n,k_-}(\tau)}\,S_{k_-}.
\]
Then $L\le N_>(\tau)\le N_\ge(\tau)\le U$. If $\max(0,\ceil L)=\floor U$, this integer is the
count and there is no atom at $\tau$.
\end{theorem}

\begin{proof}
First, $B_{n,k}(t)/D(t)=\sum_ic_i^{(k)}t/(1+d_it)$. Indeed, after multiplication by $D(t)$ both
sides are polynomials of degree at most $n$ vanishing at $0$, and at each $-1/d_i$ only the
$i$th term on the right survives and equals the left side by \eqref{rec:eq:binomcoef}; they
agree at $n+1$ points. For $0\le t\le1$, $B_{n,k}(t)=\Pr(X\ge k)$ for $X\sim\mathrm{Bin}(n,t)$, so
it is increasing with values in $[0,1]$, and
$kB_{n,k}(t)\le\mathbb E[X\mathbf 1_{X\ge k}]=ntB_{n-1,k-1}(t)$ by $r\binom nr=n\binom{n-1}{r-1}$.
Hence all sums converge absolutely and $S_k=\sum_jB_{n,k}(t_j)/D(t_j)$.

For $t\le\tau$, the expression $B_{n,k_+}(t)/D(t)-(n/k_+)B_{n-1,k_+-1}(\tau)\,t$ is $\le0$ by the
moment inequality and $D\ge1$; for $t>\tau$ it is at most $1$. So it is a minorant of
$\mathbf 1_{t>\tau}$. By monotonicity and $D\le d$,
$dB_{n,k_-}(t)/(D(t)B_{n,k_-}(\tau))\ge\mathbf 1_{t\ge\tau}$, the denominator being positive. Summing
over the atoms proves the inequalities. The factor $t$ in the correction is what makes the
minorant summable over infinitely many atoms.
\end{proof}

\begin{corollary}[Dirichlet zero counts from real samples; \stG]\label{rec:cor:countreal}
Let $\chi$ be primitive, $F_\chi(x)=\sum_\gamma(x+\gamma^2)^{-1}=x^{-1/2}\Re(\Xi_\chi'/\Xi_\chi)(\tfrac12+\sqrt x)$
as in \cref{rec:lem:square} with $t=0$, and let $a>1/4$, $T>0$. Put
\[
 t_\gamma=\frac a{a+\gamma^2},\qquad \tau=\frac a{a+T^2},\qquad M=aF_\chi(a),\qquad
 V_i=aF_\chi\bigl(a(1+d_i)\bigr).
\]
Then $L$ and $U$ of \cref{rec:thm:count} bound below and above the numbers of zeros with
$\abs\gamma<T$ and with $\abs\gamma\le T$. They use only the $n+1$ real values of
$\Re\Xi_\chi'/\Xi_\chi$ at $\tfrac12+\sqrt a$ and $\tfrac12+\sqrt{a(1+d_i)}$, all in $(1,\infty)$.
Imprimitive characters are covered through their inducing characters.
\end{corollary}

\begin{proof}
This is \cref{rec:thm:count} for the atoms $t_\gamma$, since
$t_\gamma/(1+d_it_\gamma)=a/(a(1+d_i)+\gamma^2)$ and $t_\gamma>\tau$ if and only if $\abs\gamma<T$.
\end{proof}

\begin{theorem}[Arbitrarily local sampling suffices; \stG]\label{rec:thm:countlocal}
Fix $a>1/4$, $0<\tau<1$ and an interval $(s_0,s_0+\ell)$ with $s_0=\tfrac12+\sqrt a$, $\ell>0$.
There are certificates with all sample points in this interval whose bounds satisfy
$L_n\to N_>(\tau)$ and $U_n\to N_\ge(\tau)$. If no zero has $\abs\gamma=T$, a finite certificate,
evaluated with sufficiently accurate interval enclosures, gives the exact count; for rational
$a$ and $T$ this is effective.
\end{theorem}

\begin{proof}
Choose rational $c>0$ with $s_0+\ell>\tfrac12+\sqrt{a(1+c)}$, and $d_i=ci/n^3$. Then
$1\le D_n(t)\le\exp(\sum_id_i)\to1$ uniformly on $[0,1]$. Take $\delta_n=n^{-1/4}$ and
$k_\pm=\ceil{n(\tau\pm\delta_n)}$. Binomial concentration (Chebyshev's inequality, since
$n\delta_n^2\to\infty$) gives $B_{n,k_+}(t)\to\mathbf 1_{t>\tau}$, $B_{n,k_-}(t)\to\mathbf 1_{t\ge\tau}$,
$B_{n-1,k_+-1}(\tau)\to0$ and $B_{n,k_-}(\tau)\to1$. Eventually $k_\pm\ge n\tau/2$, so
$B_{n,k_\pm}(t)\le2t/\tau$, and dominated convergence with $\sum t_j<\infty$ gives the limits.
When the two counts equal an integer $N$, the endpoints approach $N$ from the correct sides,
so eventually $L>N-1$ and $U<N+1$, and approximations with propagated error tending to zero
preserve this. At a boundary zero the test need not terminate.
\end{proof}

\begin{proposition}[Gap-dependent error; \stU]\label{rec:prop:countrate}
In \cref{rec:thm:count} suppose no atom lies in $(\tau-g,\tau+g)$ with $0<g<\min(\tau,1-\tau)$.
Take $k_\pm=\ceil{n(\tau\pm g/2)}$, $n-1\ge4/g$, $E=\exp(-(n-1)g^2/8)$ and $N=N_>=N_\ge$. Then
\[
 0\le N-L\le\frac M\tau\Bigl(1-\frac{1-E}d+E\Bigr),\qquad
 0\le U-N\le\frac M\tau\Bigl(\frac d{1-E}-1+\frac{2dE}{1-E}\Bigr).
\]
With $d-1=O(E)$, for example $d_i=Ei/n^3$, the error is $O((M/\tau)E)$.
\end{proposition}

\begin{proof}
The rounding and $n-1\ge4/g$ give $(k_+-1)/(n-1)\ge\tau+g/4$, $k_+/n\le\tau+3g/4$ and
$k_-/n\le\tau-g/4$. Hoeffding's inequality \cite{Hoeffding1963Probability} gives
$B_{n-1,k_+-1}(\tau)\le E$ and $B_{n,k_-}(\tau)\ge1-E$. Each atom $t\ge\tau+g$ contributes at least
$(1-E)/d$ to $S_{k_+}$, so $S_{k_+}\ge N(1-E)/d$. For $t\le\tau-g$ the moment inequality gives
$B_{n,k_-}(t)\le(n/k_-)tE\le2tE/\tau$, so $S_{k_-}\le N+2ME/\tau$. With $n/k_+\le1/\tau$ and
$N\le M/\tau$, substitution gives both bounds.
\end{proof}

Subtracting certified counts at zero-free heights counts the zeros in an annulus
$T_1<\abs\gamma\le T_2$. Running the certificates in parallel over rational partitions of a
height range with zero-free endpoints recovers the counting measure of $\abs\gamma$ to any
resolution, since zero heights are discrete. This does not separate $\gamma$ from $-\gamma$ for
a complex character and does not decide simplicity. The samples come from the Euler product:
for real $s>1$,
\[
 \Re\frac{\Xi_\chi'}{\Xi_\chi}(s)=\tfrac12\log(q/\pi)+\tfrac12\psi\Bigl(\frac{s+\kappa}2\Bigr)
 -\sum_{m\ge2}\frac{\Lambda(m)\Re\chi(m)}{m^s}+\mathbf 1_{q=1}\Bigl(\frac1s+\frac1{s-1}\Bigr),
\]
and for integer $X\ge3$ the discarded tail is at most
$X^{1-s}\bigl(\log X/(s-1)+(s-1)^{-2}\bigr)$. Since the coefficients of the certificate are rational
when $\tau$ and the $d_i$ are, enclosures of the samples propagate to an enclosure of $[L,U]$
whose extra width is computable in advance; small node spreads make the coefficients large, so
this is an information statement, not a stability claim.

\begin{example}\label{rec:ex:count8}
For $\zeta$ take $a=64$, $n=16$, $\tau=1/2$ (so $T=8$), $k_-=7$, $k_+=10$ and $d_i=i/(4n^3)$;
all samples lie just above $s_0=8.5$. Outward-rounded interval arithmetic, with $\zeta'/\zeta$
computed from the alternating series $\eta(s)=\sum_{m\ge0}2^{-m-1}\sum_{k\le m}(-1)^k\binom mk(k+1)^{-s}$
and $\psi$ from Binet's formula with a Bernoulli remainder, gives
$-1.314658329293\le N_\zeta(8)\le0.186602311764$, so both counts are zero; the largest input
interval has width $25750\cdot10^{-100}$. Exact rational tests on synthetic spectra return the
brackets $[1.356633,2.374871]$ ($n=32$, count $2$) and $[6.177049,7.219724]$ ($n=64$, count
$7$), and, for a spectrum with an atom of multiplicity three at the threshold, the bracket
$[1.397932,5.380918]$ with strict count $2$ and inclusive count $5$, which retains the jump. The
zeta count is classical; the example tests the certificate. The computation, its serialized
certificate, its verifier and the synthetic tests are
\begin{center}
\artifact{grh_zero_count_certificate/certified_zeta.py}\\
\artifact{grh_zero_count_certificate/zeta_certificate.json}\\
\artifact{grh_zero_count_certificate/verify_certificate.py}\\
\artifact{grh_zero_count_certificate/zero_count.py}
\end{center}
\end{example}

\paragraph{Phases on vertical lines.}
Now fix $1/2<a<b$ and a finite set $S\subset[a,b]$ of distinct offsets. For nonempty
$U\subseteq S$ put
\begin{equation}\label{rec:eq:phaseweights}
 g_U(x)=\prod_{h\in U}\frac{h^2}{x^2+h^2},\qquad
 w_h(U)=\prod_{k\in U\setminus\set h}\frac{k^2}{k^2-h^2},\qquad
 Z_U=\pi\sum_{h\in U}hw_h(U),
\end{equation}
and, with $\theta_h$ as in \eqref{rec:eq:phase} for a primitive $\chi$,
\[
 P_U(t)=\sum_{h\in U}w_h(U)h\theta_h'(t),\qquad
 C_S(A,B)=\frac1{Z_S}\sum_{h\in S}w_h(S)h\bigl(\theta_h(B)-\theta_h(A)\bigr).
\]
Let $N[A,B]$ count the zeros with $\gamma\in[A,B]$, keeping the signs of the ordinates of a
complex character.

\begin{theorem}[Counting zeros from phases; \stG]\label{rec:thm:phasecount}
For a partition $S=I\sqcup J$ into nonempty parts put
$E_S(A,B)=\frac{Z_J}{2Z_S}\bigl(P_I(A)+P_I(B)\bigr)$. Then $Z_S>0$, $E_S\ge0$ and
\[
 C_S(A,B)-E_S(A,B)\le N[A,B]\le C_S(A,B)+E_S(A,B).
\]
If neither endpoint is a zero ordinate and $m=\abs S$ distinct offsets in $[a,b]$ are split
into parts of sizes $\floor{m/2}$ and $\ceil{m/2}$, the error bound tends to zero exponentially
in $m$; in particular the test eventually returns the exact count.
\end{theorem}

\begin{proof}
By \eqref{rec:eq:phase}, $P_U(t)=\sum_\gamma g_U(t-\gamma)$, because
$g_U(x)=\sum_{h\in U}w_h(U)h^2/(x^2+h^2)$ (both sides are rational in $x^2$ with the same
residues and vanish at infinity); integration gives $Z_U=\int_\R g_U$. For $d\ge0$ put
$Q_S(d)=Z_S^{-1}\int_d^\infty g_S$. All factors decrease on $[0,\infty)$, so for $d,y\ge0$,
$g_S(d+y)\le g_I(d)g_J(y)$ and $Q_S(d)\le\frac{Z_J}{2Z_S}g_I(d)$. A single ordinate $\gamma$
contributes to $C_S$ the mass of $[A,B]$ under the probability density $g_S(t-\gamma)/Z_S$; its
error is at most $Q_S(\abs{A-\gamma})+Q_S(\abs{B-\gamma})$, which covers $\gamma\in\set{A,B}$ with
$Q_S(0)=1/2$. Summing over ordinates gives the inequalities.

With $D_k=\int_\R(1+u^2)^{-k}\dd u\sim\sqrt{\pi/k}$, comparison gives
$aD_{\abs U}\le Z_U\le bD_{\abs U}$, so $Z_J/Z_S$ is bounded for the balanced split. At an
endpoint $e$ at distance $d_e>0$ from the zeros,
$P_I(e)\le(1+d_e^2/b^2)^{-(\abs I-1)}\sum_\gamma b^2/(b^2+(e-\gamma)^2)$, and the last sum is finite.
The computed bound $E_S$ needs no value of $d_e$.
\end{proof}

A continuous logarithm of $L(\tfrac12+h+it,\chi)$ for $h>1/2$ is given by
$\sum_p\sum_k\chi(p)^k/(kp^{ks})$, the gamma logarithm is normalized on the positive axis, and the
factors $s,s-1$ of $\Xi_1$ have continuous logarithms there; so $\theta_h(B)-\theta_h(A)$ is a
difference of specified logarithms, and a principal argument of a ratio could lose whole turns.
A common absolute error $\eta$ in the scaled slopes $h\theta_h'(t)$ gives a worst-case error
$\eta\sum_h\abs{w_h(S)}$ in $P_S$, and its least possible size is as follows.

\begin{theorem}[Least amplification; \stU]\label{rec:thm:minimax}
Among all choices of $m\ge2$ distinct offsets in $[a,b]$,
\[
 \min\sum_{h}\abs{w_h(S)}=T_{m-1}\Bigl(\frac{b^2+a^2}{b^2-a^2}\Bigr)
 =\frac12\Bigl[\Bigl(\frac{b+a}{b-a}\Bigr)^{m-1}+\Bigl(\frac{b-a}{b+a}\Bigr)^{m-1}\Bigr],
\]
where $T_n(\cos\vartheta)=\cos n\vartheta$, attained when the $h^2$ are the images of the $m$
extrema of $T_{m-1}$ in $[a^2,b^2]$.
\end{theorem}

\begin{proof}
The $w_h$ are the Lagrange weights for evaluation at $0$ from the nodes $h^2$. Applying the
interpolation formula to $T_{m-1}$ composed with the affine map of $[a^2,b^2]$ onto $[-1,1]$,
whose values at the nodes are at most one in modulus, gives the lower bound. At the ordered
extrema the polynomial alternates in sign, and so do the weights because $0$ lies to the left
of all nodes; all terms then have the same sign, and the formula for $T_n$ outside $[-1,1]$
gives the value.
\end{proof}

The same data certify counts without \cref{src:thm:dgrh}, using only the classical strip
$\abs{\beta-1/2}\le D=1/2$; now $N[A,B]$ counts zeros of any real part with $\gamma\in[A,B]$.

\begin{theorem}[A count from the classical strip; \stU]\label{rec:thm:strip}
Put $K_S=\prod_{h\in S}(1-D^2/h^2)^{-1}$, $\alpha_I=\sum_{h\in I}\arcsin(D/h)$ and
$L_I=\cos(\alpha_I)\prod_{h\in I}(1+D^2/h^2)^{-1}$. If $\alpha_I<\pi/2$, then $P_I\ge0$ and
$\abs{N[A,B]-C_S(A,B)}\le(K_S/L_I)E_S(A,B)$.
\end{theorem}

\begin{proof}
The product $H_t$ exists without the critical-line hypothesis, and with $\delta=\beta-1/2$ the
partial fractions give $P_U(t)=\sum_\rho\Re g_U(t-\gamma+i\delta)$. For real $x$,
$\abs{h^2+(x+i\delta)^2}$ lies between $(1-D^2/h^2)(h^2+x^2)$ and $h^2+x^2+D^2$, and its argument is
at most $\arcsin(D/h)$ in absolute value. Hence $\abs{g_S(x+i\delta)}\le K_Sg_S(x)$ and
$\Re g_I(x+i\delta)\ge L_Ig_I(x)$. The integral of $g_S(x+i\delta)$ over a horizontal line is $Z_S$
by a contour shift, since the poles have imaginary parts $\pm h$ with $h>D$. The tails for one
zero are therefore at most $K_S$ times the real tails in the proof of \cref{rec:thm:phasecount},
and $\sum_\rho g_I(e-\gamma)\le P_I(e)/L_I$.
\end{proof}

\begin{example}\label{rec:ex:phase1018}
For $\zeta$, $[A,B]=[10,18]$, $I=\set{5/2,7/2,9/2}$ and $J=\set{3,4}$, so all observed real parts
lie in $[3,5]$, enclosed evaluations give $1.00585<C_S<1.00634$ and $0.28792<E_S<0.28922$, using
prime powers up to $256$ with explicit tails and bounded Stirling remainders. The enclosure lies
in $(0.7166,1.2956)$, so exactly one zero has height in $[10,18]$. For these five lines
$K_S/L_I<1.347$ and the enclosure of \cref{rec:thm:strip} still lies in $(0,2)$: the count holds
without RH, and the zero is simple and on the line, since reflection would pair an off-line zero
with a second one at the same height. See \artifact{dirichlet_inverse_problems/phase_counts/phase_check.py}
and \artifact{dirichlet_inverse_problems/phase_counts/phase_output.txt}.
\end{example}

The two certificates use different data: \cref{rec:cor:countreal} uses real points only and
counts by absolute height; \cref{rec:thm:phasecount} uses phases along lines and separates the
signs of the ordinates of a complex character. For arbitrarily narrow endpoint clearances the
eventual success of either uses \cref{src:thm:dgrh}. The residual certificate
\cref{li:thm:residual} counts the atoms of a spectrum in prescribed intervals from consecutive
power moments $p_K,\dots,p_{K+2m}$, that is, from Taylor coefficients at one point; the
certificate of \cref{rec:thm:count} instead uses finitely many values of the Stieltjes transform
itself, with no derivative, and the two data sets do not determine each other.

\subsection{Zeros from local Taylor data of the completed function}\label{rec:sec:taylor}
\sources{one_gap_dirichlet_zeros; zeta_zero_moment_bounds; grh_periods_sampling_and_zeros}

Two inverse problems use only derivatives of $\Xi_\chi$ at one point or on one zero-free
interval of the critical line. Under \cref{src:thm:dgrh} both become problems about positive
masses on a line. The first reads the zeros from the location of one zero of each high odd
derivative of $\log\abs{\Xi_\chi}$; the second reads them from power sums of inverse squared
distances. A third, unconditional, counts zeros in a disk from Hankel matrices of Taylor coefficients at its centre.

\paragraph{One zero-free interval.}
Let $F(t)=\Xi_\chi(\tfrac12+it)$ and let $a<b$ be such that $F$ has no zero with ordinate in
$[a,b]$. Put $G(t)=(t-a)^2(t-b)F(t)$.

\begin{theorem}[One-interval determination; \stG]\label{rec:thm:oneinterval}
For every odd $p\ge3$ there is exactly one $x_p\in(a,b)$ with
$\frac{\dd^p}{\dd t^p}\log\abs{G(t)}=0$ at $t=x_p$. The interval and the values $x_p$ along any
unbounded set of odd orders $p$ determine the multiset of nontrivial zeros of $L(s,\chi)$ and
the character $\chi$.
\end{theorem}

The proof is a real-variable inverse theorem. Write $a=c-d$, $b=c+d$, and let $\mu$ be a locally
finite multiset of real points $\gamma\notin[a,b]$ with weights $m_\gamma>0$ and
\begin{equation}\label{rec:eq:onegapsum}
 \sum_\gamma\frac{m_\gamma}{1+\gamma^2}<\infty .
\end{equation}
Put weights $A,B>0$ at $a,b$ with $A\ne B$ and, for odd $p\ge3$,
\[
 F_p(x)=\frac A{(x-a)^p}+\frac B{(x-b)^p}+\sum_\gamma\frac{m_\gamma}{(x-\gamma)^p}.
\]
On $(a,b)$ every term is decreasing and the endpoint limits are $\pm\infty$, so $F_p$ has exactly
one zero $x_p$. With the endpoints alone the zero is $y_p=c+d\tanh(h/p)$, $h=\tfrac12\log(A/B)$.

\begin{lemma}[Endpoint approximation; \stU]\label{rec:lem:endpoint}
If $R_0=\min_\gamma\abs{\gamma-c}$, then $x_p=y_p+O\bigl(\tfrac dp(d/R_0)^p\bigr)$; without exterior
mass $x_p=y_p$.
\end{lemma}

\begin{proof}
For fixed $C$ and $S>d$, uniformly on $\abs{x-c}\le C/p$,
\begin{equation}\label{rec:eq:onegaptail}
 \sum_{\abs{\gamma-c}\ge S}\frac{m_\gamma}{\abs{x-\gamma}^p}=O(S^{-p}),
\end{equation}
since replacing $\abs{x-\gamma}^{-p}$ by $\abs{c-\gamma}^{-p}$ costs a factor at most $e^{2C/S}$ for
large $p$, and $\sum m_\gamma(S/\abs{c-\gamma})^2<\infty$. Apply this at $y_p$ with $S=R_0$. On the
gap the nearer endpoint gives $\abs{F_p'}\ge p\min(A,B)d^{-p-1}$; the mean value theorem concludes.
\end{proof}

Suppose the points at smaller radii about $c$ are known, and let $E_p$ be minus their
contribution and that of the endpoints at $x_p$, so that $E_p$ is the unknown exterior
contribution. Let $R>d$ be its first occupied radius, with weights $L,U$ at $c-R,c+R$ (zero for
an empty side), $S>R$ the next radius and $\eta=R/S$. With $u_p=(x_p-c)/d$ put
\[
 k=d/R,\qquad \theta_p=p\,\recartanh(ku_p),\qquad G_p=(1-k^2u_p^2)^{-p/2},\qquad\theta_0=kh .
\]
Then $(1\pm ku_p)^{-p}=G_pe^{\mp\theta_p}$, and \eqref{rec:eq:onegaptail} with
\cref{rec:lem:endpoint} give
\begin{equation}\label{rec:eq:shell}
 \frac{R^pE_p}{G_p}=Le^{-\theta_p}-Ue^{\theta_p}+O(\eta^p),
\end{equation}
\[
 \theta_p=kh+\frac v{p^2}+O(p^{-4}),\qquad v=\frac{(k^3-k)h^3}3\ne0,\qquad G_p=1+O(p^{-1}).
\]

\begin{theorem}[Explicit inversion; \stU]\label{rec:thm:onegapinverse}
Along any unbounded set of odd orders, every occupied radius and both its weights are given by
\[
 R=\lim\abs{E_p}^{-1/p},\qquad H=\lim Z_p,\quad Z_p=\frac{R^pE_p}{G_p},\qquad
 J=\lim\frac{Z_p-H}{\theta_p-\theta_0},
\]
\[
 L=\tfrac12e^{\theta_0}(H-J),\qquad U=-\tfrac12e^{-\theta_0}(H+J).
\]
\end{theorem}

\begin{proof}
Let $f(t)=Le^{-t}-Ue^t$, so $f(\theta_0)=H$ and $f'(\theta_0)=-(Le^{-\theta_0}+Ue^{\theta_0})<0$. If
$H\ne0$ the right side of \eqref{rec:eq:shell} tends to $H$; if $H=0$ it is
$f'(\theta_0)v/p^2+O(p^{-4})+O(\eta^p)$ with nonzero leading coefficient. Either way $E_p\ne0$
eventually, which gives $R$, then $J=f'(\theta_0)$, and two linear equations give $L,U$. Remove
these weights and repeat; local finiteness leaves finitely many radii below any given one, and
an identically vanishing residual characterizes an empty remainder.
\end{proof}

\begin{theorem}[Resolution of the first discrepancy; \stU]\label{rec:thm:onegaprate}
Let $\mu,\widetilde\mu$ satisfy \eqref{rec:eq:onegapsum} with the same interval and endpoint
weights, let $R$ be the nearest radius where they differ, with weight differences
$\Delta L,\Delta U$, and put $H=\Delta Le^{-\theta_0}-\Delta Ue^{\theta_0}$,
$J=-\Delta Le^{-\theta_0}-\Delta Ue^{\theta_0}$. Then
\[
 x_p-\widetilde x_p\sim\frac d{2\sqrt{AB}}\Bigl(\frac dR\Bigr)^p\times
 \begin{cases}H/p,&H\ne0,\\ Jv/p^3,&H=0,\end{cases}
\]
and $\lim\abs{x_p-\widetilde x_p}^{1/p}=d/R$.
\end{theorem}

\begin{proof}
Evaluate $F_p-\widetilde F_p$ at $\widetilde x_p$; \eqref{rec:eq:shell} applies to the signed
difference, which is asymptotic to $R^{-p}H$, or to $R^{-p}Jv/p^2$ when $H=0$ (then $J\ne0$). For
$w$ between the two observed points $p(w-c)/d\to h$, the endpoint terms of $-d^{p+1}F_p'(w)/p$ tend
to $Ae^{-h}$ and $Be^h$, of sum $2\sqrt{AB}$, and the rest tends to zero. The mean value theorem
concludes.
\end{proof}

\begin{proof}[Proof of \cref{rec:thm:oneinterval}]
By \cref{src:thm:dgrh} the zeros of $F$ are real, and \eqref{rec:eq:onegapsum} holds because $F$ has
order one. In the Hadamard product of $F$ the exponential factors have affine real logarithms on
$\R$, so after two differentiations, termwise on the gap,
$\frac1{(p-1)!}\frac{\dd^p}{\dd t^p}\log\abs{G(t)}=F_p(t)$ with $A=2$, $B=1$; a zero at $t=0$
contributes the same power. \Cref{rec:thm:onegapinverse} recovers every zero with multiplicity,
and \cref{rec:lem:char} (or \cref{rec:rem:readchar}) the character.
\end{proof}

For $\Xi=\Xi_1(\tfrac12+it)$ (up to a constant) the interval $[-1,1]$ can be fixed in advance:
from $\Xi(t)=\tfrac12-(t^2+\tfrac14)\int_1^\infty P(v)v^{-3/4}\cos(\tfrac t2\log v)\,\dd v$ with
$P(v)=\sum_{n\ge1}e^{-\pi n^2v}\le e^{-\pi v}/(1-e^{-3\pi})$, the integral term is below $1/2$ on
$[-1,1]$. With $h=(\log2)/2$, every exterior pair $\pm r$ contributes negatively at $y_p$, so
$0<x_p<y_p<h/p$. After the nearest pairs are recovered, the signed contribution $K_p$ of the
markers and the recovered pairs at $x_p$ equals
$\sum_{r\ge R}m_r[(r-x_p)^{-p}-(r+x_p)^{-p}]>0$, where $R$ is the next positive ordinate, and
\eqref{rec:eq:shell} simplifies to $R=\lim K_p^{-1/p}$ and $m_R=\lim R^pK_p/(2\sinh(h/R))$.
For the first pair $K_p=2(1+x_p)^{-p}-(1-x_p)^{-p}$ uses only the observed point.

\begin{theorem}[Finite two-order bounds for $\Xi$; \stG]\label{rec:thm:onegapbounds}
For $r>z>0$ let $A_j(r,z)=\frac{(r-z)^{-j}-(r+z)^{-j}}{2jz}=\tfrac12\int_{-1}^1(r+vz)^{-j-1}\dd v$.
With $x=x_p$, $y=x_{p+2}$ and $D_j=K_j/(2jx_j)$, using the same recovered pairs at both orders, let
$\ell_p>x$ and $u_p^*\ge1$ solve
\[
 A_p(\ell_p,x)=D_p,\qquad Q_p(u_p^*)=D_{p+2}/D_p,\qquad Q_p(r)=\frac{A_{p+2}(r,y)}{A_p(r,x)} .
\]
Both solutions exist and are unique, and $\ell_p\le R\le u_p^*$ and $m_R\le D_p/A_p(u_p^*,x)$.
\end{theorem}

\begin{proof}
$A_p(\cdot,x)$ decreases continuously from $\infty$ to $0$, and $D_p=\sum_{r\ge R}m_rA_p(r,x)\ge A_p(R,x)$
since $m_R\ge1$. The logarithmic derivative of $A_j$ in $r$ lies between $-(j+1)/(r-z)$ and
$-(j+1)/(r+z)$, so $(\log Q_p)'(r)\le-(p+3)/(r+y)+(p+1)/(r-x)<0$ for $r\ge1$, because
$(p+3)x+(p+1)y<3h<2\le2r$. As $D_{p+2}/D_p$ is a positive weighted average of $Q_p(r)$ over $r\ge R$,
it is at most $Q_p(R)$; and $Q_p\to0$ at infinity. Finally
$D_p\ge m_RA_p(R,x)\ge m_RA_p(u_p^*,x)$.
\end{proof}

If $S>R$ is the next distinct ordinate, $u_p^*-R=O((R/S)^p)$, and if $m_R=1$ also
$R-\ell_p=O((R/S)^p/p)$; a multiplicity bound below $2$ certifies simplicity.

\begin{example}\label{rec:ex:onegap}
Orders $9$ and $11$, evaluated by interval arithmetic from the theta integral with no zero
supplied, give $14.10<\gamma_1<14.23$ and $m_{\gamma_1}=1$; orders $65$ and $67$ give
$14.1347251417337<\gamma_1<14.1347251417510$. Removing the first pair and using its locations as
markers, orders $127,129$ bound the second ordinate in $[21.02203963874,21.02203963945]$
(floating point). The programs and outputs are
\begin{center}
\artifact{one_gap_dirichlet_zeros/examples/xi_certificate.py}\\
\artifact{one_gap_dirichlet_zeros/examples/xi_certificate.txt}\\
\artifact{one_gap_dirichlet_zeros/examples/xi_certificate_high.txt}\\
\artifact{one_gap_dirichlet_zeros/examples/xi_demo.txt}
\end{center}
An instance of the cancelling case of
\cref{rec:thm:onegaprate}: on $(a,b)=(-1,1)$ with $(A,B)=(4,1)$, compare exterior weights $2,1$
at $-2,2$ with no exterior mass. Then $H=0$, $J=-2\sqrt2$, $v=-(\log2)^3/8$ and
$p^32^p(x_p-y_p)\to\sqrt2(\log2)^3/16=0.0294354987\ldots$; dyadic bisection gives the enclosures
$[0.02900627349,0.02900627350]$ at $p=15$ and $[0.02939267493,0.02939267494]$ at $p=127$
(\artifact{one_gap_dirichlet_zeros/examples/resonance_check.txt}).
\end{example}

\paragraph{Power sums of inverse squared distances.}
Let $a=x_1=\dots=x_r>x_{r+1}\ge x_{r+2}\ge\dots>0$ with $\sum x_i<\infty$, listed with
multiplicity, and $p_k=\sum_ix_i^k$. Put $h_k=\floor{p_k^2/p_{2k}}$, $c_k=p_k/h_k$, and
\[
 P_{d,k}(t)=[u^d]\frac{\prod_i(1+x_i^ku)}{(1+tu)^{h_k}}
 =\sum_{j=0}^d(-1)^j\binom{h_k+j-1}je_{d-j,k}\,t^j,
\]
where $e_{j,k}$ are the elementary symmetric functions of the $x_i^k$, computed from the $p_{k\ell}$
by Newton's identities. Let $R_{d,k}$ be the real root of $P_{d,k}$ nearest $c_k$ when it exists
and is unique. Write $q_j=x_{r+j}/a$, $b_d=\prod_{j\le d}q_j$, and let $C_d$ be the number of
$d$-element sets of tail indices whose normalized product is $b_d$ (if $v$ tail entries exceed
$q_d$ and $w$ equal it, $C_d=\binom w{d-v}$).

\begin{theorem}[Alternating errors with multiplicities; \stU]\label{rec:thm:alternating}
If the tail has at least $d$ entries, then for all large $k$, $h_k=r$, $R_{d,k}$ exists uniquely,
and
\[
 \frac{R_{d,k}}{a^k}-1=(-1)^{d+1}\frac{C_d}r\,b_d^k\bigl(1+o(1)\bigr).
\]
So even orders give eventual lower bounds and odd orders eventual upper bounds for $a^k$.
\end{theorem}

\begin{proof}
Scale to $a=1$. With $B_k=\sum_jq_j^k$ and $D_k=\sum_jq_j^{2k}$, both tending to zero,
$(r+B_k)^2-r(r+D_k)=r(2B_k-D_k)+B_k^2\ge0$, so $p_k^2/p_{2k}\downarrow r$ and $h_k=r$ eventually.
The coefficient of $u^j$ in $\prod(1+q_i^ku)$ is at most $B_k^j/j!$, so $P_{d,k}$ and its derivatives
converge coefficientwise to those of $H_d(t)=[u^d](1+u)^r(1+tu)^{-r}$, with $H_d(1)=0$ and
$H_d'(1)=(-1)^dr$. Hence $P_{d,k}$ has exactly one root near $1$, which tends to $1$, and $c_k\to1$
selects it. At $t=1$ the leading factor cancels:
$P_{d,k}(1)=\sum_{i_1<\dots<i_d}(q_{i_1}\cdots q_{i_d})^k=C_db_d^k(1+o(1))$, by domination. The mean
value theorem gives $R_{d,k}-1=-P_{d,k}(1)/P_{d,k}'(\xi_k)$ with $P_{d,k}'(\xi_k)\to(-1)^dr$.
\end{proof}

\begin{corollary}[Recovery from successive differences; \stU]\label{rec:cor:successive}
For an infinite list put $L_0=\lim p_k^{1/k}$ and $L_d=\lim\abs{R_{d,k}-R_{d+1,k}}^{1/k}$. Then
$L_0=a$, $L_d=ab_d$ and $x_{r+d}=L_0L_d/L_{d-1}$. Moreover, with $G_{0,k}=p_k$,
$G_{d,k}=\abs{R_{d,k}-R_{d+1,k}}$ and $D_{d,k}=G_{d,k+1}/G_{d,k}$, one has
$D_{0,k}D_{d,k}/D_{d-1,k}=x_{r+d}(1+O(\eta_d^k))$ for some $\eta_d<1$, which removes the tie factor
$C_d/r$.
\end{corollary}

\begin{proof}
The relative error at order $d+1$ is $o(b_d^k)$, so $G_{d,k}=(C_d/r)(ab_d)^k(1+o(1))$; take roots
and ratios. For the rate, let $\rho_j<1$ be the ratio of the next distinct $j$-fold tail product to
the largest ($\rho_j=0$ if there is none). The nonmaximal products contribute $O(\rho_j^{k-1})$, the
coefficients differ from their limits by $O(q_1^k)$, and subtraction of the next order costs
$O(q_{j+1}^k)$; hence $G_{j,k}=(C_j/r)(ab_j)^k\{1+O(\max(q_1,\rho_j)^k)\}$, and consecutive strides
give $\eta_d=\max(q_1,\rho_d,\rho_{d-1})$.
\end{proof}

\begin{theorem}[The sharp endpoint; \stU]\label{rec:thm:sharpendpoint}
Fix $d\ge2$ and suppose the tail has at least $d$ entries with $q_1>q_2>\dots>q_{d-1}>q_d$. Among
all finite or countable positive summable multisets with the moments $p_k,p_{2k},\dots,p_{dk}$ and
with largest entry of multiplicity $r$, for large $k$ the largest possible $k$th power of the
largest entry is $R_{d,k}$ when $d$ is odd, and the smallest possible is $R_{d,k}$ when $d$ is even.
Exactly one multiset attains it; it has $r$ equal leading entries and $d-1$ distinct smaller ones.
Among infinite multisets the same value is the supremum, respectively infimum.
\end{theorem}

\begin{proof}
Scale to $a=1$. A competitor's powered maximum $T$ satisfies $p_{2k}/p_k\le T\le p_k/r$, and both
bounds tend to one. $P_{d,k}$ is common to all competitors, and its value at $T$ is the degree-$d$
elementary symmetric function of the other powered entries, hence nonnegative. On a fixed
neighbourhood of $1$ containing $R_{d,k}$ and the bracket, $(-1)^dP_{d,k}'>0$ for large $k$, which
gives the inequality. For attainment let $R=R_{d,k}$, $m=d-1$,
$B_j(t)=[u^j]\prod_i(1+x_i^ku)(1+tu)^{-r}$ and $H(u)=\sum_{j\le m}B_j(R)u^j$. Since $B_d(R)=0$,
$(1+Ru)^rH(u)$ agrees with the original product through degree $d$. By separation $B_j(1)\sim b_j^k$
and $B_j(R)\sim b_j^k$ because $\abs{R-1}=O(b_d^k)$. In
$V(z)=\sum_{j\le m}(-1)^jB_j(R)z^{m-j}$, at $z=cq_i^k$ the terms of indices $i-1,i$ dominate and the
normalized value tends to $1-1/c$; so $V$ has $m$ distinct roots $z_i\sim q_i^k$, all below $R$, and
$H=\prod(1+z_iu)$. The powered multiset of $r$ copies of $R$ and the $z_i$ has the required moments.
At equality the degree-$d$ symmetric sum of a competitor's tail vanishes, so the tail has at most
$d-1$ entries and its product is $H$. For infinite competitors, append entries $\eps2^{-n}$ and
restore the first $d$ power sums by the inverse function theorem, the Jacobian being a nonsingular
Vandermonde-type matrix in the $d$ distinct powered locations.
\end{proof}

The Hankel-pencil test \cref{li:thm:ritzcert} uses all consecutive moments $p_1,\dots,p_{2N+1}$
and certifies simplicity and an enclosure of the $j$th atom as $N$ grows. Here the degree $d$ is
fixed and the stride $k$ grows, all selected roots estimate the same leading atom, the errors
encode the further atoms, and the bound is the exact endpoint for the given moments; neither
statement implies the other. The separation hypothesis is needed for the extremal statement, though not for
\cref{rec:thm:alternating}. For $(1,1,q,q,q)$ with $d=3$ and $w=q^k$, at $R_{3,k}=1+O(w^3)$ the
residual sums satisfy $S=3w+O(w^3)$, $Q=3w^2+O(w^3)$ for the entries after the two largest, while
$\sum y_i^3\ge Q^2/S$ gives $e_3(y)\ge(S^2-Q)(S^2-2Q)/(6S)>0$; so even the supremum over competitors
with leading multiplicity two is strictly below the order-three bound.

\paragraph{Application to Dirichlet zeros.}
For $F(z)=\Xi_\chi(\tfrac12+iz)$ and a real $\tau$ with $F(\tau)\ne0$, \cref{rec:lem:square} (with
$z\mapsto-v$) gives
\[
 \frac{F(\tau+\sqrt v)F(\tau-\sqrt v)}{F(\tau)^2}=\prod_\gamma\Bigl(1-\frac v{(\gamma-\tau)^2}\Bigr),
\]
so the Taylor coefficients of $\log$ of the left side at $v=0$ give every power sum of
$x_\gamma=\abs{\gamma-\tau}^{-2}$, and \cref{rec:thm:alternating,rec:thm:sharpendpoint,rec:cor:successive}
apply under \cref{src:thm:dgrh}. A multiplicity at distance $D$ counts the zeros at $\tau\pm D$
together. For the positive zeta ordinates one may use instead
$\xi(1+u)/\xi(1)=\prod_{\gamma>0}\bigl(1+(u+u^2)/(1/4+\gamma^2)\bigr)$.

\begin{example}\label{rec:ex:fourmoments}
In the coordinate $y=200/(1/4+\gamma^2)$ for the positive zeta ordinates, enclosures of
$M_4,M_8,M_{12},M_{16}$ (from the Taylor series of $\log\xi$ at $1$, regenerated from the generator
of the packet \texttt{li\_ritz\_zero\_recovery}) place $A=\max y^4$ in $[0.94,1.18]$, show
$M_4^2/M_8<2$, so $\gamma_1$ is simple, and certify $P_3'<0$, $P_4'>0$ there. Exact signs of $P_3$ and
$P_4$ at two test points give $14.13472107<\gamma_1<14.13472515$, and the residual polynomials
certify that the order-three and order-four roots are the exact extremal values for these moments.
The rational example $(1,1,\tfrac12,\tfrac12,\tfrac12,\tfrac13,\tfrac14)$, where the normalized
errors of \cref{rec:thm:alternating} visibly tend to one for $d\le4$, is checked in exact
arithmetic. The files are
\begin{center}
\artifact{zeta_zero_moment_bounds/verification/zeta_finite.py}\\
\artifact{zeta_zero_moment_bounds/verification/zeta_finite_output.txt}\\
\artifact{zeta_zero_moment_bounds/verification/check_rational.py}
\end{center}
\end{example}

\paragraph{Zeros in a disk from the Taylor expansion at its centre.}
The previous two methods use the reality of the zeros. The next one counts distinct real zeros and
conjugate pairs of a real analytic function in a disk, and recovers the multiplicities and positions
of the zeros on its boundary, from high Taylor coefficients of $\log F$ at the centre alone, with no
hypothesis on the zeros. Let $c\in\R$, $R>0$, let $F$ be holomorphic near $\overline D(c,R)$ with
$\overline{F(\bar z)}=F(z)$ and $F(c)\ne0$, and put
\[
 b_n=-\frac{R^n}{(n-1)!}(\log F)^{(n)}(c),\qquad d=\ceil{\sqrt k},\qquad H_k=(b_{2k+i+j})_{0\le i,j<d}.
\]
For $t>0$ let $P_k(t)$ and $N_k(t)$ count the eigenvalues of $H_k$ greater than $t$ and less than $-t$.

\begin{theorem}[Hankel inertia at the centre; \stU]\label{rec:thm:census}
Suppose $D(c,R)$ contains $r$ distinct real zeros and $s$ distinct conjugate pairs of nonreal zeros.
Then $(P_k(k),N_k(k))=(r+s,s)$ for all large $k$. If $a_j$ counts the real boundary zeros of
multiplicity $j$ and $c_j$ the nonreal boundary conjugate pairs of multiplicity $j$, then for each
$j\ge1$, eventually
\[
 P_k\bigl((j-\tfrac12)d\bigr)-P_k\bigl((j+\tfrac12)d\bigr)=a_j+c_j,\qquad
 N_k\bigl((j-\tfrac12)d\bigr)-N_k\bigl((j+\tfrac12)d\bigr)=c_j .
\]
If each $b_n$ in $H_k$ is replaced by $b_n+d^{-1/4}b_{n+1}$, the interior counts persist, and the
positive boundary-scale eigenvalues can be labelled by the boundary orbits so that
$\lambda_k=md+md^{3/4}x+O(\sqrt d)$, where $m$ is the multiplicity and $x$ the real part of the
normalized reciprocal $u$ of either member; the orbit is $c+R(x\pm i\sqrt{1-x^2})$.
\end{theorem}

\begin{proof}
With $f(z)=F(c+Rz)/F(c)$, distinct zeros $z_\ell$ in the closed unit disk, $u_\ell=1/z_\ell$ and
multiplicities $m_\ell$, dividing by $\prod(1-u_\ell z)^{m_\ell}$ and the Cauchy estimate give
$b_n=\sum_\ell m_\ell u_\ell^n+e_n$, $\abs{e_n}\le Ch^{-n}$ for some $h>1$. With $v_d(u)=(1,u,\dots,u^{d-1})^{\mathsf T}$,
$H_k=A_k+B_k+E_k$, where $A_k=\sum_{\abs{u_\ell}>1}m_\ell u_\ell^{2k}v_d(u_\ell)v_d(u_\ell)^{\mathsf T}$, $B_k$ is the
same sum over $\abs{u_\ell}=1$, and $\norm{E_k}\le Cdh^{-2k}$. Write $A_k=WD_kW^{\mathsf T}$ over the reals, with a
column $v_d(u)$ and scalar $mu^{2k}>0$ for a real node, and columns $\sqrt2\Re v_d,\sqrt2\Im v_d$ with a
block of eigenvalues $\pm m\abs u^{2k}$ for a pair. A Vandermonde minor bounds $\sigma_{\min}(W)$ below, so
by congruence $A_k$ has inertia $(r+s,s)$ and nonzero eigenvalues growing exponentially, while
$\norm{B_k+E_k}=O(d)=o(k)$; Weyl's inequalities give the interior counts. For a boundary node $w$ the
projection of $v_d(w)$ onto the column space $U$ of $W$ is $O(1)$, since that space is spanned by the
reversed vectors $u^{1-d}v_d(u)$ with bounded inner products with $v_d(w)$. For distinct unit $w,w'$,
$\abs{v_d(w)^*v_d(w')}\le2/\abs{1-\bar ww'}$ and $\norm{v_d(w)}^2=d$, so with $B_k=LJ_kL^{\mathsf T}$ and $Q$ the
projection onto $U^\perp$, $(QL)^{\mathsf T}(QL)=dI+O(1)$ and $H_k=A_k+QB_kQ+O(\sqrt d)+E_k$, where $QB_kQ$ has
nonzero spectrum $dJ_k+O(1)$; the blocks of $J_k$ have eigenvalues $m$ (real node) and $\pm m$ (pair).
For the modification, each node weight is multiplied by $1+d^{-1/4}u$, the boundary blocks acquire
eigenvalues $m\abs{1+d^{-1/4}u}=m(1+d^{-1/4}x+O(d^{-1/2}))$, and distinct orbits of equal multiplicity
separate at order $d^{3/4}$.
\end{proof}

For $F(t)=\Xi_1(\tfrac12+it)$, or $F_\chi(t)=e^{-i\theta/2}\Xi_\chi(\tfrac12+it)$ for a primitive $\chi$ with
root number $e^{i\theta}$ (a real entire function), $P_k(k)-N_k(k)$ eventually counts the distinct
critical-line ordinates in $(c-R,c+R)$, and nonreal pairs correspond to zeros off the line; under
\cref{src:thm:dgrh} the pair terms vanish (\stG), and if exactly one endpoint of the interval is a zero,
of multiplicity $m$, the boundary eigenvalue satisfies $\lambda_k/d\to m$. No rule at a universal
finite order exists: $1$ and $1-(z/a)^{M+1}$, $0<a<1$, agree to order $M$ at $0$ but have different counts.
For $f(z)=(1-2z)^2(1+4z^2)(1-z)^3(1+z^2)^2(1-z/2)$ exact rational inertia at $k=1024$, $d=32$
detects the double conjugate pair and the triple real zero on the unit circle
(\artifact{grh_periods_sampling_and_zeros/checks/zero_census.txt}). Compared with
\cref{li:thm:contour}, which bounds the imaginary parts of the zeros in a domain from moments of $F'/F$ integrated along its boundary, and
with \cref{li:thm:ritzcert}, which needs a positive spectrum, the data here are Taylor coefficients at
one point and nonreal zeros are allowed.

\subsection{The character from filtered Li coefficients}\label{rec:sec:li}
\sources{dirichlet_li_rigidity}

For a primitive $\chi$ of conductor $f$ with zero multiset $Z_\chi$ use the complex Li coefficients
of Lagarias \cite{Lagarias2007Li}, $u_\rho=1-1/\rho$ and
$\lambda_\chi(k)=\lim_{T\to\infty}\sum_{\rho\in Z_\chi,\abs\rho\le T}(1-u_\rho^k)$. For an integer
stride $h\ge1$ and a complex polynomial $A(z)=\sum_{j=0}^Ja_jz^j$ with $A(1)=1$ the observation is
\[
 X_{\chi,h,A}(n)=\sum_{j=0}^Ja_j\lambda_\chi(hn+j)\qquad(n\ge0).
\]
Thus $A=z^r$ encodes an unknown offset, and several coefficients a weighted mixture of offsets.

\begin{theorem}[Uniform separation; \stU]\label{rec:thm:liseparation}
For distinct triples $(\chi,h,A)$ and $(\psi,k,B)$ and an integer $D\ge0$ there are
$\eps,\eta>0$ and $N_0$ such that for every interval $I$ of at least $N_0$ consecutive nonnegative
integers and every complex polynomial $P$ of degree at most $D$,
\[
 \#\set{n\in I:\abs{X_{\chi,h,A}(n)-X_{\psi,k,B}(n)-P(n)}>\eps}\ge\eta\abs I .
\]
\end{theorem}

\begin{corollary}[Incomplete observations; \stU]\label{rec:cor:liincomplete}
If $\abs{I_j}\to\infty$ and the data on $I_j$ have the form $X_{\chi,h,A}(n)+P_j(n)+e_j(n)$ with
$\sup_j\deg P_j<\infty$ and, for every $t>0$, a proportion of indices with $\abs{e_j(n)}>t$ tending
to zero, then $(\chi,h,A)$ is unique.
\end{corollary}

\begin{proof}
Two representations differ by a polynomial of bounded degree, and outside the two error sets at
level $\eps/2$ the residual difference is at most $\eps$, contradicting
\cref{rec:thm:liseparation}.
\end{proof}

For example, on the windows $[2^{j^2},2^{j^2}+j^2)$, whose union has density zero, samples of
$\lambda_\chi(11n+2)$ and $\lambda_\chi(11n+5)$ mixed with weights $1/4,3/4$ determine $\chi$, the
stride $11$, both offsets and both weights, even with an arbitrary quadratic background on each
window and arbitrary replacements at square indices.

The proof uses the classical facts that $\Xi_\chi$ has order one with zero multiset $Z_\chi$, that
$\rho\mapsto1-\bar\rho$ preserves $Z_\chi$ with multiplicities, and that
\begin{equation}\label{rec:eq:licount}
 N_\chi(T)=\#\set{\rho\in Z_\chi:\abs{\Im\rho}\le T}=\frac T\pi\log\frac{fT}{2\pi e}+O_\chi(\log(2+T))
\end{equation}
\cite{BennettMartinOBryantRechnitzer2021Counting}, so that $\sum\abs\rho^{-2}<\infty$ and a zero at
height $T$ has multiplicity $O_\chi(\log(2+T))$. Let $\Delta x(n)=x(n+1)-x(n)$ and $d=\max(2,D+1)$.
Applying $\Delta^d$ to the common truncation and removing the cutoff,
\begin{equation}\label{rec:eq:liexp}
 \Delta^dX_{\chi,h,A}(n)=\sum_wb_{\chi,h,A}(w)w^n,\qquad
 b_{\chi,h,A}(w)=-(w-1)^d\sum_{\rho\in Z_\chi:\,u_\rho^h=w}A(u_\rho),
\end{equation}
with summable coefficients, bounded nodes and $1$ as the only accumulation point; contributions
with $w=1$ vanish.

\begin{lemma}[Detection on translated intervals; \stU]\label{rec:lem:lishell}
Let $s(n)=\sum_wc_ww^n$ with distinct bounded nodes accumulating only at $1$ and
$M=\sum_w\abs{c_w}<\infty$, and suppose $c_w\ne0$ for some $\abs w\ge1$. Let $c_v\ne0$ have maximal
node modulus $R=\abs v\ge1$. Then $\frac1N\sum_{n=L}^{L+N-1}s(n)v^{-n}\to c_v$ uniformly in $L\ge0$,
and on every sufficiently long interval at least the fraction $\abs{c_v}/(4M)$ of the indices
satisfy $\abs{s(n)}\ge\abs{c_v}R^n/2$.
\end{lemma}

\begin{proof}
For every $\delta>0$ only finitely many nodes have $\abs w\ge1+\delta$, so the maximum exists. For
$w\ne v$ the average of $(w/v)^n$ is at most $\min(1,2/(N\abs{1-w/v}))$ uniformly in $L$; sum over
finitely many nodes and bound the tail by its $\ell^1$ norm. For large $N$ the average has modulus
at least $3\abs{c_v}/4$, and each summand has modulus at most $M$.
\end{proof}

\begin{proposition}[Decoding; \stU]\label{rec:prop:lidecode}
If the coefficients \eqref{rec:eq:liexp} of two triples agree at every $\abs w\ge1$, then their
strides, conductors and filters agree, and their zero multisets differ at only finitely many
points.
\end{proposition}

\begin{proof}
For $w\ne1$ near $1$ an occupied node has the single preimage $u=\exp((\mathrm{Log}\,w)/h)$, since
all but finitely many $u_\rho$ lie in a neighbourhood of $1$ on which $u\mapsto u^h$ is injective. On
an occupied exterior or unit node there, $H(w)=-b(w)/(w-1)^d=mA(u)$ with $m=m_\chi(\rho)$, and
$\bigl\lvert\abs{H(w)}-m\bigr\rvert\le m\abs{A(u)-1}=O(\log(2+\abs{\Im\rho})/\abs\rho)\to0$, so rounding
recovers $m$. Reflection gives $u_{1-\bar\rho}=1/\bar u_\rho$, which supplies the interior nodes
with their multiplicities after rounding. For high zeros
$\abs{\mathrm{Log}\,w}=h/\abs\gamma+O_h(\abs\gamma^{-2})$, so by \eqref{rec:eq:licount} the number
$C(\delta)$ of nodes with $\delta\le\abs{\mathrm{Log}\,w}<\delta_0$ is
\[
 C(\delta)=\frac h{\pi\delta}\Bigl(\log\frac1\delta+\log\frac{fh}{2\pi e}\Bigr)+O\bigl(\log(1/\delta)\bigr),
\]
whose two terms recover $h$ and then $f$. With $h$ common, the unfolded $u$ coincide, so $A(u)=B(u)$
at infinitely many points and $A=B$; every sufficiently high zero is common with its multiplicity,
by $\rho=1/(1-u)$.
\end{proof}

\begin{proof}[Proof of \cref{rec:thm:liseparation}]
By \cref{rec:prop:lidecode} and \cref{rec:lem:char}, distinct triples have a coefficient difference
$c_w\ne0$ at some $\abs w\ge1$. Apply \cref{rec:lem:lishell} to
$s=\Delta^d(X_{\chi,h,A}-X_{\psi,k,B})$, which kills every $P$ of degree at most $D$. The binomial
coefficients of $\Delta^d$ have absolute sum $2^d$, so $\abs{s(n)}\ge\abs{c_v}R^n/2$ forces some
$n+j$, $0\le j\le d$, with $\abs{X-X'-P}\ge\abs{c_v}/2^{d+1}$ there; each witness is used at most $d+1$
times. Take $\eps=\abs{c_v}/2^{d+2}$ and $\eta=\abs{c_v}/(8M(d+1))$.
\end{proof}

Under \cref{src:thm:dgrh} every node has $\abs w=1$, $\Delta^dX_{\chi,h,A}$ is bounded with an
absolutely convergent Fourier series, and every coefficient is a Bohr mean,
\begin{equation}\label{rec:eq:limean}
 b(w)=\lim_{N\to\infty}\frac1N\sum_{n=L}^{L+N-1}\Delta^dX_{\chi,h,A}(n)\,\bar w^{\,n}\qquad(\abs w=1),
\end{equation}
uniformly in $L$; then
$h=\lim_{\delta\downarrow0}\pi\delta C(\delta)/\log(1/\delta)$ and
$f=\frac{2\pi e}h\exp\bigl(\lim_{\delta\downarrow0}[\pi\delta C(\delta)/h-\log(1/\delta)]\bigr)$
(\stG). With unbounded outliers, apply \eqref{rec:eq:limean} to the radial projection onto a disc of
radius $K>\sum_w\abs{b(w)}$ of the differenced data; outside $d+1$ translates of the bad set the
projection error is at most $2^dt$, so the means are unchanged. This uses infinite data; division by
$(w-1)^d$ amplifies finite-sample errors near $1$.

\begin{corollary}[Real observations; \stU]\label{rec:cor:lireal}
With real filter coefficients, observations of $\Re X_{\chi,h,A}$ under the conditions of
\cref{rec:cor:liincomplete} determine $h$, $A$ and the unordered pair $\set{\chi,\bar\chi}$, and
conjugation is exactly the ambiguity.
\end{corollary}

\begin{proof}
Since $\lambda_{\bar\chi}(n)=\overline{\lambda_\chi(n)}$, twice the real observation is the filtered
Li sum of $Z_\chi\cup Z_{\bar\chi}$, which satisfies reflection, the same multiplicity bound and twice
\eqref{rec:eq:licount}. The argument above gives $h$, $A$ and the union up to a finite divisor. The
proof of \cref{rec:lem:char}, applied to the products of the two gamma factors, gives
$L(s,\chi)L(s,\bar\chi)=L(s,\psi)L(s,\bar\psi)$, so $(\chi+\bar\chi)(p)=(\psi+\bar\psi)(p)$ for
$p\nmid f$; every reduced class contains a prime, and orthogonality gives equality of the pairs.
\end{proof}

\subsection{Zero multisets from prime sums: Bohr means and sharp upper limits}\label{rec:sec:prime}
\sources{prime_zero_reconstruction; dirichlet_inverse_problems; three_scale_reconstruction}

Three prime-side observations determine zero multisets with multiplicities: the autocorrelation
of one complex-weighted prime signal (all primitive characters at once), the lag averages of a
single twisted sum sampled at factorials, and the sharp upper limits of integer-coefficient
dilation filters sampled at $x=6^n$. They use different data, and none implies another.

\begin{lemma}[A smoothed explicit formula; \stG]\label{rec:lem:triangle}
Let $\chi$ be primitive of conductor $q$, and $\delta_\chi=\mathbf 1_{q=1}$. For $x\ge1$,
\begin{equation}\label{rec:eq:triangle}
 \sum_{n\le x}\Lambda(n)\chi(n)(1-n/x)=\delta_\chi\frac x2-\sum_\rho\frac{x^\rho}{\rho(\rho+1)}
 +O\bigl(x^{1/4}\log(q+2)\bigr)
\end{equation}
with an absolute constant, and for fixed $\chi$ the error is $O_\chi(1+\log x)$. Moreover
$\sum_{\abs{\Im\rho}>Y}\abs{\rho(\rho+1)}^{-1}\ll\log(q(Y+2))/(Y+1)$ for $Y\ge1$,
$\sum_\rho\abs{\rho(\rho+1)}^{-1}\ll\log(q+2)$, and every interval $[t-1,t+1]$ contains
$O(\log(q(\abs t+2)))$ ordinates, all with absolute constants.
\end{lemma}

\begin{proof}
With $a$ the parity, $\Xi=\Xi_\chi$ satisfies $\Re(\Xi'/\Xi)(\tfrac12+it)=0$ away from zeros, by the
functional equation, and the Hadamard product gives
$\frac{\Xi'}\Xi(s)-\frac{\Xi'}\Xi(2+it)=\sum_\rho\bigl(\frac1{s-\rho}-\frac1{2+it-\rho}\bigr)$. On the
critical line each $1/(s_0-\rho)$ is purely imaginary, so
$\Re\frac{\Xi'}\Xi(2+it)=\sum_\rho\frac{3/2}{(3/2)^2+(t-\Im\rho)^2}$, and the left side is
$O(\log(q(\abs t+2)))$ by the Euler series, Stirling and the conductor term. Evaluating at a
non-ordinate $t'$ within $1/2$ of $t$ gives the local count. Summing over unit intervals gives
$\abs{L'/L(\tfrac14+it,\chi)}\ll\log(q(\abs t+2))$, and on heights $t_M^\pm\in\pm[M,M+1]$ at
distance $\gg1/\log(q(M+2))$ from all ordinates, $L'/L\ll\log^2(q(M+2))$ on $\tfrac14\le\Re s\le2$.
The inverse Mellin transform of $y^s/(s(s+1))$ on $\Re s=2$ is $(1-y^{-1})_+$, so the sum equals
$\frac1{2\pi i}\int_{(2)}-\frac{L'}L(s,\chi)\frac{x^s}{s(s+1)}\dd s$. Shifting to $\Re s=1/4$ crosses
the pole at $1$ (residue $x/2$) and the zeros (residues $-mx^\rho/(\rho(\rho+1))$); the new integral
is $\ll x^{1/4}\int\log(q(\abs t+2))(1+t^2)^{-1}\dd t$. The tail bounds follow from the local count
and $\abs{\rho(\rho+1)}\asymp1+\abs{\Im\rho}^2$. For fixed $\chi$ shift instead to $\Re s=-1/4$, where
the functional equation and Stirling give $L'/L\ll_\chi\log(2+\abs{\Im s})$; the additional residue
at $s=0$ is $b-r(\log x-1)$ if $-L'/L=-r/s+b+O(s)$, and the integral is $O_\chi(x^{-1/4})$.
\end{proof}

\paragraph{One signal for all characters.}
Enumerate the primitive characters as $(\chi_j)_{j\ge0}$, by conductor and then in a fixed order,
with $\chi_0=1$, conductors $q_j$ and $w_j=q_j+1$, and write
$m_j(\gamma)=\ord_{s=1/2+i\gamma}L(s,\chi_j)$. The decoding rests on a code that is independent of
\cref{src:thm:dgrh}. Let $p_j$ be the $j$th prime, $c_0=1$, $k_0=0$ and $c_j=2^{-k_j}\sqrt{p_j}$, where,
with $D_n=2^{\max_{j\le n}k_j}$ and $d_n=D_n^{-1}(3D_n(n+1))^{-(2^n-1)}$, $k_j$ is the least
nonnegative integer with
\[
 2^{-k_j}p_jw_j\le2^{-j-2}\min\Bigl(1,\min_{0\le n<j}\frac{d_n}{n+1}\Bigr).
\]
Then $\sum_{j\ge1}c_jw_j\le1/4$ and $\sum_{j>n}c_jw_j\le2^{-n-2}d_n/(n+1)$. Let $\mathcal M$ be the set
of sequences $(m_j)$ of nonnegative integers with $m_j\le Hw_j$ for some $H$.

\begin{lemma}[Prefix separation; \stU]\label{rec:lem:prefix}
If integers $u_0,\dots,u_n$, not all zero, satisfy $\abs{u_j}\le(n+1)w_j$, then
$\abs{\sum_{j\le n}c_ju_j}\ge d_n$. Consequently $E(m)=\sum_jc_jm_j$ is injective on $\mathcal M$; if
$m_j\le Hw_j$ and $n+1\ge H$, the first $n+1$ entries of $m$ are the unique prefix $v$ with
$0\le v_j\le(n+1)w_j$ and $\abs{A'-\sum_{j\le n}c_jv_j}<d_n/2$, for any real $A'$ with
$\abs{A'-E(m)}<d_n/8$.
\end{lemma}

\begin{proof}
The field $K_n=\Q(\sqrt{p_1},\dots,\sqrt{p_n})$ has degree $2^n$ with independent sign changes (if
$\sqrt{p_n}\in K_{n-1}$ it would be a rational multiple of a product of earlier roots, as a common
eigenvector of the sign changes, and squaring contradicts the parity of the $p_n$-valuation). So
$z=D_n\sum c_ju_j$ is a nonzero algebraic integer whose conjugates are at most
$\tfrac94D_n(n+1)<3D_n(n+1)$, and its norm is a nonzero integer. For decoding, the true prefix lies
within $H\sum_{j>n}c_jw_j\le d_n/4$ of $E(m)$, hence within $3d_n/8$ of $A'$, and any other prefix is
at least $d_n$ away from it. Two distinct sequences of $\mathcal M$ have different prefixes for large
$n$, so $E$ is injective.
\end{proof}

\begin{theorem}[Universal zero reconstruction; \stG]\label{rec:thm:universal}
With these coefficients put $W(n)=\sum_jc_j\chi_j(n)$, so that $\Re W\ge3/4$ and $\abs W\le5/4$,
\[
 S(x)=\sum_{n\le x}\Lambda(n)W(n)(1-n/x),\qquad F(t)=e^{-t/2}\bigl(S(e^t)-e^t/2\bigr).
\]
The limits $C(h)=\lim_T\frac1T\int_0^TF(t+h)\overline{F(t)}\,\dd t$ and
$P(\gamma)=\lim_U\frac1U\int_0^UC(h)e^{-i\gamma h}\dd h$ exist, and with $\rho_\gamma=\tfrac12+i\gamma$,
\[
 \abs{\rho_\gamma(\rho_\gamma+1)}\sqrt{P(\gamma)}=\sum_{j\ge0}c_jm_j(\gamma)\qquad(\gamma\in\R).
\]
Hence $C$ determines the labelled nontrivial zero multisets of all primitive Dirichlet
$L$-functions, and through induction those of all Dirichlet $L$-functions in the open strip.
\end{theorem}

\begin{proof}
Since $\sum c_j\log(q_j+2)<\infty$, \cref{rec:lem:triangle} summed over $j$ gives
$F(t)=G(t)+O(e^{-t/4})$ with $G(t)=\sum_{\gamma}a_\gamma e^{i\gamma t}$ over the countable union of
ordinates, $a_\gamma=-A_\gamma/(\rho_\gamma(\rho_\gamma+1))$, $A_\gamma=\sum_jc_jm_j(\gamma)\ge0$, and
$B=\sum_\gamma\abs{a_\gamma}<\infty$; all terms are nonnegative before the phases, so regrouping is
free, although an ordinate may occur for infinitely many characters. Expanding the product in the
absolutely summable double series and using
$\frac1T\int_0^Te^{i(\gamma-\eta)t}\dd t\to\mathbf 1_{\gamma=\eta}$ gives
$C(h)=\sum_\gamma\abs{a_\gamma}^2e^{i\gamma h}$, and likewise $P(\gamma)=\abs{a_\gamma}^2$. By the local
count, $m_j(\gamma)\le C_1\log(q_j(\abs\gamma+2))\le H_\gamma(q_j+1)$ for an integer
$H_\gamma\ge C_1(1+\log(\abs\gamma+2))$, so $(m_j(\gamma))_j\in\mathcal M$ and \cref{rec:lem:prefix}
decodes it.
\end{proof}

The full complex correlation is used; its even real part or its value at lag zero would not
suffice. Primes alone also suffice: with
$S_{\mathrm p}(x)=\sum_{p\le x}(\log p)W(p)(1-p/x)$ and
$b=\tfrac23\sum_{j:\chi_j^2\text{ principal}}c_j$, the signal
$F_{\mathrm p}(t)=e^{-t/2}(S_{\mathrm p}(e^t)-e^t/2)+b$ has the same autocorrelation (\stG). Indeed
partial summation and the prime number theorem in progressions give
$\sum_{p\le y}\chi_j(p)^2\log p\,(1-p^2/y^2)=\tfrac23\epsilon_jy+o_j(y)$, uniformly bounded by
Chebyshev, so dominated convergence gives the prime squares $b\sqrt x+o(\sqrt x)$, while higher
powers contribute $O(x^{1/3}\log^2x)$; thus $F_{\mathrm p}-F=o(1)$.

For an observed ordinate with $K_\gamma=\abs{\rho_\gamma(\rho_\gamma+1)}$, an approximation
$P'\ge0$ with $\abs{P'-P(\gamma)}<d_n^2/(64K_\gamma^2)$ gives $\abs{K_\gamma\sqrt{P'}-A_\gamma}<d_n/8$,
so the prefix of length $n+1\ge H_\gamma$ is decoded; the thresholds $d_n$ are extremely small.
For a finite set $E$ of frequencies with tail $\eps=\sum_{\gamma\notin E}\abs{a_\gamma}$ and minimal
spacing $\Delta$, $\bigl\lvert\frac1T\int_0^TF(t+h)\overline{F(t)}\dd t-C(h)\bigr\rvert\le
O(1/T)+3B\eps+2B^2/(T\Delta)$ for $h\ge0$; without control of $\Delta$ this is no complexity bound.

The finite-family version of the same statistic detects zeros off the line.

\begin{proposition}[Finite-family converse; \stU]\label{rec:prop:finiteconverse}
Let $J$ be a finite set of primitive characters with weights $c_j>0$, $S_J(x)$ the corresponding
sum with the kernel $(1-n/x)$, and $d$ the weight of $\chi_0$ (zero if absent). Then
$S_J(x)-dx/2=O(\sqrt x)$ if and only if all nontrivial zeros of the $L(s,\chi_j)$, $j\in J$, lie on
the critical line.
\end{proposition}

\begin{proof}
One direction is \cref{rec:lem:triangle}. Conversely, for $\Re s>1$,
$\int_1^\infty(S_J(x)-dx/2)x^{-s-1}\dd x=\frac{\sum_jc_j(-L'/L)(s,\chi_j)}{s(s+1)}-\frac d{2(s-1)}$, and the
bound continues the left side holomorphically to $\Re s>1/2$. At a zero $\rho$ there the residue
of the right side is $-\sum_jc_j\ord_\rho L(s,\chi_j)/(\rho(\rho+1))\ne0$. Reflection excludes the
left half of the strip.
\end{proof}

No converse for the countable aggregate is asserted: meromorphic continuation of a countable sum
through accumulating zeros is not automatic.

\paragraph{Factorial samples of one twisted sum.}
For primitive $\chi$, with $m_\chi(t)$ the multiplicity at $\tfrac12+it$ and an integer lag
$h\ge1$, put
\[
 B_\chi(x)=\frac{\sum_{n\le x}\Lambda(n)\chi(n)(1-n/x)-\delta_\chi x/2}{\sqrt x},\qquad
 c_\chi(t)=-\frac{m_\chi(t)}{(\tfrac12+it)(\tfrac32+it)},
\]
\[
 C_{\chi,h}(N)=\frac1N\sum_{n=1}^NB_\chi((n+h)!)\,\overline{B_\chi(n!)}.
\]
By \cref{rec:lem:triangle}, $B_\chi(x)=\sum_tc_\chi(t)e^{it\log x}+O_\chi((1+\log x)/\sqrt x)$ with
$\sum_t\abs{c_\chi(t)}<\infty$.

\begin{lemma}[Factorial phases; \stU]\label{rec:lem:factorial}
For real $u\ne0$ and $N\ge2$,
$\bigl\lvert\sum_{n=1}^Ne^{iu\log(n!)}\bigr\rvert\ll\sqrt N\bigl(\sqrt{\abs u}+1/\sqrt{\abs u}\bigr)$, and
inserting a weight $n^{iv}$, $v$ fixed, keeps the bound with a constant depending on $v$.
\end{lemma}

\begin{proof}
On $[M,2M]$ the second derivative of $u\log\Gamma(x+1)/(2\pi)$ has constant sign and size
$\asymp\abs u/M$, since $(\log\Gamma)''(x+1)=\sum_{j\ge0}(x+1+j)^{-2}$. Van der Corput's
second-derivative estimate bounds each dyadic block by $O(M\sqrt{\abs u/M}+\sqrt{M/\abs u})$; sum the
blocks, and use partial summation with $\abs{(x^{iv})'}=\abs v/x$ for the weight.
\end{proof}

\begin{theorem}[Lag averages at factorials; \stG]\label{rec:thm:factorial}
For every integer $h\ge1$,
\[
 C_{\chi,h}(N)=\sum_t\frac{\abs{c_\chi(t)}^2}{1+iht}e^{iht\log N}+o(1),
\]
the series converging absolutely and uniformly in $\log N$; the sequence $C_{\chi,h}(N)$ has no
limit; and for every real $v$,
\[
 \lim_{X\to\infty}\frac1{\log X}\sum_{N\le X}\frac{C_{\chi,h}(N)}{N^{1+ihv}}=\frac{\abs{c_\chi(v)}^2}{1+ihv}.
\]
Hence $m_\chi(v)=\sqrt{(\tfrac14+v^2)(\tfrac94+v^2)}\,\abs{c_\chi(v)}$ is recovered for every $v$, and
the averages at $h=1$ determine $\chi$.
\end{theorem}

\begin{proof}
Replace each $B$ by its series; the error vanishes in the mean. For a pair of modes,
$e^{it\log((n+h)!)-iu\log(n!)}=e^{i(t-u)\log(n!)}n^{iht}(1+O_{h,t}(1/n))$. For $t\ne u$ the mean
tends to zero by \cref{rec:lem:factorial}; for $t=u$,
$\frac1N\sum_{n\le N}n^{iht}=N^{iht}/(1+iht)+O_{h,t}((1+\log N)/N)$. Absolute convergence controls
the tails. The logarithmic mean of $N^{iw}$ is $\mathbf 1_{w=0}$, which gives the second formula. Every
primitive $L$-function has zeros of arbitrarily large height, so some coefficient at a nonzero
frequency is nonzero and $C_{\chi,h}(N)$ cannot converge. The character follows from
\cref{rec:lem:char}.
\end{proof}

\begin{corollary}[Distances and joint laws; \stG]\label{rec:cor:factorialjoint}
For every real $v$ and primitive $\chi,\psi$,
\[
 \lim_{N\to\infty}\frac1N\sum_{n\le N}B_\chi(n!)(n!)^{-iv}=c_\chi(v),\qquad
 \lim_{N\to\infty}\frac1N\sum_{n\le N}\abs{B_\chi(n!)-B_\psi(n!)}^2
 =\sum_t\frac{\abs{m_\chi(t)-m_\psi(t)}^2}{(\tfrac14+t^2)(\tfrac94+t^2)},
\]
and its square root is a metric on primitive characters. For primitive $\chi_1,\dots,\chi_r$, reals
$b_j$, $V(y)=(B_{\chi_j}(e^{y+b_j}))_j$ and every continuous $H:\C^r\to\C$,
$\lim_N\frac1N\sum_{n\le N}H(V(\log n!))=\lim_Y\frac1Y\int_0^YH(V(y))\dd y$, both limits existing.
\end{corollary}

\begin{proof}
The first two statements follow from the series and \cref{rec:lem:factorial}; zero distance gives
equal multisets and \cref{rec:lem:char} equal characters. For the joint laws, the vector stays in a
compact box; approximate $H$ uniformly by polynomials in the coordinates and their conjugates, and
these by trigonometric polynomials after truncating the series. A nonzero mode has continuous mean
zero and, by \cref{rec:lem:factorial}, factorial mean zero.
\end{proof}

At fixed multiplicative shifts the logarithmic separation is constant, which is why the joint laws
agree with the continuous ones; at lag $h$ in the factorial index it is $h\log n+O_h(1/n)$, and the
diagonal term $n^{iht}$ retains an oscillation in $N$.

\paragraph{Sharp upper limits on a sparse grid.}
Fix $q$, weights $b_\chi\ge0$ on the characters modulo $q$ with $b_{\bar\chi}=b_\chi$, and put
$w_{\boldsymbol b}(m)=\sum_\chi b_\chi\chi(m)\in\R$ and
$P_{\boldsymbol b}(x)=\sum_{m\le x}\Lambda(m)w_{\boldsymbol b}(m)(x/m-1)$. For $r\in\set{2,3,6}$ let
\[
 (\recA_rf)(x)=(r+1)f(x)-f(rx)-rf(x/r),\qquad g_r(t)=r+1-2\sqrt r\cos(t\log r)>0,
\]
$W_{\boldsymbol b}(t)=\sum_\chi b_\chi m_\chi(t)$, where $m_\chi(t)$ is the order of $L(s,\chi)$ at
$\tfrac12+it$, $\eta(0)=1$, $\eta(t)=2$ for $t>0$, $g(t)=(g_2,g_3,g_6)(t)$, and
\[
 \mu_{\boldsymbol b}=\sum_{t\ge0}\frac{\eta(t)W_{\boldsymbol b}(t)g_2(t)^2}{1/4+t^2}\delta_{g(t)},\qquad
 M_{ijk}=\limsup_{n\to\infty}6^{-n/2}\bigl(\recA_2^{i+2}\recA_3^j\recA_6^kP_{\boldsymbol b}\bigr)(6^n).
\]
The support of $\mu_{\boldsymbol b}$ lies in $K=\prod_r[(\sqrt r-1)^2,(\sqrt r+1)^2]$.

\begin{lemma}[Positive spectral expansion; \stG]\label{rec:lem:threescale}
For every real polynomial $H$, as $x\to\infty$,
\[
 x^{-1/2}\bigl(\recA_2^2H(\recA_2,\recA_3,\recA_6)P_{\boldsymbol b}\bigr)(x)
 =\sum_{t\ge0}\frac{\eta(t)W_{\boldsymbol b}(t)g_2(t)^2}{1/4+t^2}H(g(t))\cos(t\log x)+O_{q,\boldsymbol b,H}(x^{-1/4}),
\]
absolutely and uniformly in $\log x$; if $H\ge0$ on $K$ all coefficients are nonnegative.
\end{lemma}

\begin{proof}
$P_{\boldsymbol b}(x)=\frac1{2\pi i}\int_{(2)}D_{\boldsymbol b}(s)\frac{x^s}{s(s-1)}\dd s$ with
$D_{\boldsymbol b}=-\sum_\chi b_\chi L'/L(s,\chi)$, since $(y-1)_+$ has Mellin kernel $y^s/(s(s-1))$.
The operators act on $x^s$ by $G_r(s)=r+1-r^s-r^{1-s}$, with $G_r(0)=G_r(1)=0$, so $G_2^2$ removes the
possible double pole at $1$; equivalently $\recA_r$ annihilates $x$ and maps $x\log x$ to a multiple
of $x$. Shift to $\Re s=1/4$ along heights at distance $\gg_q1/\log(\abs T+2)$ from all ordinates,
where $L'/L\ll_q\log^2(\abs T+2)$; the new line contributes $O(x^{1/4})$. The residue at
$\rho=\tfrac12+it$ is $W_{\boldsymbol b}(t)x^\rho G_2(\rho)^2H(G(\rho))/(\rho(1-\rho))$ with
$\rho(1-\rho)=1/4+t^2$ and $G_r(\rho)=g_r(t)$. Pair $t$ with $-t$ using $W_{\boldsymbol b}(-t)=W_{\boldsymbol b}(t)$.
\end{proof}

\begin{theorem}[Three-scale extremal reconstruction; \stG]\label{rec:thm:threescale}
Every $M_{ijk}$ is finite and $M_{ijk}=\int_Ku_2^iu_3^ju_6^k\,\dd\mu_{\boldsymbol b}(u)$. For a polynomial
$H\ge0$ on $K$, the sharp upper limit of the filtered sum on $6^n$ equals $\int_KH\dd\mu_{\boldsymbol b}$, and
equals the upper limit over all real $x$. For $t\ge0$ and
$Q_{N,t}(u)=\prod_r\bigl(1-(u_r-g_r(t))^2/(16r)\bigr)^N=\sum c_{ijk}(N,t)u_2^iu_3^ju_6^k$,
\[
 W_{\boldsymbol b}(t)=\frac{1/4+t^2}{\eta(t)g_2(t)^2}\lim_{N\to\infty}\sum_{i,j,k\le2N}c_{ijk}(N,t)M_{ijk},
\]
the sums decreasing to the limit. The moments determine $\mu_{\boldsymbol b}$, and the moment matrices
are positive semidefinite.
\end{theorem}

\begin{proof}
By \cref{rec:lem:threescale} the sampled quantity is $\sum_ta_t\cos(nt\log6)+o(1)$ with $a_t\ge0$ and
$\sum a_t=A<\infty$, so the upper limit is at most $A$. For any finite set of $t$, the torus sequence
$(e^{imt\log6})$ has a convergent subsequence, and quotients of two late terms give arbitrarily
large $n$ with all retained cosines near $1$; this gives the value $A$. If $g(t)=g(t')$, then
$\cos(at)=\cos(at')$, $\cos(bt)=\cos(bt')$ and $\cos((a+b)t)=\cos((a+b)t')$ with $a=\log2$, $b=\log3$;
the third equality forces a common sign $\varepsilon$ in the first two congruences, and the
irrationality of $a/b$ gives $t'=\varepsilon t$, so $g$ is injective on $[0,\infty)$. Each factor of
$Q_{N,t}$ lies in $[0,1]$ on $K$ and equals $1$ only at $u_r=g_r(t)$, so $Q_{N,t}$ decreases to the
indicator of $\set{g(t)}$ and dominated convergence gives the atom $\mu_{\boldsymbol b}(\set{g(t)})$;
$g_2\ge(\sqrt2-1)^2$ allows division. Uniqueness of $\mu_{\boldsymbol b}$ follows from
Stone--Weierstrass.
\end{proof}

Each filter of fixed degree uses finitely many arguments $2^{n+a}3^{n+b}$, integers for large $n$,
and no centering constant. The identity-class weights $b_\chi=1/\varphi(d)$ modulo each $d\mid q$ give
the aggregate orders of the primitive characters of conductor exactly $q$ by Möbius inversion:
$\sum_{\chi\bmod q\text{ primitive}}m_\chi(t)=\sum_{d\mid q}\mu(q/d)\varphi(d)W_d^{\mathrm{id}}(t)$. The
weights $\chi_0+\Re\chi\ge0$ give $m_\chi$ for real $\chi$ and $(m_\chi+m_{\bar\chi})/2$ otherwise.

\begin{proposition}[Oriented multiplicities; \stG]\label{rec:prop:oriented}
For a nonreal $\chi$ modulo $q$ and $t_0\in\R$, let
$P^\chi(x)=\sum_{m\le x}\Lambda(m)\chi(m)(x/m-1)$, choose $r\in\set{2,3}$ with $\sin(t_0\log r)\ne0$ if
$t_0\ne0$, write $T_rf(x)=f(rx)$, and let $F_{N,t_0}(x)$ be the real part of
\[
 x^{-1/2}\recA_2^2Q_{N,\abs{t_0}}(\recA_2,\recA_3,\recA_6)
 \bigl(2I-r^{-1/2}e^{-it_0\log r}T_r-r^{1/2}e^{it_0\log r}T_{1/r}\bigr)P^\chi(x).
\]
Its sharp upper limit decreases, as $N\to\infty$, to
$\frac{g_2(t_0)^2}{1/4+t_0^2}m_\chi(-t_0)\bigl[2-2\cos(2t_0\log r)\bigr]$, with a positive bracket. So
real sharp upper limits determine the oriented zero multiset of $L(s,\chi)$.
\end{proposition}

\begin{proof}
The contour argument without pairing gives, with $g$ extended evenly,
\[
 x^{-1/2}(\recA_2^2HP^\chi)(x)=\sum_{t\in\R}\frac{m_\chi(t)g_2(t)^2}{1/4+t^2}H(g(t))e^{it\log x}+O(x^{-1/4}).
\] The last factor has multiplier $2-2\cos((t-t_0)\log r)\ge0$ on
$x^{1/2+it}$, so the real part is a cosine series with nonnegative coefficients and the recurrence
argument gives its sum. As $N\to\infty$ the localizer keeps $t=\pm t_0$, and the multiplier vanishes
at $t=t_0$. Replace $t_0$ by $-t_0$; at $t_0=0$ omit the last factor.
\end{proof}

Two scales do not suffice in general: for $t=\pi(2/\log3-1/\log2)$ and $t'=\pi(2/\log3+1/\log2)$
one has $0<t<t'$, $g_2(t)=g_2(t')$, $g_3(t)=g_3(t')$ and $g_6(t)\ne g_6(t')$, so positive
summable cosine spectra with an atom at $t$ or at $t'$ have the same two-scale endpoints. This is
an obstruction for general spectra; it does not assert that $t$ or $t'$ is an ordinate of a zero.
The finite experiment for $q=1$ at $n\le6$ in
\artifact{three_scale_reconstruction/verification.txt} illustrates the filters and is not an
approximation with certified error.

\subsection{Zero divisors from sparse samples of prime sums}\label{rec:sec:sparse}
\sources{prime_observations_and_progression_free_sets}

Two further observations sample a single weighted prime sum on an explicit sparse set of radii.
The first, at the powers of two, recovers the character-labelled multiplicities at a fixed modulus
under \cref{src:thm:dgrh}, although on this grid ordinates differing by $2\pi/\log2$ oscillate
identically. The second uses only primes $p\equiv1\pmod q$, or the primes splitting completely in a
Galois field, at $O(\log X\log\log X)$ radii up to $X$, and recovers the zeros with no hypothesis on
them.

\paragraph{Labelled multiplicities at the powers of two.}
Fix $q$, order its $h=\varphi(q)$ characters $\chi_0,\dots,\chi_{h-1}$ with $\chi_0$ principal, let $m_j(\gamma)$ be the
order of $L(s,\chi_j)$ at $\tfrac12+i\gamma$, and put
\[
 T_2=\frac{2\pi}{\log2},\quad\theta=2^{1/h},\quad a(n)=\sum_j\theta^j\chi_j(n),\quad
 B=\frac{\log(102q)+10}{\log(7/6)},\]\[\quad D=2hB,\quad A=h+4\log(512D^h),\quad L=\frac{4\sqrt A}{T_2}.
\]
With $g(y)=\frac{T_2}{2\sqrt{\pi A}}\exp(-T_2^2y^2/(4A))$, $k=g+g'/T_2$, $W(z)=(1+z/T_2)\exp(Az^2/T_2^2)$ (so that
$\int k(y)e^{-zy}\dd y=W(z)$) and $t_n=n\log2$, observe
\[
 Y_n=e^{t_n/2}W(\tfrac12)-\sum_{\substack{m\ge2\\\abs{\log m-t_n}\le L\sqrt{t_n}}}\frac{\Lambda(m)a(m)}{\sqrt m}k(t_n-\log m).
\]

\begin{theorem}[Labelled zeros from dyadic samples; \stG]\label{rec:thm:dyadic}
The sequence $(Y_n)_{n\ge1}$ determines, for every real $\gamma$, the vector $(m_j(\gamma))_{j<h}$, that is, the
zero divisor in $0<\Re s<1$ of every character modulo $q$ with its label and multiplicity. The
conclusion persists after any $o(1)$ perturbation of the $Y_n$.
\end{theorem}

\begin{proof}
\emph{Bounded multiplicities.} Jensen's formula on the disk of radius $7/4$ about $2+it$, with
$\abs L\le q\abs s/\Re s$ for a primitive nonprincipal inducer and $(s-1)\zeta(s)/(s+1)$ for $\zeta$, gives
$m_j(T_2u)\le B(1+\abs u)$. \emph{Expansion.} Mellin inversion on $\Re s=2$ and a shift to $\Re s=\tfrac14$,
along heights at distance $\gg_q1/\log(2+\abs v)$ from all ordinates, give, with $c(u)=\sum_j\theta^jm_j(T_2u)$ and
$\mathcal U$ its support,
\[
 Y_n=\sum_{u\in\mathcal U}c(u)(1+iu)e^{-Au^2}e^{2\pi inu}+o(1),
\]
absolutely and uniformly in $n$; by \cref{src:thm:dgrh} every crossed zero is on the critical line, the
pole contributes $e^{t_n/2}W(\tfrac12)$, and the truncation of the prime sum to the window costs $o(1)$
because $T_2^2L^2/(4A)=4$. \emph{Folding.} For $r\in\R/\Z$,
$Z(r)=\lim_N\frac1N\sum_{n\le N}Y_ne^{-2\pi inr}=\sum_{u\in r+\Z}c(u)(1+iu)e^{-Au^2}$, which vanishes exactly when the
class has no zeros. \emph{Unfolding.} For $\abs u\le R$, $c(u)$ lies in the finite set of
$\sum_jb_j\theta^j$, $0\le b_j\le B(1+R)$, which determines the digits since $X^h-2$ is Eisenstein, and distinct
members differ by at least $\Delta_R=[D(1+R)]^{-(h-1)}$, the norm of a nonzero algebraic integer being at
least one. Write the class as $\epsilon+\Z$ or $\epsilon+\tfrac12+\Z$ with $\abs\epsilon\le\tfrac14$ and recover the singleton
$\set\epsilon$, then the pairs $\set{\epsilon\pm j}$ (respectively $\set{\epsilon\pm(j+\tfrac12)}$) in order, subtracting
the recovered groups: for a pair $u,v$ the functional $P_{u,v}(z)=(\Im z-v\Re z)/(u-v)$ sends $1+iu$ to $1$ and
$1+iv$ to $0$, and after multiplying by $e^{Au^2}$ the later groups contribute less than $\tfrac14\Delta_{R_0}$, by
the choice of $A$. Rounding to the finite set gives $c(u)$ and its digits.
\end{proof}

An error at most $e^{-AR^2}\Delta_R/(4(1+R))$ in $Z(r)$ still allows recovery of all groups through radius $R$.
For $q=3$ the coefficients are $m_0+\sqrt2\,m_1$, which separates a common simple zero ($1+\sqrt2$) from a
double zero of $\zeta$ ($2$); the exact checks in $\Q(\sqrt2)$ are
\begin{center}
\artifact{prime_observations_and_progression_free_sets/dyadic/check_dyadic.json}
\end{center}
Compared with
\cref{rec:thm:threescale}, which reads aggregate weighted multiplicities from sharp upper limits on $6^n$,
and \cref{rec:thm:universal}, which uses a continuous autocorrelation, this observation is one complex
sequence on the regular grid $2^n$ with exact frequency collisions.

\paragraph{Completely split primes at sparse radii.}
Fix $q\ge1$, $\varphi=\varphi(q)$, $Z_q(s)=\prod_{\chi\bmod q}L(s,\chi)$, and let $\Theta_q$ be the supremum of the real
parts of its zeros in $0<\Re s<1$. With $H_1(s)=\exp[-\cos(s-1)]$ and
$K(u)=\frac{e^u}\pi\int_0^\infty e^{-\cosh v}\cos(uv)\dd v$, observe at $t=\log x$
\[
 P_q(t)=\varphi\sum_{p\equiv1\bmod q}(\log p)K(t-\log p),\qquad E_q(t)=P_q(t)-e^{t-1},
\]
at the radii $\mathcal T=\bigcup_{j\ge1}\bigl(h_j\Z\cap[2^{j-1},2^{j+2}]\bigr)$, $h_j=[64(j+1)]^{-1}$, of which
$O(T\log(2+T))$ lie below $T$.

\begin{theorem}[Zeros and the growth exponent from split primes; \stU]\label{rec:thm:mellin}
For every fixed $q$,
\[
 \limsup_{t\to\infty,\ t\in\mathcal T}\frac{\log(1+\abs{E_q(t)})}t=\Theta_q .
\]
So GRH for the characters modulo $q$ is equivalent to $E_q(t)=O_q(e^{t/2})$ on $\mathcal T$, and to
$O_{q,\eps}(e^{(1/2+\eps)t})$ there for every $\eps>0$. The samples determine the multiset of zeros of $Z_q$ in
$0<\Re s<1$; perturbations $O(e^{\alpha t})$, $\alpha<1$, preserve the poles and residues in $\Re s>\max(\tfrac12,\alpha)$,
and for $\alpha<\tfrac12$ the whole multiset, the central point included.
\end{theorem}

\begin{proof}
On $0\le\Re s\le2$, $\abs{H_1(\sigma+iv)}\le\exp[-\cos(1)\cosh v]$, and $K$ is the inverse Laplace transform of
$H_1$ on any line $0\le c\le2$, so $P_q$ is that of $H_1A_q$, $A_q(s)=\varphi\sum_{p\equiv1}(\log p)p^{-s}$. Differentiating
on $\Re s=2$ gives $\abs{E_q^{(k)}(u)}\ll_qe^{2u}(10\log(k+2))^k$ for $u\ge0$ and large $k$. \emph{Interpolation.}
For large $t$, $j=\floor{\log_2t}$, $h=h_j$, $N=\floor{t\sqrt{\log(t+2)}}$, interpolate the samples at the $2N+1$
nodes $h(b+i)$, $b=\floor{t/h}$, $\abs i\le N$, all in $\mathcal T$. Each Lagrange basis value at $t$ is at most
$N+1$, so perturbations $O(e^{\alpha u})$ stay $O_\eps(e^{(\alpha+\eps)t})$, and the remainder is
$\ll\exp(2t+o(t))(10eh\log(k+2))^k=O(\exp[-c_0t\sqrt{\log t}])$, since $10e\log2/64<1/3$. \emph{Continuation.}
With $J=\set{a:a^2=1}$, $D=-Z_q'/Z_q=\varphi\sum_{n\equiv1}\Lambda(n)n^{-s}$ and
$B_J(w)=\varphi\sum_{n\bmod q\in J}\Lambda(n)n^{-w}$, one has $A_q(s)=D(s)-B_J(2s)+U(s)$ with $U$ holomorphic on $\Re s>\tfrac13$,
and $B_J$ regular on $\Re w\ge1$ except for its pole at $1$ of residue $\abs J$. So for $\tfrac12\le\Re\rho<1$,
\[
 \Res_{s=\rho}A_q(s)=-\ord_\rho Z_q-\tfrac12\abs J\mathbf 1_{\rho=1/2},
\]
and the Laplace transform $\int_0^\infty E_q(t)e^{-st}\dd t=H_1(s)A_q(s)-\frac{e^{-1}}{s-1}-V_q(s)$, $V_q$ holomorphic on
$\Re s<2$, continues with residues $H_1(\rho)\Res_\rho A_q$; replacing $E_q$ by the interpolant changes it by an
entire function, $H_1$ never vanishes, and reflection gives the zeros left of the line.
\emph{Exponent.} If $b$ is the sampled limsup, interpolation gives $E_q=O(e^{(b+2\eps)t})$, so every zero
right of $\tfrac12$ and the pole at $\tfrac12$ force $b\ge\Theta_q$. Conversely, for $\Theta_q<d<1$ shift the integral
for $P_q$ to $\Re s=d$, where $D=O(\log^2(2+\abs{\Im s}))$, crossing only the pole at $1$ (contribution $e^{t-1}$);
the double-exponential decay of $H_1$ kills the horizontals. Under the critical-line hypothesis shift
to $\tfrac13<d<\tfrac12$ along heights avoiding the ordinates, which gives
$E_q(t)=-e^{t/2}\bigl(\tfrac{\abs J}2H_1(\tfrac12)+\sum_{\Re\rho=1/2}H_1(\rho)e^{it\Im\rho}\bigr)+O(e^{dt})$ with a bounded bracket.
\end{proof}

The square-prime correction $\tfrac12\abs J$ at $\tfrac12$ is unavoidable and makes the prime series always
singular there. Truncating to $x^{1/3}\le p\le x^2$ costs $O(x^{1/3}\log x)$, an exponent below $\tfrac12$.

\begin{corollary}[A prescribed power interval; \stU]\label{rec:cor:mellinwindow}
For $0<\delta\le1$ and $a=1/\delta$, the conclusions of \cref{rec:thm:mellin} hold for the finite sums over
$x^{1-\delta}\le p\le x^{1+\delta}$, $p\equiv1\pmod q$, with weight $K_a(u)=\frac{ae^u}\pi\int_0^\infty e^{-\cosh v}\cos(auv)\dd v$, mesh
$h_{j,a}=1/(64a(j+1))$ and centring $x/e$; there are $O_\delta(\log X\log\log X)$ radii through $X$.
\end{corollary}

\begin{proof}
The multiplier $H_a(s)=\exp[-\cos((s-1)/a)]$ has no zeros, and in $1-a\le\Re s\le1+a$ it is at most
$\exp[-\cos(1)\cosh(v/a)]$ at height $v$; the mesh compensates the factor $a^k$ in the derivative bound, and shifting the
$K_a$-contour to $1\mp a$ bounds the discarded tails by
$\ll_ax^{1-a\delta}\log(2x)$, that is, by $\log(2x)$.
\end{proof}

\begin{theorem}[Galois fields; \stU]\label{rec:thm:mellingalois}
For a finite Galois $L/\Q$ with group $G$, $n=\abs G$, replace $\varphi$ by $n$ and the primes $p\equiv1\pmod q$ by the
unramified primes splitting completely in $L$. Then \cref{rec:thm:mellin} and \cref{rec:cor:mellinwindow}
hold with $\zeta_L$ and its zero abscissa $\Theta_L$, and with central correction half the number of
elements of $G$ of order dividing two.
\end{theorem}

\begin{proof}
For $C\le G$ let $\pi_C(g)$ count the cosets of $G/C$ fixed by $g$, so that $-\zeta_{L^C}'/\zeta_{L^C}$ has coefficient
$(\log p)\pi_C(g_p^m)$ at $p^m$. With $r$ the number of subgroups of order two,
$n\mathbf 1_{g^2=1}=2\sum_{\abs C=2}\pi_C(g)+(1-r)\operatorname{reg}(g)$: at $g=1$ this is numerical, an involution lies in exactly
one such subgroup after each of the $n$ conjugations, and other $g$ give zero. So
$B=2\sum_{\abs C=2}(-\zeta_{L^C}'/\zeta_{L^C})+(1-r)(-\zeta_L'/\zeta_L)$ has first-prime coefficient $n(\log p)\mathbf 1_{g_p^2=1}$ and residue $1+r$
at $1$, and the proof of \cref{rec:thm:mellin} applies with $\abs J=1+r$; the zero count and local partial
fractions for $\zeta_L$ follow from its Hadamard product \cite{LagariasOdlyzko1977Effective}. For the endpoint,
each $\zeta_L/\zeta_{L^C}$ with $[L:L^C]=2$ is entire by the quadratic Hecke factorization, so the critical-line
hypothesis for $\zeta_L$ places the zeros of these $\zeta_{L^C}$ on the line.
\end{proof}

For the splitting field of $T^3-2$, with group $S_3$, complete splitting at $p>3$ means three distinct
roots modulo $p$, and the central correction is $2$; this is a sparse finite-observation criterion for
RH of a degree-six field whose zeta function contains a two-dimensional Artin factor
(\artifact{prime_observations_and_progression_free_sets/mellin/check_mellin.txt}). By
\cref{rec:thm:twolags} two fixed lags of a different, continuous profile also measure the zero edge; the
present samples are discrete, sparse and restricted to one residue class.

\subsection{Prime powers and the modulus from a three-cutoff extremal profile}\label{rec:sec:profile}
\sources{prime_power_rigidity}

The weight $1/n-1/x$ of the source remainder (\cref{src:sec:mechanism}) is kept, but the residue
class is allowed to vary and the auxiliary twist is dropped. Fix $q\ge1$, $G_q=(\Z/q\Z)^\times$,
$d=\varphi(q)$, and for $a\in G_q$ put
\[
 A_a(x)=\sum_{\substack{n\le x\\n\equiv a\ (q)}}\Lambda(n)\Bigl(\frac1n-\frac1x\Bigr),\qquad
 c_a=\lim_{x\to\infty}\Bigl(A_a(x)-\frac{\log x}d\Bigr),\qquad
 Y_a(t)=e^{t/2}\Bigl(A_a(e^t)-\frac td-c_a\Bigr),
\]
and, for $t\ge\abs h$,
\[
 B_{a,h}(t)=2Y_1(t)-Y_a(t+h)-Y_{a^{-1}}(t-h),\qquad F_q(a,h)=\limsup_{t\to\infty}B_{a,h}(t).
\]
For $n=p^k$ let $v=v_p(q)$, $m=q/p^v$ and $P_q(a;p^k)=\mathbf 1_{a\equiv p^k\ (m)}/\varphi(p^v)$, a
probability distribution on $G_q$, concentrated at $p^k\bmod q$ when $p\nmid q$.

\begin{lemma}[The weighted explicit formula; \stU]\label{rec:lem:ramp}
Let $Z_\chi$ be the zero multiset of $L(s,\chi^*)$ in the open strip. As $t\to\infty$,
\[
 Y_a(t)=\frac1d\sum_{\chi\bmod q}\overline{\chi(a)}\sum_{\rho\in Z_\chi}\frac{e^{(\rho-1/2)t}}{\rho(1-\rho)}
 +O_q\bigl((1+t)e^{-t/2}\bigr).
\]
For a single primitive character the difference between its normalized sum and its zero series is
real analytic in $t>0$.
\end{lemma}

\begin{proof}
For primitive $\xi$, the sum $\sum_{n\le x}\Lambda(n)\xi(n)(1/n-1/x)$ is
$\frac1{2\pi i}\int_{(2)}-\frac{L'}L(s,\xi)\frac{x^{s-1}}{s(s-1)}\dd s$. Shifting left to
$\Re s=-J-1/2$ on heights avoiding the zeros, a zero of order $r$ gives
$rx^{\rho-1}/(\rho(1-\rho))$, the pole at $1$ gives a constant (or $\log x-1-\recgE$ for $\zeta$), the
point $0$ gives $x^{-1}$ times a polynomial of degree at most one in $\log x$, and a trivial zero
$-j$ gives $-x^{-j-1}/(j(j+1))$; the last series extends holomorphically to $\Re t>0$. The left
integrals tend to zero as $J\to\infty$, and $\sum_\rho\abs{\rho(1-\rho)}^{-1}<\infty$. Passing from
$\chi^*$ to $\chi$ changes the sum by $O_q((1+\log x)/\sqrt x)$ after normalization, and
orthogonality averages over characters.
\end{proof}

\begin{theorem}[Prime-power recovery; \stG]\label{rec:thm:primepower}
For $a\in G_q$, $F_q(a,h)$ is finite and nonnegative, continuous on $h>0$ and real analytic off the
logarithms of prime powers, and at every integer $n>1$
\[
 F_q'(a,\log n+)-F_q'(a,\log n-)=
 \begin{cases}-\dfrac{2\log p}{\sqrt{p^k}}P_q(a;p^k),&n=p^k,\\[3pt]0,&\text{otherwise}.\end{cases}
\]
In fact $F_q(a,h)=2\bigl(K_q(1,0)-K_q(a,h)\bigr)$ with
$K_q(a,u)=\frac1d\sum_\chi\sum_{\rho\in Z_\chi}\overline{\chi(a)}e^{i\gamma u}/(1/4+\gamma^2)$. The same
$F_q$ results with $e^t$ restricted to $b^j$, $j\in\N$, for any fixed $b>1$, and one sequence $j_\ell$
realizes all these upper limits for all moduli and residues, uniformly on bounded sets of lags.
\end{theorem}

\begin{proof}
The coefficient of $e^{i\gamma t}/(d(1/4+\gamma^2))$ in $B_{a,h}$ is
$2-\overline{\chi(a)}e^{i\gamma h}-\chi(a)e^{-i\gamma h}=\abs{1-\overline{\chi(a)}e^{i\gamma h}}^2\in[0,4]$.
So the upper limit is at most the sum of these nonnegative coefficients. Enumerating all ordinates of
all Dirichlet characters, the powers of $(e^{i\gamma_1\log b},\dots,e^{i\gamma_N\log b})$ return near
the identity at arbitrarily large indices, and a diagonal sequence works for every ordinate; with
absolute summability this gives equality, without any independence of ordinates. For the jumps use
the exact primitive expansion of the proof of \cref{rec:lem:ramp}: a prime term
$\Lambda(n)\chi^*(n)(e^{h/2}/n-e^{-h/2})\mathbf 1_{h\ge\log n}$ is continuous with derivative jump
$\Lambda(n)\chi^*(n)/\sqrt n$, and the rest is analytic near $h$. For $n=p^k$, only characters trivial
on $G_{p^v}$ have $\chi^*(p)\ne0$, and orthogonality on $G_m$ gives
$\frac1d\sum_\chi\overline{\chi(a)}\chi^*(p^k)=P_q(a;p^k)$.
\end{proof}

For example, modulo $12$ the profiles at $a=5$ and $a=11$ both detect the prime $2$, which lies in
no reduced class, with $P_{12}(5;2)=P_{12}(11;2)=1/2$. Summing the jumps over $a$ and multiplying by
$-\sqrt n/2$ gives $\Lambda(n)$. Since a perturbation of the weights with running sum $E(u)=o(\sqrt u)$
changes $Y_a(\log x)$, after recentering, by at most $2\sup_{u\ge x}\abs{E(u)}/\sqrt u$, the profile is
unchanged by erasing any finite initial segment of the prime table.

\begin{theorem}[Rigidity of the modulus; \stG]\label{rec:thm:rigid}
On $X_q=G_q\times\R$ put
\[
 D_q\bigl((a,s),(b,t)\bigr)=\sqrt{F_q(a^{-1}b,t-s)} .
\]
This is a metric. With
$q^\flat=q/2$ for $q\equiv2\pmod4$ and $q^\flat=q$ otherwise, $(X_q,D_q)$ is isometric to $(X_r,D_r)$
if and only if $q^\flat=r^\flat$. The surjective isometries of $X_q$ are exactly
$(a,t)\mapsto(a_0a^\epsilon,b+\epsilon t)$ with $a_0\in G_q$, $b\in\R$, $\epsilon=\pm1$.
\end{theorem}

\begin{proof}
By the formula for $F_q$, $D_q$ is a Hilbert distance with coordinates
$\overline{\chi(a)}e^{i\gamma t}/\sqrt{d(1/4+\gamma^2)}$, so it satisfies the triangle inequality and is
translation invariant. For $0<u<\log2$ there are no prime terms, and the trivial zeros give
$\partial_uK_q(1,u)=\tfrac12\log u+O_q(1)$, hence $F_q(1,u)=u\log(1/u)+O_q(u)$. A point at distance
zero from the identity fixes all spectral coordinates; a nonzero time displacement would move the
unbounded slope at $(1,0)$ to a point with finite slopes, contradicting \cref{rec:thm:primepower},
and with no displacement $\chi(a)=1$ for all $\chi$ since every primitive inducer has zeros.

Choose a zeta ordinate $\gamma_0\ne0$. Along a $D_q$-continuous path the coordinate $e^{i\gamma_0t}$ has
a continuous phase; subtracting it leaves a continuous map into the countable space
$G_q\times(2\pi/\gamma_0)\Z$, which is constant. So an isometry acts by
$(a,t)\mapsto(\sigma(a),f_a(t))$ with homeomorphisms $f_a$, and the small-scale law, through
$g(y)=y^2/(2\log(1/y))$, gives $\abs{f_a(t+u)-f_a(t)}/\abs u\to1$, whence $f_a(t)=\epsilon_at+b_a$.
Profiles between two lines are nonconstant by Dirichlet's theorem and the jump formula, so
orientations agree; after a simultaneous inversion all are positive. For two lines $a,b$ and large
primes $p\equiv a^{-1}b\pmod q$ the target has a kink at $\log p+b_b-b_a$, so $Cp\in\Z$ with
$C=e^{b_b-b_a}$; two primes and B\'ezout give $C\in\Z$, and the inverse gives $C^{-1}\in\Z$. After
normalizing $(1,0)$, every prime $p\nmid qr$ forces $\sigma(p\bmod q)=p\bmod r$. For $q=r$ this proves the
description of the isometries. In general, with $L=\lcm(q,r)$, the reductions $G_L\to G_q,G_r$ have
the same kernel, so $\varphi(p^{v_p(q)})=\varphi(p^{v_p(r)})$ for all $p$, which forces $q^\flat=r^\flat$.
Conversely, for odd $m$ the reduction $G_{2m}\to G_m$ is an isomorphism preserving primitive inducers
and every term of $F$.
\end{proof}

Moduli $8$ and $12$ have isomorphic unit groups, but $5\equiv13\pmod8$ while $5\not\equiv13\pmod{12}$,
so their distances differ; moduli $5$ and $10$ give the same distances.

Without \cref{src:thm:dgrh} the same statistic measures the edge of the zero set. Let $\Theta_q$ be the
supremum of $\Re\rho$ over $\bigcup_\chi Z_\chi$, so $\Theta_q\ge1/2$.

\begin{theorem}[Two fixed lags; \stU]\label{rec:thm:twolags}
For every $a\in G_q$ and $0<h_1<h_2$,
\[
 \limsup_{t\to\infty}\frac1t\log\Bigl(1+\max_{j=1,2}\abs{B_{a,h_j}(t)}\Bigr)=\Theta_q-\tfrac12 .
\]
In particular, boundedness of $B_{a,h_1}$ and $B_{a,h_2}$ on a tail is equivalent to Dirichlet GRH
at modulus $q$.
\end{theorem}

\begin{proof}
\Cref{rec:lem:ramp} gives the upper bound. Conversely, with $z=s-\tfrac12$, for $\Re z>1/2$,
$\int_0^\infty Y_a(t)e^{-zt}\dd t=\frac{D_a(s)}{s(s-1)}-\frac{1/d}{(s-1)^2}-\frac{c_a}{s-1}$, where
$D_a=\sum_{n\equiv a}\Lambda(n)n^{-s}$. At a zero $\rho$ with $\eta=\Re\rho-\tfrac12>0$, multiplicities
$m_\chi$, $M=\sum m_\chi>0$ and $C=\sum m_\chi\overline{\chi(a)}$, the transform of $B_{a,h}$ (on a common
tail) has residue $H(h)/(d\rho(1-\rho))$ at $z_0=\rho-\tfrac12$, with
$H(h)=2M-e^{z_0h}C-e^{-z_0h}\overline C$. With $s_0=\abs C/M\le1$ and $u=\eta h$,
$\abs{H(h)}^2/(4M^2)\ge(1-s_0\cosh u)^2$, and with $A=\cosh\eta h_1<B=\cosh\eta h_2$,
$\max(\abs{1-s_0A},\abs{1-s_0B})\ge(B-A)/(B+A)>0$. So one of the two transforms has a pole at $z_0$,
which is impossible if the growth exponent were below $\eta$. Reflection gives the equivalence.
\end{proof}

The second lag is what protects against cancellation at a shared zero: one lag can hit a point
where $H(h)=0$. For $q=1$, with $K(h)=K_1(1,h)$,
\[
 K(h)=Y_1(h)-e^{-h/2}\Bigl(\log(2\pi)-\sum_{k\ge1}\frac{e^{-2kh}}{2k(2k+1)}\Bigr),\qquad
 K(0)=2+\recgE-\log(4\pi),
\]
and one-sided differences at $\log2,\log4,\log6$ computed from finitely many prime powers
reproduce the predicted jumps. For example, the two one-sided slopes at $\log2$ are
approximately $0.450030537632177$ and $-0.530227605836371$, and their difference agrees with
$-2\log2/\sqrt2$. The check is
\artifact{prime_power_rigidity/check.py}, with output \artifact{prime_power_rigidity/checks.txt}.

\subsection{The modulus and the twist from the signed zero spectrum}\label{rec:sec:spectrum}
\sources{dirichlet-grh-research; exponent_seven_and_grh_transfers; grh-torsion-and-zero-data}

Here the data are on the zero side and the unknowns are arithmetic. The function $F$ below is the
signed counterpart of the function $K_q$ of \cref{rec:thm:primepower}: its coefficients are
differences of multiplicities of an untwisted and a twisted family, and its derivative jumps
again sit at logarithms of primes.

The Mellin residues of the smoothed prime remainder in the source proof record only the
difference between two sums of zero multiplicities; the relevant declaration of the source is
\code{New.isPolar_primeRemainder} (see \cref{src:sec:mechanism}). Contributions at a common zero can cancel exactly,
and the names of the characters have disappeared. Nevertheless the signed data determine both
arithmetic parameters, and for the parameters of the lattice argument a finite part of the
data does so with a certificate.

\paragraph{Setting.}
Two classes of parameters occur.
\begin{itemize}
\item The \emph{general class}: $q>0$ with $8\mid q$, and a prime $d\equiv7\pmod8$ coprime
to $q$.
\item The \emph{lattice class}: $18480\mid q$ and $d$ a prime with $d\equiv q/2-1\pmod q$.
Then $(d,q)=1$, $8\mid q$, $d\equiv7\pmod8$ and $q\mid2(d+1)$, so that $q\le2(d+1)$; this is
a subclass of the general class, and it is the one used in the source.
\end{itemize}
Write $h=\varphi(q)$ and $\eta(n)=(n/d)$. For each $\chi\bmod q$ let $\chi^*$ be its inducing
primitive character. For a real ordinate $\gamma$ put $\rho_\gamma=\tfrac12+i\gamma$ and
\begin{equation}\label{rec:eq:mgamma}
 m(\gamma)=\frac1h\sum_{\chi\bmod q}\Bigl(\ord_{\rho_\gamma}L(s,\chi^*)-
   \ord_{\rho_\gamma}L(s,\chi^*\eta)\Bigr),\qquad
 b_\gamma=\frac{m(\gamma)}{1/4+\gamma^2},\qquad
 F(t)=\sum_\gamma b_\gamma e^{i\gamma t}.
\end{equation}
Removing imprimitive Euler factors preserves orders in the open critical strip, so $b_\gamma$
is precisely the residue in the source, and by \cref{src:thm:dgrh} every nontrivial zero in
these sums is of the form $\rho_\gamma$. The input to the inverse problem is the signed
coefficient at each frequency, with coincident frequencies combined; both signs of the
ordinate are used.

\begin{proposition}[The spectral function; \stG]\label{rec:prop:ode}
In the general class the series $F$ is absolutely and uniformly convergent, real and even.
For a prime power $n=p^k$ let $Q_p=p^{v_p(q)}$, $q^{(p)}=q/Q_p$ and
\[
 a(p^k)=\frac{1-\eta(p^k)}{\varphi(Q_p)}\,\mathbf 1_{p^k\equiv1\bmod q^{(p)}} .
\]
Then, as distributions on $\R$, with $\delta_u$ the unit point mass at $u$,
\begin{equation}\label{rec:eq:ode}
 F''-\tfrac14F=(\log d)\,\delta_0-\frac2h\cosh(t/2)\,\dd t
 +\sum_{p^k\ge2}\frac{\log p}{p^{k/2}}\,a(p^k)\bigl(\delta_{\log p^k}+\delta_{-\log p^k}\bigr).
\end{equation}
Equivalently, for every real test function $g\in C_c^\infty(\R)$, with
$\widehat g(\gamma)=\int g(t)e^{i\gamma t}\,\dd t$,
\begin{equation}\label{rec:eq:weil}
 -\sum_\gamma m(\gamma)\widehat g(\gamma)
 =(\log d)\,g(0)-\frac2h\int g(t)\cosh(t/2)\,\dd t
 +\sum_{n=p^k}\frac{\log p}{\sqrt n}\,a(n)\bigl[g(\log n)+g(-\log n)\bigr].
\end{equation}
\end{proposition}
\begin{proof}
The classical zero count $N(T)=O_{q,d}(T\log(2+T))$ ensures $\sum_\gamma\abs{b_\gamma}<\infty$
by partial summation. Both families are closed under complex conjugation, whence
$b_{-\gamma}=b_\gamma\in\R$. Since $F''-F/4=-\sum_\gamma m(\gamma)e^{i\gamma t}$, the two
displays are equivalent, and it suffices to check \eqref{rec:eq:weil}. The Guinand--Weil
explicit formula for a primitive character of conductor $f$ contributes $\log(f/\pi)\,g(0)$,
the gamma term for its parity, and a negative sum over prime powers with coefficients
$\Lambda(n)\chi^*(n)/\sqrt n$ at one sign of $\log n$ and the conjugate coefficients at the
other. The primitive principal character also contributes
$\widehat g(i/2)+\widehat g(-i/2)$ from the poles at $0,1$ of the completed zeta function.
This is the normalization of \cite[\S2.1]{ChirrePereiraDeLaat2021Primes}; see also
\cite{Booker2015Lfunctions}. Equivalently the formula follows by shifting the contour of the
logarithmic derivative of the completed primitive $L$-function.

As $q>2$, half the characters have each parity, and the odd character $\eta$ interchanges
the parities. All gamma terms therefore cancel between the two families. The two conductors
for $\chi^*$ and $\chi^*\eta$ are $f$ and $fd$, since $(f,d)=1$, which gives $\log d$ at the
origin after the signed subtraction. Only the untwisted family contains the primitive
principal character, which gives the density after division by $h$. At $p^k$ the average of
the untwisted primitive character values is
$\mathbf 1_{p^k\equiv1\bmod q^{(p)}}/\varphi(Q_p)$: a nonprincipal $p$-component makes the
primitive value zero, and orthogonality over the characters modulo $q^{(p)}$ evaluates the
remaining sum. Multiplication by $1-\eta(p^k)$ gives the displayed $a$, also at $p=d$.
\end{proof}

For a point $u$ write $J(u)=F'(u+)-F'(u-)$, with $J(u)=0$ where $F$ is differentiable.

\begin{theorem}[Jumps and uniqueness; \stG]\label{rec:thm:jumps}
In the general class, $F$ is analytic away from the atoms of \eqref{rec:eq:ode}, and
\begin{equation}\label{rec:eq:mass}
 J(0)=\log d,\qquad
 J(\log p)=\frac{\log p}{\sqrt p}\bigl(1-\eta(p)\bigr)\mathbf 1_{p\equiv1\bmod q}
 \quad\text{for every prime }p\nmid qd .
\end{equation}
Thus $F'$ has a jump at $\log p$ precisely if $p\equiv1\pmod q$ and $\eta(p)=-1$, in which
case the jump is $2\log p/\sqrt p$. Let $S$ be the set of primes $p$ at which $F'$ has a
positive jump at $\log p$. For all sufficiently large $X$,
\[
 d=\exp J(0),\qquad q=\gcd\set{p-1:\ p\in S,\ p>X}.
\]
Consequently the data $\set{(\gamma,b_\gamma)}$ determine $(q,d)$, even after they have been
changed at finitely many frequencies; more generally the same reconstruction works after
adding any continuously differentiable function to $F$.
\end{theorem}
\begin{proof}
The right side of \eqref{rec:eq:ode} is smooth apart from a locally finite set of point
masses. Integrating across one such point shows that the jump of $F'$ equals its mass. Away
from the atoms $F$ satisfies an ordinary differential equation with analytic coefficients,
hence is analytic there. This gives \eqref{rec:eq:mass}; evenness then gives
$F(t)=F(0)+\tfrac12\log d\,\abs t+O_{q,d}(t^2)$ near zero.

Every entry $p-1$ with $p\in S$ is divisible by $q$ after the finitely many primes dividing
$qd$ have been discarded. Suppose a prime $\ell$ divided the gcd after division by $q$. We
produce infinitely many primes in $S$ with $\ell\nmid(p-1)/q$. If $\ell\mid q$, require
$p\equiv1+q\pmod{q\ell}$ and independently choose a nonzero quadratic nonresidue modulo $d$.
If $\ell\nmid qd$, then $\ell$ is odd; choose a unit other than $1$ modulo $\ell$, the
residue $1$ modulo $q$, and a nonresidue modulo $d$. If $\ell=d$, the conditions
$p\equiv1\pmod q$ and $\eta(p)=-1$ already imply $p\not\equiv1\pmod\ell$. Each construction
gives a reduced class by the Chinese remainder theorem and infinitely many primes by
Dirichlet's theorem. Thus no $\ell$ divides the quotient. All these jumps survive a $C^1$
perturbation, including any finite trigonometric polynomial.
\end{proof}

\paragraph{The spectral function on actual prime sums.}
Let $R$ be the remainder of the source proof,
\[
 R(x)=\sum_{\substack{n\le x\\n\equiv1\bmod q}}\Lambda(n)\bigl(1-\eta(n)\bigr)
   \Bigl(\frac1n-\frac1x\Bigr)-\frac{\log x}h-c,
\]
where $c$ gives $R(x)\to0$ by the prime number theorem for fixed progressions; in the
source this is the declaration \code{New.primeRemainder}.

\begin{proposition}[\stG]\label{rec:prop:remainder}
In the general class,
\begin{equation}\label{rec:eq:remainder}
 \sqrt x\,R(x)=F(\log x)+O_{q,d}\Bigl(\frac{1+\log x}{\sqrt x}\Bigr)\qquad(x\ge1).
\end{equation}
In the lattice class the comparison is exact for every $x\ge1$:
\[
 F(0)=-2c-\log d,\qquad
 \sqrt x\,R(x)=F(\log x)+\frac{c+\log d-h^{-1}\log x}{\sqrt x} .
\]
\end{proposition}
\begin{proof}
Use $F'(0+)=(\log d)/2$ and solve \eqref{rec:eq:ode} on $t\ge0$. Write
$S^*(x)=\sum_{p^k\le x}\log(p)\,a(p^k)(p^{-k}-x^{-1})$. The homogeneous solutions are
$\cosh(t/2)$ and $\sinh(t/2)$, the particular solution for the continuous forcing is
$-2t\sinh(t/2)/h$, and an atom of mass $B$ at $u>0$ adds $2B\sinh((t-u)/2)$ for $t\ge u$.
Consequently
\[
 \frac{F(\log x)}{\sqrt x}=S^*(x)-\frac{1-x^{-1}}h\log x+\frac{F(0)}2(1+x^{-1})
 +\frac{\log d}2(1-x^{-1}).
\]
The difference between $S^*(x)$ and the sum in $R(x)$ comes exclusively from powers of
$p\mid q$. It is $C_{\mathrm{bad}}+O_q((1+\log x)/x)$, where
$C_{\mathrm{bad}}=\sum_{p\mid q}\sum_{k\ge1}\log(p)\,a(p^k)p^{-k}$; this follows by bounding
a geometric tail for each such $p$ and at most $\log x/\log p$ terms of size $2\log(p)/x$.
Taking $x\to\infty$ gives $c+C_{\mathrm{bad}}+(F(0)+\log d)/2=0$, proving
\eqref{rec:eq:remainder}. In the lattice class the source proves $\eta(p)=1$ for every
$p\mid q$ (\code{New.Parameters.chiK_eq_one_of_dvd}). Then $C_{\mathrm{bad}}=0$ and every
correction at $p\mid q$ vanishes.
\end{proof}

\begin{example}\label{rec:ex:9239}
The pair $(q,d)=(18480,9239)$ lies in the lattice class. The two prime jumps at $258721$ and
$794641$ have $p-1=14q$ and $p-1=43q$, respectively, so their gcd is already $q$. Trial
division and Euler's criterion verify both primes and both nonresidues; the output is in
\artifact{dirichlet-grh-research/critical-line/prime-jumps-output.txt}.
\end{example}

\begin{corollary}[Oscillation at the critical scale; \stG]\label{rec:cor:oscillation}
In the general class the logarithmic mean of $\sqrt x\,R(x)$ is $b_0$, and for every
$\gamma\ne0$ with $b_\gamma\ne0$,
\[
 \liminf_{x\to\infty}\sqrt x\,R(x)\le b_0-\abs{b_\gamma},\qquad
 \limsup_{x\to\infty}\sqrt x\,R(x)\ge b_0+\abs{b_\gamma}.
\]
Infinitely many nonzero frequencies have $b_\gamma\ne0$.
\end{corollary}
\begin{proof}
The mean is obtained by termwise averaging. The averages of $F(t)$ times
$1\pm\cos(\gamma t)$ over $[0,T]$ tend to $b_0\pm b_\gamma$, and each weight is nonnegative
with mean one. Equation \eqref{rec:eq:remainder} transfers the bounds to the prime sum. The
cusp of $F$ at zero is not removed by any finite trigonometric polynomial, so infinitely
many nonzero frequencies survive. The mean includes any genuine central zero; the argument
uses neither simplicity nor linear independence of the ordinates.
\end{proof}

\begin{corollary}[Four central values at once; \stG]\label{rec:cor:chowla}
Take $q=8$ and let $d\equiv7\pmod8$ be prime. Write $\chi_D(n)=(D/n)$ for the primitive
quadratic character of fundamental discriminant $D$, and $R_{8,d}$ for the remainder above.
Then
\[
 B_d:=\lim_{X\to\infty}\frac1{\log X}\int_1^XR_{8,d}(x)\,\frac{\dd x}{\sqrt x}
 =-\sum_{D\in\set{-d,\,4d,\,-8d,\,8d}}\ord_{s=1/2}L(s,\chi_D)\ \in\ -2\Z_{\ge0}.
\]
Thus $B_d>-2$ is equivalent to the simultaneous nonvanishing at $1/2$ of these four
quadratic $L$-functions.
\end{corollary}
\begin{proof}
Equation \eqref{rec:eq:remainder} identifies the mean with $b_0$. The untwisted functions
are $\zeta$ and the $L$-functions of the characters of discriminants $-4,8,-8$. The
alternating eta series gives $\zeta(1/2)<0$, and the alternating beta series gives
$L(1/2,\chi_{-4})>0$. For $8$ and $-8$ the signs on each block of four successive odd
integers are respectively $(+,-,-,+)$ and $(+,+,-,-)$; the sums of these blocks are positive
because $y^{-1/2}$ is strictly decreasing and strictly convex. The twist is $\chi_{-d}$ and
the four resulting discriminants are the ones displayed. Finally each completed quadratic
$L$-function has root number $+1$, so each central order is even.
\end{proof}

Proving the strict lower bound $B_d>-2$ for every such $d$ would establish this subfamily of
Chowla's conjecture on central nonvanishing \cite{Soundararajan2000Nonvanishing}. The
one-sided estimate $R(x)\ll x^{-1/2+\eps}$ of the source allows every fixed value of $B_d$
and supplies no such bound.

\paragraph{Reading a jump from one high band.}
The reconstruction can discard every zero of height below $T$, and its errors can be made
explicit.

\begin{lemma}[\stG]\label{rec:lem:Fbound}
In the general class $\norm F_\infty\le6\log q+3\log d+48$. In the lattice class
$\norm F_\infty\le9\log d+60$.
\end{lemma}
\begin{proof}
For a primitive character of conductor $f$, let $\Xi$ be its completed function, with an
additional factor $s(s-1)$ in the zeta case. Its Hadamard logarithmic derivative and
functional equation give
\[
 \sum_\rho\Re\frac1{2-\rho}=\Re\frac{\Xi'}{\Xi}(2)\le\tfrac12\log f+4 .
\]
Indeed the constant in the real part of the Hadamard logarithmic derivative vanishes by
evaluating at a regular point on the critical line, where both real parts vanish by GRH and
reflection. At $2$ the gamma contribution is at most $1/2$, since the digamma function is
increasing and $\psi(2)<1$. The logarithmic Euler derivative has absolute value at most
$\sum_{n\ge2}(\log n)/n^2<2$ by the integral test, and the zeta correction is $1+1/2$.
Discard $-\tfrac12\log\pi$. On GRH, $(1/4+\gamma^2)^{-1}\le6\Re(2-\rho)^{-1}$. Adding the
two families therefore gives $\norm F_\infty\le6\log q+3\log d+48$. In the lattice class
$q\le2(d+1)\le4d$.
\end{proof}

Use the convention $\widehat g(s)=\int g(t)e^{-ist}\,\dd t$ in the next statement and set
\[
 g(t)=\cos(3t/2)\Bigl(\frac{\sin(t/16)}{t/16}\Bigr)^8,\qquad g(0)=1 .
\]
Each sinc factor has frequency support $[-1/16,1/16]$, so $\widehat g$ is supported on
$[-2,-1]\cup[1,2]$; it is an explicit piecewise polynomial, obtained by eight convolutions
of interval indicators followed by two translations. Define the band reading
\begin{equation}\label{rec:eq:AT}
 A_T(u)=-\frac1T\sum_{T\le\abs\gamma\le2T}m(\gamma)\,\widehat g(\gamma/T)\,e^{i\gamma u}.
\end{equation}

\begin{proposition}[Explicit band reading; \stG]\label{rec:prop:band}
Let $(q,d)$ be in the general class and $u\in\R$. Choose $0<r\le1$ such that $(u-2r,u+2r)$
contains no potential atom of \eqref{rec:eq:ode} except possibly $u$; at $u=0$ one may take
$r=\log(2)/3$, and at $u=\log p$ one third of
$\min\set{\log(p/(p-1)),\log((p+1)/p)}$. For $Tr\ge1$ put
\[
 B_u(T)=\frac{512}{7T}\cosh\bigl((\abs u+2r)/2\bigr),\qquad
 K_u(T)=\frac{2\cdot16^8}{7r^7T^8}\Bigl(108T^2+\frac{40T}r+\frac6{r^2}+\frac14\Bigr).
\]
Then
\begin{equation}\label{rec:eq:error}
 \abs{A_T(u)-J(u)}\le B_u(T)+\norm F_\infty K_u(T),
\end{equation}
and in the lattice class $\norm F_\infty$ may be replaced by $9\log d+60$.
\end{proposition}
\begin{proof}
Write $\mu=F''-\tfrac14F$. The reading $A_T(u)$ equals the pairing
$\langle\mu,g(T(\cdot-u))\rangle$, defined as
$\int F(t)(D^2-\tfrac14)[g(T(t-u))]\,\dd t$; termwise integration follows from the absolute
uniform convergence of $F$. Multiply the test by a cutoff equal to one on $\abs{t-u}\le r$,
zero on $\abs{t-u}\ge2r$, and between these equal to $1-10z^3+15z^4-6z^5$ at
$z=\abs{t-u}/r-1$. Its first two derivatives are bounded by $2/r$ and $6/r^2$. On the local
part only the selected mass and the density of \eqref{rec:eq:ode} contribute. Since
$\int\abs g\le2(16+16/7)=256/7$, the density costs at most $B_u(T)$. For the complementary
part integrate by parts onto the test. Writing $g(t)=16^8t^{-8}\cos(3t/2)\sin^8(t/16)$ shows
that for $\abs t\ge1$ its derivatives of orders $0,1,2$ are bounded by $16^8\abs t^{-8}$
times $1,10,108$, respectively: the trigonometric factor has frequency at most $2$ and total
absolute exponential coefficients at most $1$. The product rule and
$\int_{\abs{t-u}\ge r}\abs{t-u}^{-8}\,\dd t=2/(7r^7)$ give $K_u$. The cutoff is $C^2$, which
suffices by approximation, and the integrable differentiated tests justify passage to the
full real line. This last step permits localization even though the positive prime masses
and the smooth density on the right of \eqref{rec:eq:ode} individually grow exponentially at
infinity.
\end{proof}

\begin{remark}\label{rec:rem:fastband}
A test with vanishing moments removes the term $B_u(T)$ at the cost of explicit constants.
Take a real even $w\in C_c^\infty(\set{s:1<\abs s<2})$ with $\int_\R w=2\pi$ and
$g(v)=(2\pi)^{-1}\int w(s)e^{isv}\,\dd s$. Then $g(0)=1$ and $\int v^jg(v)\,\dd v=0$ for
every $j\ge0$, since $w$ vanishes near zero, and the same localization gives, for every
fixed $u\in\R$ and $A>0$ in the general class,
\[
 J(u)=-\frac1T\sum_{T<\abs\gamma<2T}m(\gamma)w(\gamma/T)e^{i\gamma u}+O_{q,d,u,w,A}(T^{-A}):
\]
the atom contributes $J(u)$, the smooth compactly supported density contributes
$O_A(T^{-A})$ by Taylor expansion and the vanishing moments, and away from $u$ the $L^1$
norms of the test and of its first two derivatives are $O_A(T^{-A})$ by Schwartz decay.
\end{remark}

\paragraph{Errors in the spectral data.}
The formulas quantify errors in the data. Let $E_T$ bound the total absolute missing,
inserted or altered weight and let $\delta_T\le1$ bound the displacement of matched
ordinates; use records in a padded band, so that every entry that can cross an endpoint is
included. If $M_T$ bounds the absolute weight of the matched entries, a valid absolute error
for the reading \eqref{rec:eq:AT} is
\begin{equation}\label{rec:eq:noise}
 e_u(T)=\frac{\norm{\widehat g}_\infty E_T}T+\frac{M_T\delta_T}T
 \Bigl(\abs u\,\norm{\widehat g}_\infty+\frac{\norm{\widehat g'}_\infty}T\Bigr)
 +e_{\mathrm{arithmetic}},
\end{equation}
because $\widehat g(v/T)e^{ivu}$ is Lipschitz in $v$ with constant
$\norm{\widehat g'}_\infty/T+\abs u\norm{\widehat g}_\infty$. Both norms can be bounded from
the explicit spline. For each fixed packet $M_T=O_{q,d}(T\log(2+T))$ by zero counting, so
$E_T=o(T)$ and $\delta_T=o(1/\log T)$ preserve every fixed jump in the limit; for actual
certificates $M_T$ can be bounded by the observed absolute weight plus the certified weight
budget. The readings use signed entries and do not require pairing two nearby ordinates
before summation.

\begin{corollary}[\stG]\label{rec:cor:discrepancy}
Two distinct parameter pairs of the general class have absolute signed multiplicity
discrepancy $\gg T$ in each sufficiently high dyadic band $T<\abs\gamma<2T$.
\end{corollary}
\begin{proof}
Their jump functions differ at some point by \cref{rec:thm:jumps}, and \cref{rec:rem:fastband}
at that point bounds the discrepancy from below.
\end{proof}

\begin{theorem}[A certified stopping rule from one band; \stG]\label{rec:thm:band}
Let $(q,d)$ be in the lattice class. From a sufficiently high band $T\le\abs\gamma\le2T$ of
the aggregate signed data one can recover $(q,d)$ with a computable certificate. No upper
bound for $q$ or $d$ is part of the input. On an insufficient band the procedure may return
\emph{undetermined}. If the band readings \eqref{rec:eq:AT} have certified absolute errors
$e_u(T)$ tending to zero for each fixed observation site $u$, the procedure succeeds on
every sufficiently high band.
\end{theorem}
\begin{proof}
Let $a$ be the real part of the observed $A_T(0)$ and write $K=K_0(T)$,
$D=B_0(T)+60K+e_0(T)$. When $9K<1$, \eqref{rec:eq:error} confines $\log d$ to
\[
 \Bigl[\frac{a-D}{1+9K},\ \frac{a+D}{1-9K}\Bigr].
\]
Use outward-rounded arithmetic and test the primes in the exponential of this interval. If
there is a unique eligible prime, it is $d$; otherwise decline this band. For fixed true $d$
the interval shrinks to $\log d$. List the divisors $q'$ of $2(d+1)$ with $18480\mid q'$ and
$d\equiv q'/2-1\pmod{q'}$.

For every pair of distinct candidate moduli, find a prime $p$ coprime to all candidates and
to $d$, with $\eta(p)=-1$, for which their indicators in \eqref{rec:eq:mass} differ.
Enumerating primes terminates: choose a prime power $\ell^a$ occurring to a higher exponent
in one modulus than in the other. Modulo the least common multiple of the candidates choose
$1$ at all other components and a unit congruent to $1$ modulo the lower $\ell$-power but
distinct from $1$ modulo $\ell^a$. If the lower power is $1$, then $\ell$ is odd, since $8$
divides each candidate. Set a nonresidue at $d$ independently. Dirichlet's theorem supplies
infinitely many primes in this class. For each witness intersect the two possible masses,
$0$ and $2\log p/\sqrt p$, with the certified interval around the observed
$\Re A_T(\log p)$ from \eqref{rec:eq:error} plus the measurement error. Retain compatible
moduli. Each witness eventually separates its pair because $p$ is fixed and $B_u,K_u,e_u$
tend to zero. A surviving singleton is certified.
\end{proof}

A single accepted band certifies both integers. This assumes a genuine packet in the stated
parameter class; authenticating a putative spectrum is a separate problem. The affine bound
\eqref{rec:eq:error} is what supplies a test for accepting a band before either parameter is
known.

\begin{example}\label{rec:ex:band}
Take $d=46199$ and either $q=18480$ or $q=92400$. A genuine band at $T=10^{12}$ suffices,
even with absolute reading error $10^{-7}$ at each of the two sites. Use $r=1/8$ at zero and
$r=1/(4p)$ at $p=110881$. The origin interval contains only $d$, which leaves exactly these
two candidate moduli. The witness $p$ has $\eta(p)=-1$ and $p-1=6\cdot18480$. The two
possible masses differ by more than $2/334$. This is an a priori guarantee for either
packet; no zero ordinate at height $10^{12}$ was computed. The error margins are verified
by rational arithmetic in the following files.
\begin{center}
\artifact{exponent_seven_and_grh_transfers/other-results/check_decoder_example.py}\\
\artifact{exponent_seven_and_grh_transfers/other-results/decoder_example.txt}
\end{center}
\end{example}

\paragraph{Recovery from an initial segment, with a height bound.}
When the signed list is known through a height $T$, probes that annihilate the smooth
density of \eqref{rec:eq:ode} exactly give a bound for the height that suffices.

\begin{theorem}[Finite recovery with a height bound; \stG]\label{rec:thm:recovery}
Let $(q,d)$ be in the lattice class. Knowing the signed list $\set{(\gamma,m(\gamma))}$
through a sufficiently high finite height determines $q$ and $d$. The algorithm needs no
prior numerical upper bound on either. For approximate input it accepts enclosures for the
finite sums it uses; with complete lists and arbitrarily refinable certified errors it
terminates. For exact input a sufficient height is $O(d^2\log^3(2d))$ with an absolute
implied constant.
\end{theorem}

The height estimate is independent of the cost of computing or certifying the input lists.

\begin{proof}
\emph{Probes.} In this proof $\widehat g(\gamma)=\int g(t)e^{i\gamma t}\,\dd t$, as in
\eqref{rec:eq:weil}. If a real compactly supported $g$ kills the $\cosh(t/2)$ integral, has
value one at $u$, and its support isolates $u$ among the potential atoms, then by
\eqref{rec:eq:weil} the quantity $-\sum_\gamma m(\gamma)\widehat g(\gamma)$ equals $J(u)$,
with $J$ as in \eqref{rec:eq:mass}. The origin and prime identities alone would only give
an asymptotic reconstruction; to certify a finite stopping time one must estimate the unseen
zero contribution before knowing the conductors.

\emph{A bound containing the unknown on both sides.} For a primitive character of conductor
$f$ and parity $a\in\set{0,1}$, let $B_f$ denote the sum of $(1+\gamma^2)^{-1}$ over its
nontrivial zeros, with multiplicity. Hadamard factorization and the functional equation give
at $s=2$
\[
 \sum_\gamma\frac{3/2}{9/4+\gamma^2}=\Re\frac{L'}L(2)+\tfrac12\log(f/\pi)
 +\tfrac12\psi(1+a/2)+\tfrac32\mathbf 1_{f=1}.
\]
This is the real logarithmic derivative of the completed entire function, with $s(s-1)$
included for zeta \cite{LamzouriLiSoundararajan2015Conditional}; GRH gives the indicated
denominators. Multiplying the summand by $3/2$ majorizes $(1+\gamma^2)^{-1}$. For $f>1$,
$\abs{L'/L(2)}\le\sum_{n\ge2}(\log n)/n^2<1$ (isolate $2$ and $3$, integrate from $3$). For
$f=1$ the logarithmic derivative is negative. Since $\psi(1+a/2)\le\psi(2)<\log\pi$, in
both cases $B_f\le\tfrac34\log f+\tfrac94$. Averaging absolute values before taking the
signed difference yields
\[
 W:=\sum_\gamma\frac{\abs{m(\gamma)}}{1+\gamma^2}
 \le\tfrac32\log q+\tfrac34\log d+\tfrac92\le\tfrac94\log d+7,
\]
using $q\le3d$. For any chosen test $g$, suppose $A$ bounds
$\sup_{\abs\gamma>T}(1+\gamma^2)\abs{\widehat g(\gamma)}$, and let $[\ell_-,\ell_+]$ enclose the real
part of the negative sum through height $T$. The imaginary part of the exact full sum
vanishes. For the origin probe put $y=\log d$. Whenever $9A/4<1$, the arithmetic certificate
is
\begin{equation}\label{rec:eq:certificate}
 \frac{\ell_--7A}{1+9A/4}\ \le\ y\ \le\ \frac{\ell_++7A}{1-9A/4}.
\end{equation}
Exponentiate with outward rounding and stop the first stage when there is just one
admissible prime in the resulting interval. After that stage the error for each new probe
is at most $A(\tfrac94\log d+7)$, independent of $q$. A finite interval can likewise be
certified by summing enclosures of $m(\gamma)\widehat g(\gamma)$; omitted or uncertain
entries require an enclosure of their possible total contribution. Without such an input
error guarantee a list may omit a decisive term.

\emph{Two explicit probes.} A simple choice at location $v$ is the translate by $v$ of
\[
 g_a(t)=\frac{f_a''(t)-f_a(t)/4}{f_a''(0)-1/4},\qquad f_a(t)=\bigl(1-(t/a)^2\bigr)_+^8 .
\]
It equals one at its center and kills both $e^{t/2}$ and $e^{-t/2}$ by integration by parts.
Use $a=1/2$ at zero, or $a=1/[3(p+1)]$ at $\log p$; the distance to the nearest other
logarithm of a positive integer is at least $\log(1+1/p)\ge1/(p+1)$. Four integrations by
parts permit $A=\norm{g_a^{(4)}}_1(1+T^2)/T^4$.

For the sharper height estimate, take $b=1/6$ at zero and $b=1/[10(p+1)]$ at $\log p$. Set
$k=\floor{bT/2}\ge4$. Convolve the uniform probability density on $[-b/k,b/k]$ with itself
$k$ times and divide by its value at zero to obtain $w$. It is symmetric and unimodal
(averaging any symmetric unimodal density on a centered interval preserves this property).
Before normalization its maximum is at least $1/(2b)$ by total mass and support. Define
\[
 g(t)=w(t-v)-\frac{w(t-v-2b)+w(t-v+2b)}{2\cosh b}.
\]
The support lies within $3b$ of the center, $g(v)=1$, and translation of $e^{\pm t/2}$ by
$\pm2b$ cancels its integral. The Fourier transform of the unnormalized $w$ is
$\recsinc(b\gamma/k)^k$, where $\recsinc(x)=\sin(x)/x$. For $\abs\gamma>T$ its absolute
value is at most $2^{-k}(T/\abs\gamma)^k$. The translation multiplier is at most two, so one
can use
\[
 A=4b(1+T^2)\,2^{-\floor{bT/2}} .
\]
The probes have sufficient differentiability for the explicit formula: convolution with
smooth approximate identities gives the same identity in the limit, with spectral domination
by a multiple of $(1+\gamma^2)^{-1}$. Increasing $T$ and refining enclosures therefore
certifies $d$ at the origin; one can, for example, choose $T$ dyadically and refine the
finite-sum enclosure to width $1/T^2$ or smaller. For exact data $T=O(\log(2d))$ suffices at
the origin, since adjacent integer logarithms near $d$ are separated by $\gg1/d$.

\emph{Separating the finitely many moduli.} Form $M=2(d+1)$. The remaining candidates are
the divisors $q$ of $M$ with $18480\mid q$ and $M/q$ odd. For two distinct candidates
$q_1<q_2$ there is a unit $y\bmod M$ congruent to one modulo $q_1$ and different from one
modulo $q_2$: choose a prime $\ell$ with $v_\ell(q_1)<v_\ell(q_2)$, prescribe the residue
$1+\ell^{v_\ell(q_1)}$ modulo the full $\ell$-power in $M$, and one at all other prime
powers. If $v_\ell(q_1)=0$, this residue is two and $\ell$ is odd. Now require
$p\equiv y\bmod M$ and $\eta(p)=-1$. These conditions specify one coset of index
$2\varphi(M)$ modulo $Md$, since $(M,d)=1$. Dirichlet's theorem gives a prime in this
coset, and a search with integer tests finds one. The two possible values of $J(\log p)$
are separated by $2\log p/\sqrt p$. A certified interval excluding one value eliminates a
candidate. Continuing until only one remains proves finite recovery.

\emph{The height.} Theorem~1.4 of \cite{LamzouriLiSoundararajan2015Conditional} bounds the
least prime in a coset of index $I>1$ modulo $Q\ge20000$ by $O(I^2\log^2Q)$ under Dirichlet
GRH; the explicit alternatives are $p\le10^9$ or
$p\le((I-1)\log Q+3(I+1)+\tfrac52(\log\log Q)^2)^2$. Here $I\le2M$ and $Q=Md$, so each
witness can be taken with $p\ll d^2\log^2(2d)$. For a sufficiently large absolute constant,
$T$ of order $p\log(2pd)$ makes the spline tail times $W$ less than half the two-value
separation. All witnesses use the same spectral list, so the maximum of these heights
proves the claimed bound.
\end{proof}

\begin{example}\label{rec:ex:recovery}
Take $d=109909799$. Sixteen candidate moduli survive the divisor tests. Even knowing the
totient would leave the pair $1127280,1201200$ at totient $230400$. The prime $34945681$ is
a quadratic nonresidue modulo $d$, and $\gcd(p-1,2(d+1))=1127280$; its jump distinguishes
the two. At $T=768$ the spline at the origin has a tail below $1.045\cdot10^{-12}$ for this
$d$. At $T=6\cdot10^{10}$ the witness probe has a tail below $5.220\cdot10^{-11}$, against a
jump greater than $0.005875$. These are calculations on known parameters; an observed
spectral list of these heights was not computed. Primes, characters, candidate enumeration
and tail arithmetic are checked in the following files.
\begin{center}
\artifact{grh-torsion-and-zero-data/check_recovery.py}\\
\artifact{grh-torsion-and-zero-data/verification/check_recovery.txt}
\end{center}
\end{example}

The band rule of \cref{rec:thm:band} and the initial-segment rule of
\cref{rec:thm:recovery} use different data and neither contains the other: the first
discards all zeros below $T$ and comes with the noise model \eqref{rec:eq:noise}, the second
uses the whole list through $T$ and comes with an explicit sufficient height. Both are
confined to the lattice class, where $q\mid2(d+1)$ leaves finitely many candidate moduli
once $d$ is known; \cref{rec:thm:jumps} and \cref{rec:rem:fastband} cover the general class
in the limit.

\subsection{Omitted primes from normalized zero spacings}\label{rec:sec:gaps}
\sources{dirichlet_zero_recovery}

The zeros of an imprimitive $L$-function on $\Re s=0$ come from the omitted Euler factors. Let
$\chi$ modulo $q$ be induced by $\chi^*$ of conductor $f$, and $P=\set{p\mid q:p\nmid f}$, nonempty,
$m=\abs P$. Then $L(s,\chi)=L(s,\chi^*)\prod_{p\in P}(1-\chi^*(p)p^{-s})$, and by the functional
equation and nonvanishing on $\Re s=1$ every zero on the positive imaginary axis comes from the
finite product. Writing $\chi^*(p)=e^{2\pi i\alpha_p}$ with rational $\alpha_p\in[0,1)$, the
ordinates contributed by $p$ are the positive points of $(\Z+\alpha_p)/\lambda_p$,
$\lambda_p=\log p/(2\pi)$; these lattices are disjoint at positive height by unique factorization,
and each factor has simple zeros. So the intensity is $\Lambda=\frac1{2\pi}\sum_{p\in P}\log p$.
List the positive ordinates as $t_1<t_2<\cdots$, put $g_j=\Lambda(t_{j+1}-t_j)$ and
$S_P(x)=\lim_N\frac1N\#\set{j\le N:g_j>x}$; common rescaling of the heights leaves $S_P$ unchanged.
Put $w_p=\log p/\sum_{r\in P}\log r$, $b=1/\max_pw_p$ and $A_P(x)=\prod_{p\in P}(1-w_px)$.

\begin{proposition}[Gap law for the omitted factors; \stU]\label{rec:prop:gaplaw}
For every $x\ge0$ the limit $S_P(x)$ exists, is independent of the phases $\alpha_p$, and equals
$-A_P'(x)$ for $0\le x<b$ and $0$ for $x\ge b$. If $p_*$ is the largest prime of $P$, the gap law has
an atom at $b$ of mass $w_{p_*}\prod_{r\ne p_*}(1-w_r/w_{p_*})$.
\end{proposition}

\begin{proof}
This is the Barra--Gaspard spacing law for superposed lattices, as stated by Marklof and
Str\"ombergsson \cite{MarklofStrombergsson2013Gaps}. Directly: from the $n$th point
$(n+\alpha_p)/w_p$ of the normalized $p$-lattice, its position in the $r$-lattice is the fractional
part of $nw_r/w_p+\alpha_pw_r/w_p-\alpha_r$. Since $1$ and the $w_r/w_p$, $r\ne p$, are rationally
independent (a relation would give an integer relation between logarithms of distinct primes),
these positions equidistribute in the unit cube by Weyl's criterion. For $0\le x<b$ the next
$r$-point lies farther than $x$ with limiting frequency $1-w_rx$, jointly for all $r$, and the next
$p$-point is $1/w_p>x$ away; starting points of type $p$ have frequency $w_p$. Hence
$S_P(x)=\sum_pw_p\prod_{r\ne p}(1-w_rx)=-A_P'(x)$. Every gap is at most $b$, and the left limit at $b$
gives the atom.
\end{proof}

\begin{theorem}[Recovery from short gaps; \stU]\label{rec:thm:gaps}
Let $\chi,\widetilde\chi$ have nonempty omitted sets $P,Q$, and suppose $S_P=S_Q$ on $(0,\eps)$ for
some $\eps>0$. Then $\abs P=\abs Q$, and if $\abs P\ge3$, then $P=Q$. No absolute density, modulus,
conductor or character value is used.
\end{theorem}

\begin{proof}
Near zero $A_P(x)=1-\int_0^xS_P$ is a polynomial, so the germ determines $A_P$, whose degree is $m$
and whose roots are the $1/w_p$. Hence the sorted weights agree, and with
$c=\sum\log q_i/\sum\log p_i$ the ordered primes satisfy $q_i=p_i^c$. For $m\ge3$, if $c$ were
irrational, the six exponentials theorem \cite{Waldschmidt2005Further} applied to
$\log p_1,\log p_2,\log p_3$ and $1,c$ would make one of the primes $p_i,q_i$ transcendental. So
$c=a/d$ in lowest terms, $q_i^d=p_i^a$, and unique factorization gives $q_i=p_i$.
\end{proof}

For a squarefree principal modulus with at least three prime divisors $P$ is its prime divisor set,
so the short-gap law recovers the modulus. For one prime the normalized law is a point mass at $1$;
for two primes $p_1<p_2$ and $q_1<q_2$ the laws agree exactly when
$\log q_1/\log p_1=\log q_2/\log p_2$, which the four exponentials conjecture would rule out for
distinct pairs. The inverse map is not continuous.

\begin{proposition}[Approximation by unrelated prime sets; \stU]\label{rec:prop:gapinstab}
For every $P$ with $m\ge2$ there are pairwise disjoint sets $Q_n$ of $m$ primes, disjoint from $P$,
with $\min Q_n\to\infty$, whose gap laws converge weakly to that of $P$ and whose short-gap
polynomials converge coefficientwise.
\end{proposition}

\begin{proof}
Order the weights $w_1<\dots<w_m$ of $P$ and, for large $T$, choose by Bertrand's postulate
primes $q_i\in[\floor{e^{Tw_i}},2\floor{e^{Tw_i}}]$; then $\log q_i=Tw_i+O(1)$ and the
normalized weights are $w_i+O_m(1/T)$. Let $T$ grow fast enough along a sequence; the cutoffs
converge to $b$ and the survival functions converge at every $x\ne b$.
\end{proof}

\begin{example}\label{rec:ex:gaps30}
The principal character modulo $30$ and the character modulo $210$ induced by the quadratic
character modulo $7$ both have $P=\set{2,3,5}$, with phases $(0,0,0)$ and $(0,\tfrac12,\tfrac12)$.
With $e_j$ the elementary symmetric functions of $\log2/\log30,\log3/\log30,\log5/\log30$,
\[
 S_P(x)=1-2e_2x+3e_3x^2\approx1-0.63021791x+0.09344797x^2\qquad(0\le x<b),
\]
with $b=\log30/\log5\approx2.11328275$, and about $8.55\%$ of the gaps equal $b$. The empirical
frequencies of $541{,}314$ and $541{,}316$ gaps below height $10^6$ agree with the law within
$1.2\cdot10^{-5}$ at the six tested events; see
\begin{center}
\artifact{dirichlet_zero_recovery/gaps/verify_gaps.py}\\
\artifact{dirichlet_zero_recovery/gaps/example_output.txt}
\end{center}
\end{example}

\subsection{Characters and multiplicity vectors in short windows at large height}\label{rec:sec:localwin}
\sources{local_euler_recovery; translated_dirichlet_zero_divisors; dirichlet_prime_frequency_inverse; prime_pole_windows}

The zeros are now observed only in one window $\abs{\Im\rho-T}\le H$ at a large height $T$. A test
function holomorphic across the closed critical strip and decaying fast enough there makes the
Guinand--Weil formula local, and the radius needed to suppress the unseen zeros is
$\pi^{-1}\log\log\log T+O(1)$. One family of tests reads a single Euler coefficient uniformly in the
conductor; a second, binomial, family gives full rank and both signs of every balanced
combination at the same radius, and survives translation of the functions. Several statements
use complex zeros and no hypothesis on their location. In this subsection
$\widehat f(v)=\int_\R f(t)e^{-ivt}\dd t$, $z_\rho=(\rho-\tfrac12)/i$, and for a primitive $\chi$ of
conductor $q$ and parity $a$,
\[
 G_\chi(t)=\log(q/\pi)+\Re\frac{\Gamma'}\Gamma\Bigl(\frac{1/2+a+it}2\Bigr),\qquad
 G_a^\circ(t)=\Re\frac{\Gamma'}\Gamma\Bigl(\frac{1/2+a+it}2\Bigr).
\]
For tests $h$ holomorphic and rapidly decreasing in a strip wider than $\abs{\Im z}\le\tfrac12$ the
explicit formula reads, with each zero counted with multiplicity
\cite{ChirrePereiraDeLaat2021Primes,Davenport2000Multiplicative},
\begin{equation}\label{rec:eq:gwstrip}
 2\pi\sum_\rho h(z_\rho)=2\pi\mathbf 1_{q=1}\bigl\{h(\tfrac i2)+h(-\tfrac i2)\bigr\}+\int_\R h(t)G_\chi(t)\dd t
 -\sum_{n\ge2}\frac{\Lambda(n)}{\sqrt n}\bigl\{\chi(n)\widehat h(\log n)+\overline{\chi(n)}\widehat h(-\log n)\bigr\}.
\end{equation}
It needs no hypothesis on the zeros; complex tests are allowed by linearity, and the tests below
are reached from compactly supported smooth Fourier transforms by truncation, using the
unit-interval zero count.

\paragraph{One Euler coefficient from one window.}
For $Q\ge2$ and $T\ge2$ put
\[
 \ell=\log\bigl(e^eQ(T+3)\bigr),\qquad R=32\log\ell,\qquad D=D(Q,T)=2+\pi^{-1}\log(4R),
\]
so that $D=\pi^{-1}\log\log\log(QT)+O(1)$. With $w(z)=\cosh(\pi z)$ let
\[
 k_R(z)=c_R\exp\Bigl[-\cosh(\pi z/2)-\frac{Rw(z)}{R+w(z)}\Bigr],
\]
with $c_R$ chosen so that $\int_\R k_R=1$. For a prime $p$ and an integer $m\ge4$ put $u=\log p$,
$b=8p$, $J_b=\mathbf 1_{[-b,b]}/(2b)$, $g=J_b^{*m}$, $K=g*k_R$, and
\[
 h(z)=\frac1{iu}\frac{\dd}{\dd z}\bigl(e^{iuz}K(z-T)\bigr),\qquad
 \mathcal E_p=\frac{\sqrt p}{\log p}\Bigl[\int_\R h(t)G_0^\circ(t)\dd t-2\pi\sum_{\abs{\Im\rho-T}\le H}h(z_\rho)\Bigr].
\]
The integral is known before any zero is observed; for a product of $k$ factors it is multiplied
by $k$.

\begin{theorem}[Local Euler coefficient; \stU]\label{rec:thm:localcoeff}
Let $q_\chi\le Q$, let $p$ be prime and $m\ge4$ an integer, and put $B=8pm$, $H=2B+D(Q,T)$, with
$T\ge H+2$. Then
\[
 \abs{\mathcal E_p-\chi(p)}\le C\,2^{-m}
\]
for an absolute constant $C$. The only $\chi$-dependent input is the multiset of complex zeros
$\rho$ with $\abs{\Im\rho-T}\le H$; under \cref{src:thm:dgrh} their ordinates with multiplicities
suffice. Ramified primes, where $\chi(p)=0$, are included.
\end{theorem}

\begin{lemma}[Decay on the closed strip; \stU]\label{rec:lem:lekernel}
Uniformly for $R\ge2$ and real $v$, $\abs{\widehat k_R(v)}\ll_Ne^{-\abs v/2}(1+\abs v)^{-N}$ for each
$N\ge0$. For fixed $M>0$, $j=0,1$ and $z=x+iy$ with $\abs y\le\tfrac12$,
\[
 \abs{k_R^{(j)}(z)}\ll_Me^{-M\abs x},\qquad \abs{k_R^{(j)}(z)}\ll_Me^{-R/4}e^{-M\abs x}\quad(\abs x\ge D).
\]
\end{lemma}

\begin{proof}
In this strip $\Re w\ge0$, $\abs{R+w}\ge\max(R,\abs w)$ and
$\Re\frac{Rw}{R+w}=R\frac{R\Re w+\abs w^2}{\abs{R+w}^2}\ge0$; when $\abs w\ge R$ the last expression is at
least $R/4$, and $\abs{w(x+iy)}\ge\abs{\sinh(\pi x)}$ covers $\abs x\ge D$. Also
$\Re\cosh(\pi z/2)\ge2^{-1/2}\cosh(\pi x/2)$. The normalizer and its reciprocal are bounded
uniformly (on $\abs x\le1$ use $Rw/(R+w)\le w$). Each derivative of fixed positive order of
$Rw/(R+w)=R-R^2/(R+w)$ is $O_j(\exp(C_j\abs x))$ on the strip uniformly in $R$, and the
double-exponential first factor absorbs it. This gives the pointwise bounds and uniform $L^1$
bounds for derivatives of each fixed order on both edges of the strip. Shift the Fourier
contour to $\Im z=-\operatorname{sgn}(v)/2$ and integrate by parts. For fixed $R$ the kernel is
analytic and rapidly decreasing on $\abs{\Im z}\le\tfrac34$, since the nearest zeros of $R+w$ have
$\abs{\Im z}=1$; this justifies the shift and the use of \eqref{rec:eq:gwstrip}.
\end{proof}

\begin{lemma}[Frequency response; \stU]\label{rec:lem:lefreq}
One has $\widehat h(0)=0$, $\widehat h(u)=1$,
$\widehat h(v)=\frac vue^{-i(v-u)T}\widehat k_R(v-u)\recsinc(b(v-u))^m$, and
\[
 \frac{\sqrt p}u\Bigl(\sum_{n\ge2,\,n\ne p}\frac{\Lambda(n)}{\sqrt n}\abs{\widehat h(\log n)}
 +\sum_{n\ge2}\frac{\Lambda(n)}{\sqrt n}\abs{\widehat h(-\log n)}\Bigr)\ll2^{-m}.
\]
\end{lemma}

\begin{proof}
Differentiation, translation and convolution give the formula. For $p/2\le n\le2p$, $n\ne p$,
$b\abs{\log(n/p)}\ge4\abs{n-p}$, the factors $(\sqrt p/\sqrt n)(\Lambda(n)/u)(\log n/u)$ are bounded,
and $\abs{\widehat k_R}\le1$, so these terms contribute $O(\sum_{j\ge1}(4j)^{-m})$. In the remaining
sums $\abs{v-u}\ge\log2$, and before the sinc factors
$\sum_n\frac{\Lambda(n)}{\sqrt n}\frac{\log n}ue^{-\abs{\log n-u}/2}(1+\abs{\log n-u})^{-4}\ll\sqrt p$: for
$n\le p$ bound it by $p^{-1/2}\sum_{n\le p}\Lambda(n)$, and on $2^jp<n\le2^{j+1}p$ by
$O(\sqrt p(1+j)^{-3})$ by Chebyshev; negative frequencies have the same majorant since
$\abs{-\log n-u}\ge\abs{\log n-u}$. Two sinc factors give $O(p^{-2})$ and each further one is at most
$1/2$.
\end{proof}

\begin{proof}[Proof of \cref{rec:thm:localcoeff}]
In \eqref{rec:eq:gwstrip}, $\widehat h(0)=0$ kills the conductor term, and \cref{rec:lem:lefreq} leaves
the positive frequency $\log p$, with value $\chi(p)\log p/\sqrt p$, up to $O(2^{-m})$ after
normalization. \emph{Unseen zeros.} The density $g$ is supported in $[-B,B]$ with
$\norm g_\infty\le(2b)^{-1}$, so $\abs K+\abs{K'}\ll b^{-1}$ on the closed strip, and if
$d_x=\abs{x-T}-B\ge D$, \cref{rec:lem:lekernel} gives
$\abs{K(x+iy-T)}+\abs{K'(x+iy-T)}\ll_Mb^{-1}e^{-R/4}e^{-Md_x}$ for $\abs y\le\tfrac12$. The modulation
costs at most $\sqrt p$ on the strip. A unit interval at height $y$ contains
$O(\log(q_\chi(\abs y+3)))$ zeros \cite{BennettMartinOBryantRechnitzer2021Counting}; on the $j$th unit
interval outside the window $d_x\ge B+D+j$ and the count is $O(\ell+j+1)$, since $H<T$. The omitted
contribution is therefore $\ll\frac p{bu}e^{-R/4}\sum_{j\ge0}(\ell+j+1)e^{-(B+D+j)}\ll e^{-B}$, and the
argument keeps $\Re\rho$ throughout. \emph{Parity and poles.} Reflection for $\psi$ at
$z=\tfrac14+it/2$ gives $G_1^\circ-G_0^\circ=\pi\operatorname{sech}(\pi t)$. As $\int\abs h\ll1$ and
$\operatorname{sech}(\pi t)\ll e^{-\pi T/4}$ off $\abs t\le T/4$, where the box support is at distance at least
$T/4$, the parity error is $O(e^{-cT})$. At $z=\pm i/2$ the distance from the support is at least
$T/2$, so the pole readings are $O(e^{-T/2})$; with $M=4$ these errors are $O(2^{-m})$ since
$T\ge16pm$.
\end{proof}

\begin{theorem}[Identification and factorization from one window; \stG]\label{rec:thm:localid}
There are absolute constants $C_0,C_1$. Fix $d\ge1$, put $U=d^2\log^2(2Q)\log(2d)$ and
$H=C_0U+D(Q,T)$. If $T\ge C_1U$, the multiset of ordinates in $[T-H,T+H]$ of the zeros of
$\prod_{j=1}^kL(s,\chi_j)$, $0\le k\le d$, $q_{\chi_j}\le Q$, determines $k$ and the multiset
$\set{\chi_1,\dots,\chi_k}$ of primitive characters, repetitions allowed, from the total
multiplicities alone.
\end{theorem}

\begin{proof}
Subtract two putative products as multisets of factors, leaving distinct primitive
$\chi_1,\dots,\chi_s$, $s\le2d$, with nonzero integer differences $a_j$, $\norm a_1\le2d$. Take $X=Ad^2\log^2(2Q)$
with $A$ large and a smooth $W\ge0$ supported in $(1,2)$ with $\int W=1$. By \cref{rec:lem:gram},
\[
 \frac1X\sum_p(\log p)W(p/X)\chi_i(p)\overline{\chi_j(p)}=\mathbf 1_{i=j}+O\bigl(\log(2Q)/\sqrt X\bigr),
\]
so the Gram matrix
differs from the identity entrywise by at most $1/(2s)$, and the weighted mean square of
$\sum_ja_j\chi_j(p)$ is at least $\tfrac12\norm a_2^2$. The total weight divided by $X$ is bounded, so some prime $X<p<2X$ has
\begin{equation}\label{rec:eq:localsep}
 \Bigl\lvert\sum_ja_j\chi_j(p)\Bigr\rvert\ge c_*\norm a_2\ge c_*>0,
\end{equation}
with $c_*$ absolute, also when $p$ ramifies in some factors. Take $m=\ceil{A'\log(2d)}\ge4$ with
$2dC2^{-m}<c_*/2$; for $C_0,C_1$ large the window of \cref{rec:thm:localcoeff} for this $p$ lies in
$[T-H,T+H]$ and $T\ge16pm+D+2$. To allow different degrees use only the zero sum in $\mathcal E_p$:
integrating by parts, $\bigl\lvert\frac{\sqrt p}u\int hG_0^\circ\bigr\rvert\le\frac{\sqrt p}{u^2}\int K(t-T)\abs{G_0^{\circ\prime}(t)}\dd t
\ll\sqrt p/(Tu^2)$, since $G_0^{\circ\prime}(t)\ll(1+\abs t)^{-1}$. Each zero-only reading differs from the
corresponding sum of character values by at most $d(C2^{-m}+C'\sqrt p/(Tu^2))<c_*/2$ after enlarging
$A',C_1$, and equal spectra in the window would contradict \eqref{rec:eq:localsep}, including against
the empty product.
\end{proof}

For $d=1$ the conductor, parity, order and all Euler factors are determined, and a finite decoder
enumerates candidates. A partial matching of two lists pairs occurrences bijectively; let $N$ be
the number of unpaired occurrences and $S$ the sum of the displacements of the paired ones.

\begin{corollary}[Local matching cost; \stG]\label{rec:cor:localmatch}
For distinct primitive characters of conductor at most $Q$ and the window of \cref{rec:thm:localid}
with $d=1$, every partial matching of their lists satisfies
$S+N/\log(2+\log(2Q))\gg\log(2Q)$. With only exact coincidences paired, the symmetric difference
counted with multiplicity is $\gg\log(2Q)\log(2+\log(2Q))$.
\end{corollary}

\begin{proof}
Use the prime of \eqref{rec:eq:localsep}, $p\asymp\log^2(2Q)$, and the statistic with its sum extended
to the whole window; the tail estimate covers the complement of the larger window as well. The
readings differ by at least $c_*/2$, while $\norm h_\infty\ll p^{-1}$ and $\norm{h'}_\infty\ll(\log p)/p$ on
the real axis. Hence $1\ll(N/u+S)/\sqrt p$.
\end{proof}

At $T=1000$ the statistic, computed in floating point from self-generated zero lists of $\zeta$,
$L(s,\chi_{-4})$ and the quartic character modulo $5$ with $\chi(2)=i$, reproduces $\chi(p)$ for
$p=2,3,5,7$ with errors below $10^{-4}$, and four zero-only readings select
$\zeta(s)L(s,\chi)^2$ among $20$ candidate products, with residual $0.000203$ against $1.000077$ for
the next candidate;
\begin{center}
\artifact{local_euler_recovery/numerical/local_readout.py}\\
\artifact{local_euler_recovery/numerical/output.txt}
\end{center}

\paragraph{A binomial cutoff and full rank at the endpoint.}
Fix distinct primitive characters $\chi_1,\dots,\chi_d$ of conductors $q_j$, $Q=\lcm(q_j)$, and put
$m(\rho)=(\ord_\rho L(s,\chi_j))_j$ for a nontrivial zero of at least one of them, and
$m(\gamma)=m(\tfrac12+i\gamma)$ under \cref{src:thm:dgrh}. Write $\log_3T=\log\log\log T$. Constants
may depend on the family.

\begin{lemma}[A cutoff across the critical strip; \stU]\label{rec:lem:binomkernel}
Fix $W\ge1$ and an integer $J\ge4$. Let $k_M(z)=c_M\operatorname{sech}^2(z)\bigl(1+M^{-1}\cosh(\pi z)\bigr)^{-M}$ with
$\int_\R k_M=1$, and let $K=K_{M,W}$ be the convolution of $k_M$ with $J$ uniform probability
densities on $[-W/J,W/J]$. Uniformly in integers $M\ge2$,
\[
 \abs{K(x+iy)}\ll_{W,J}e^{-2\abs x},\qquad \abs{K(x+iy)}\ll_{W,J}2^{-M/2}e^{-2\abs x}\quad
 \bigl(\abs x\ge W+\pi^{-1}\operatorname{arsinh}M\bigr)
\]
for $\abs y\le\tfrac12$, and $\abs{\widehat K(\xi)}\ll e^{-\abs\xi/2}\min\set{1,(J/(W\abs\xi))^J}$. The function $K$
is even, nonnegative on $\R$, admissible in \eqref{rec:eq:gwstrip}, and $\widehat K(0)=1$. The same
bounds hold for $K^{(j)}$, $0\le j\le J-4$, with constants depending on $j,W,J$.
\end{lemma}

\begin{proof}
The constants $c_M$ are bounded above and below, by comparison with $\operatorname{sech}^2x$ and by integration
over $[-1,1]$ using $(1+y/M)^M\le e^y$. In the closed strip $w=\cosh(\pi z)$ has $\Re w\ge0$ and
$\abs w^2=\sinh^2(\pi x)+\cos^2(\pi y)$, so $\abs{1+w/M}\ge1$, and $\abs{1+w/M}\ge\sqrt2$ when $\abs w\ge M$,
which holds for $\abs x\ge\pi^{-1}\operatorname{arsinh}M$; and $\abs{\operatorname{sech}^2(x+iy)}\ll e^{-2\abs x}$. The nearest
poles of the binomial factor have $\abs{\Im z}=1$ (if $\cosh\pi z=-M$ then $\sin\pi y=0$), so the
kernel is analytic with exponential decay on $\abs{\Im z}\le\tfrac34$, and shifting the Fourier
integral to $\Im z=-\operatorname{sgn}(\xi)/2$ gives $\abs{\widehat k_M(\xi)}\ll e^{-\abs\xi/2}$. Averaging translates
over the compact support $[-W,W]$ gives the pointwise bounds, and the multiplier
$(\sin(W\xi/J)/(W\xi/J))^J$ the last one. For $K^{(j)}$ put one derivative on each of $j$ uniform
factors; each is a signed measure on two endpoints.
\end{proof}

At centre $U$, $T/2\le U\le2T$, take $M=\ceil{10\log\log T}$ and $H=W+\pi^{-1}\operatorname{arsinh}M=\pi^{-1}\log_3T+O_W(1)$.
Multiplying $K(z-U)$ by a fixed finite sum of exponentials $e^{i\lambda z}$ costs a fixed factor on
the strip, and by the lemma and the unit-interval count the zeros with $\abs{\Im\rho-U}>H$
contribute $O_W(2^{-M/2}\log(2T))=o(1)$; the pole terms are $O_W(e^{-2U})$; and Stirling's formula on
$\abs{t-U}\le U/2$ gives
\begin{equation}\label{rec:eq:binomgamma}
 \int h(t)G_\chi(t)\dd t=\log\bigl(qU/(2\pi)\bigr)\int h(t)\dd t+O_W(T^{-1})
\end{equation}
for these tests. For $t\mapsto K(t-U)e^{i\omega t}$ the transform at $\xi$ is $e^{iU(\omega-\xi)}\widehat K(\xi-\omega)$,
so, for any fixed finite list of frequencies, the prime sum in \eqref{rec:eq:gwstrip} tends, as
$W\to\infty$, to its terms with $\pm\log n$ in the list, uniformly in $U$ and $M$; the large-$n$
majorant $C_\omega\log n/(n(1+\log n)^J)$ is summable. The primes are fixed first, then $W$, then $T$.

\begin{theorem}[Local independence of multiplicity vectors; \stU]\label{rec:thm:localrank}
There are $C,T_0$ such that for $T\ge T_0$ the vectors $m(\rho)$, $\abs{\Im\rho-T}\le\pi^{-1}\log_3T+C$,
indexed by distinct complex zeros, span $\C^d$. In particular every nonzero complex linear
combination of the $L'/L(s,\chi_j)$ has a pole in this height range. For a family closed under
conjugation the same holds for $T\le-T_0$.
\end{theorem}

\begin{proof}
The induced characters on $(\Z/Q\Z)^\times$ are distinct, so the value table has rank $d$; choose $d$
independent rows and a prime $p_i$ in each corresponding class, so that
$B_{ij}=\chi_j(p_i)\log(p_i)/\sqrt{p_i}$ is invertible. For $\lambda=\log p_i$ use
$h_i(t)=K(t-T)\bigl(e^{i\lambda t}-e^{i\lambda T}\widehat K(\lambda)\bigr)$, whose integral vanishes by evenness of $K$,
so that \eqref{rec:eq:binomgamma} is $O_W(1/T)$ despite the growing common gamma term. Frequency
selection gives $2\pi\sum_{\abs{\Im\rho-T}\le H}h_i(z_\rho)m_j(\rho)=-B_{ij}+o_W(1)+o_{T;W}(1)$, the unmodulated
correction having a bounded prime sum and coefficient $o_W(1)$, and the tail including zeros off
the line. Choose $W$, then $T_0$, to make the error smaller than half the least singular value of
$B$; the left matrix has rank $d$ and its rows lie in the span of the $m(\rho)$. At a zero the residue
of $\sum_jv_jL'/L(s,\chi_j)$ is $v\cdot m(\rho)$. Conjugation gives negative heights.
\end{proof}

This sharpens \cref{zer:thm:rank}, whose windows have length $K_q\log T$, and the radii
$b\log_3T+O_{q,b}(1)$, $b>1/\pi$, obtained earlier for the characters of one modulus. The
coefficient $1/\pi$ is reached with an additive constant; its optimality is open.

\begin{lemma}[Distinguishing primes; \stU]\label{rec:lem:frame}
For every $D_0>0$ there is a finite set $P$ of primes coprime to $Q$ with
$\sum_{p\in P}\frac{(\log p)^2}p\bigl\lvert\sum_ju_j\chi_j(p)\bigr\rvert^2>D_0^2\norm u_2^2$ for all nonzero $u\in\C^d$.
\end{lemma}

\begin{proof}
In each reduced class modulo $Q$ the sum of $1/p$ diverges; choose primes so that
\[
 \sum(\log p)^2/p>D_0^2/\varphi(Q)
\]
in every class, and use orthogonality.
\end{proof}

\begin{theorem}[Both signs at the endpoint; \stG]\label{rec:thm:bothsigns}
Let $d\ge2$. There are $C,c>0$ and $T_0$ such that for $T\ge T_0$, every real unit vector
$u\perp\mathbf 1$ and $s\in\set{\pm1}$,
\[
 \sum_{\abs{\gamma-T}\le\pi^{-1}\log_3T+C}\bigl(s\,u\cdot m(\gamma)\bigr)_+\ge c,
\]
and for each fixed $b>1/\pi$ and large $T$,
$\sum_{\abs{\gamma-T}\le b\log_3T}(s\,u\cdot m(\gamma))_+\gg_b\log_3T$. In particular, for every nonzero $n\in\Z^d$
with $\sum_jn_j=0$, $\prod_jL(s,\chi_j)^{n_j}$ has a zero and a pole in every window of the first
kind, after all cancellations, and $\gg\norm n_2\log_3T$ of each in the windows of the second kind.
\end{theorem}

\begin{proof}
Put $c_u=\sum_ju_j\log q_j$ and $b_p(u)=\frac{\log p}{\sqrt p}\sum_ju_j\chi_j(p)$, and fix $P$ by
\cref{rec:lem:frame} with $\norm{b(u)}_2>\norm{(\log q_j)_j}_2+3$ for every such $u$. For $a=(a_0,(a_p)_{p\in P})$
of norm one let $A_a(t)=a_0+\sum_pa_pe^{it\log p}$ and $h_{a,U}(t)=K(t-U)\abs{A_a(t)}^2$, extended as
$A_a(z)\overline{A_a(\bar z)}$. The frequencies of $\abs{A_a}^2$ are $0$ (coefficient one), $\pm\log p$ (from the
constant term) and $\log(p/r)$ for distinct $p,r\in P$, which are not signed prime-power logarithms;
so $\int h_{a,U}=1+o_W(1)$. Combine \eqref{rec:eq:gwstrip} over $j$ with coefficients $u_j$: by
\eqref{rec:eq:binomgamma} the gamma contribution is $c_u\int h_{a,U}+O_W(T^{-1})$, the common term
$\log(U/2\pi)$ cancelling because $\sum u_j=0$, and
\[
 2\pi\sum_{\abs{\gamma-U}\le H}h_{a,U}(\gamma)\,u\cdot m(\gamma)=c_u-2\Re\Bigl(\overline{a_0}\sum_pa_pb_p(u)\Bigr)+o_W(1)+o_{T;W}(1),
\]
uniformly in $T/2\le U\le2T$, $a$ and $u$. By \cref{src:thm:dgrh} each $h_{a,U}(\gamma)$ is nonnegative.
Take $a_0=1/\sqrt2$ and $a_p=\pm\overline{b_p(u)}/(\sqrt2\norm{b(u)}_2)$: the main terms are $c_u\mp\norm{b(u)}_2$,
of either sign with margin greater than three. Since $\abs{A_a}^2\le1+\abs P$ and $K\ge0$,
\begin{equation}\label{rec:eq:weightedsign}
 \sum_{\abs{\gamma-U}\le H}K(\gamma-U)\bigl(s\,u\cdot m(\gamma)\bigr)_+\ge c_0=\frac1{2\pi(1+\abs P)}\qquad(s=\pm1),
\end{equation}
and $U=T$ with $\norm K_\infty<\infty$ gives the first claim. For the second, with $R=b\log_3T$ and
$M,W,H$ fixed at their values for $T$, integrate \eqref{rec:eq:weightedsign} over $\abs{U-T}\le R-H$:
each ordinate receives total weight at most $\int K=1$, so the sum is at least $2c_0(R-H)$.
\end{proof}

\Cref{zer:thm:short} gives, for the same unshifted families, mass $\gg H\log\log H$ on intervals of
length $H\ge(2/\pi+\eps)\log_3T$; the theorem above removes the $\eps$ at the price of a bounded
mass, and its tests survive translation of the functions (\cref{rec:thm:moving}).

\begin{corollary}[Exact balancing; \stG]\label{rec:cor:balance}
In each window of the first kind in \cref{rec:thm:bothsigns} the points
$v_\gamma=m(\gamma)/\sum_jm_j(\gamma)$ have a convex hull in the hyperplane $\sum_jx_j=1$ containing the
ball of radius $c_1/\log T$ about $\mathbf 1/d$, and some at most $d$ of these ordinates with
positive integer weights have a weighted sum of multiplicity vectors with all coordinates equal.
\end{corollary}

\begin{proof}
The weight $\sum_{\abs{\gamma-T}\le H}K(\gamma-T)\sum_jm_j(\gamma)$ is $O(\log T)$, so
\eqref{rec:eq:weightedsign} gives $\max_\gamma u\cdot v_\gamma\gg1/\log T$ for every unit $u\perp\mathbf 1$;
separation gives the ball, Carath\'eodory's theorem a rational convex combination equal to
$\mathbf 1/d$ on at most $d$ points, and clearing denominators the integer weights.
\end{proof}

\begin{corollary}[A cost for moving or deleting zeros; \stG]\label{rec:cor:binomtransport}
For two fixed distinct primitive characters and $b>1/\pi$, charge $1$ for each unmatched
occurrence and $\min(1,\abs{\gamma-\gamma'})$ for each matched pair in a partial matching of their
zero lists in $[T-b\log_3T,T+b\log_3T]$. Every matching costs $\gg_b\log_3T$.
\end{corollary}

\begin{proof}
Take $u=(1,-1)/\sqrt2$ and the positive test above; by the tail bound the indicator of $[-H,H]$
may be replaced by a taper $\eta_H$, equal to one on $[-H,H]$ and linear to zero on
$H\le\abs x\le H+1$. For $J'=[T-R+H+1,T+R-H-1]$, $R=b\log_3T$, the function
$f(t)=\abs{A_a(t)}^2\int_{J'}K(t-U)\eta_H(t-U)\dd U$ is nonnegative, supported in the window, and has
supremum and Lipschitz constant bounded independently of $T$. Integrating the signed estimate over
$J'$ gives $\sum f(\gamma_1)-\sum f(\gamma_2)\gg\abs{J'}\gg\log_3T$, while a matched pair contributes at most
$\max(2\norm f_\infty,\operatorname{Lip}f)\min(1,\abs{\gamma_1-\gamma_2})$ and an unmatched zero at most $\norm f_\infty$.
\end{proof}

\paragraph{Translated functions and coalescing shifts.}
Consider $L(s+i\tau_j,\chi_j)$ with real shifts $\tau=(\tau_j)$ in a fixed compact set $E$, allowing
repeated primitive characters but with $(\chi_j,\tau_j)\ne(\chi_k,\tau_k)$ on $E$ for $j\ne k$; the
multiplicity at $\tfrac12+it$ is $m_{\chi_j}(t+\tau_j)$.

\begin{theorem}[Uniform translated families; \stG]\label{rec:thm:moving}
\Cref{rec:thm:localrank}, \cref{rec:thm:bothsigns} and \cref{rec:cor:balance} hold for these shifted
functions uniformly in $\tau\in E$, the rank statement without \cref{src:thm:dgrh}; for a pair, so does \cref{rec:cor:binomtransport}, and its
lower bound persists after the infimum over all allowed translations.
\end{theorem}

\begin{proof}
Testing the original list with $h(t-\tau_j)$ replaces the prime coefficient by $\chi_j(n)n^{-i\tau_j}$
and the gamma density by $G_{\chi_j}(t+\tau_j)$; poles and errors are uniform for bounded shifts. It
suffices to choose the primes of \cref{rec:lem:frame} uniformly. For fixed $\tau$ the matrix
$G_X(j,k)=\sum_{p\le X,p\nmid Q}\frac{(\log p)^2}p\chi_j(p)\overline{\chi_k(p)}p^{-i(\tau_j-\tau_k)}$ satisfies
$G_X/(\tfrac12\log^2X)\to I_d$, by the prime number theorem in progressions and partial summation
($\int_2^X(\log x)x^{-1-iv}\dd x$ is $\tfrac12\log^2X+O(1)$ for $v=0$ and $O_v(\log X)$ otherwise); where the
characters coincide the shifts differ. The condition $\lambda_{\min}(G_X)>D_0^2$ is open in $\tau$ and
persists as $X$ grows, since $G_X$ increases as a quadratic form, and compactness gives one $X$ for
all of $E$. All later arguments use finitely many frequencies with uniformly bounded coefficients.
\end{proof}

For a primitive $\chi$, $r\ge1$ and real $h\ne0$ put
\[
 F_{r,h}(s)=\prod_{k=0}^rL(s+ikh,\chi)^{(-1)^{r-k}\binom rk},
\]
so that $\zeta(s+ih)/\zeta(s)$ and $\zeta(s+2ih)\zeta(s)/\zeta(s+ih)^2$ are the cases $r=1,2$; at $h=0$ all
factors cancel. Zeros and poles are counted after cancellation.

\begin{theorem}[Uniformity as the shifts coalesce; \stG]\label{rec:thm:coalesce}
Fix $\chi$, $r\ge1$ and $B_0>0$. There are $C,T_0$ independent of $h$ such that for $T\ge T_0$ and
every $0<\abs h\le B_0$, $F_{r,h}$ has a zero and a pole on the critical line in
$\abs{\Im s-T}\le\pi^{-1}\log_3T+C$; if $N_\pm$ count them there, $N_\pm\ge c\abs h^r$, and in
$\abs{\Im s-T}\le b\log_3T$, $b>1/\pi$ fixed, $N_\pm\gg_{\chi,r,B_0,b}\max\set{1,\abs h^r\log_3T}$.
\end{theorem}

\begin{proof}
For a test $f$ put $\mathcal D_{r,h}(f)=h^{-r}\sum_kd_k\sum_\gamma f(\gamma-kh)$, $d_k=(-1)^{r-k}\binom rk$, zeros with
multiplicity; its signed measure has atom $\ord_{1/2+it}F_{r,h}$ at $t$, and
$h^{-r}\sum_kd_kg(x+kh)=\int_{[0,1]^r}g^{(r)}(x+h(v_1+\dots+v_r))\dd v$ controls it although $\abs{d_k/h^r}$ is large.
Choose distinct primes $p_1,p_2\nmid q_\chi$ and put $b_p(h)=\chi(p)\frac{\log p}{\sqrt p}\bigl(\frac{e^{-ih\log p}-1}h\bigr)^r$,
extended continuously to $h=0$. They have no common zero on $[-B_0,B_0]$ (that would make
$\log p_1/\log p_2$ rational), so $\max_{p=p_1,p_2}\abs{b_p(h)}\ge\delta>0$. Use $J=r+4$ in
\cref{rec:lem:binomkernel}, and for the maximizing $p$ the two positive tests
$f=K(t-U)\bigl(1\pm\Re(\alpha e^{it\log p})\bigr)$, $\alpha=\overline{b_p(h)}/\abs{b_p(h)}$. In
$\mathcal D_{r,h}(f)$ the coefficient $\chi(n)$ is multiplied by $((n^{-ih}-1)/h)^r$, of size at most
$(\log n)^r$, the off-frequency majorant $O_p((\log n)^{r+1}/(n(1+\log n)^J))$ is summable, and the
peak gives $\mp\abs{b_p(h)}+o_W(1)$ uniformly in $h$. The gamma density
$h^{-r}\sum_kd_kG_\chi(t+kh)$ is $O(T^{-r})$ near $U$ and bounded globally, since
$G_\chi^{(r)}(t)=O((1+\abs t)^{-r})$; the pole terms are $O(e^{-2U})$. For the truncation take a smooth
taper $\eta_H$ with derivatives through order $r$ bounded independently of $H$ and
$g=(1-\eta_H(t-U))f$: every derivative through order $r$ is $O(2^{-M/2}e^{-2\abs{t-U}})$, and applying the
integral identity before summing over zeros gives $\abs{\mathcal D_{r,h}(g)}=o(1)$ with no factor
$\abs h^{-r}$. Hence the tapered nonnegative tests have readings of opposite signs bounded below,
and since $0\le1\pm\Re(\cdot)\le2$, $\sum_{\abs{t-U}\le H+1}K(t-U)(s\ord F_{r,h})_+\gg\abs h^r$ for $s=\pm1$. The
supremum of $K$ gives both signs in the endpoint window, integration over $U$ the larger window,
and each positive count is an integer.
\end{proof}

\begin{theorem}[Microscopic steps; \stG]\label{rec:thm:micro}
Fix $r\ge1$, $\chi$ and $0<\kappa<2\pi$. For any $C>\frac\pi2+\frac{\pi\kappa}{2\pi-\kappa}$ and all large $T$, $F_{r,h}$ has
a zero and a pole in $\abs{\Im s-T}\le C/\log\log T$ simultaneously for all $0<\abs h\le\kappa/\log T$;
in particular $[T,T+1]$ contains $\gg_{\kappa,\chi,r}\log\log T$ zeros and as many poles, uniformly in these
steps.
\end{theorem}

\begin{proof}
Put $\lambda=\log T$, $\ell=\log\log T$. Under \cref{src:thm:dgrh}, $\abs{S_\chi(t)}\le(\tfrac14+o(1))\lambda/\ell$ for
$\abs{t-T}=O(1)$ \cite{CarneiroChandeeMilinovich2015Note}, so for fixed $a>0$ and $\abs{x-T}\le2$,
$\bigl\lvert N_\chi((x,x+a/\ell))-\frac a{2\pi}\frac\lambda\ell\bigr\rvert\le(\tfrac12+o(1))\frac\lambda\ell$. A central interval of length
$a/\ell$, $a>\pi$, contains a zero $\gamma_0$. For $A>\pi\kappa/(2\pi-\kappa)$, each of the two adjacent intervals
of length $A/\ell$ from $\gamma_0$ has at least $A\lambda/(\kappa\ell)-O(1)$ sites of the grid $\gamma_0+nh$ and at most
$(A/(2\pi)+\tfrac12+o(1))\lambda/\ell$ zeros, and the first coefficient is larger; so each side has at least
$r$ empty sites. With $v(n)=m_\chi(\gamma_0+nh)$, $v(0)>0$, the order-$r$ divided difference on $r$
negative empty indices and $0$ is positive by the Lagrange formula, and on $r-1$ negative empty
indices, $0$ and one positive empty index it is negative. Every divided difference on integer
nodes is a convex combination of divided differences on $r+1$ consecutive sites in their span
(interpolate through a missing interior integer $y$ with weights $(y-x_0)/(x_r-x_0)$ and
$(x_r-y)/(x_r-x_0)$ and induct), and on consecutive sites it is $\Delta^rv(n)/r!$, the order of
$F_{r,h}$ at the corresponding ordinate. All these ordinates lie within $(a/2+A)/\ell$ of $T$, and
$a/2+A<C$ is possible. For the unit interval use disjoint windows with a slightly larger $\kappa'<2\pi$.
\end{proof}

At spacing exactly $2\pi/\log T$ a lattice is invariant under the step, so a density comparison
alone gives nothing at $\kappa=2\pi$.

\begin{corollary}[Nonintegral multiples of the mean spacing; \stG]\label{rec:cor:nonres}
Fix $K_0,\delta>0$. If $\nu=\abs h\log T/(2\pi)$ satisfies $0<\nu\le K_0$ and $\operatorname{dist}(\nu,\set{1,2,\dots})\ge\delta$,
then $F_{r,h}$ has a zero and a pole in $\abs{\Im s-T}\le C_{\chi,r,K_0,\delta}/\log\log T$, uniformly in $h$.
\end{corollary}

\begin{proof}
For $\nu\le\tfrac12$ use \cref{rec:thm:micro}. Otherwise take consecutive half-open bins of length
$s=\abs h$, $r+1$ blocks of $k=\floor{A/(s\ell)}$ bins from $T$, and bin counts $a_n$. The counting formula
gives block means $\nu+O(1/A)+o(1)$, the endpoint error $O(\log T/\ell)$ being divided by
$k\asymp A\log T/(\nu\ell)$; for $A$ large each block mean lies strictly between $\floor\nu$ and $\ceil\nu$, so
each block has counts below and above $\nu$. Choosing one bin per block with $a_n-\nu$ of alternating
signs makes every term of the order-$r$ Lagrange divided difference of one strict sign, and the
reverse choice the other. The consecutive-site identity gives $\Delta^ra_n$ of both signs, and for
$h=s>0$, $\Delta^ra_n$ is the signed divisor count of $F_{r,h}$ in bin $n$, translation by $ks$ preserving the
bins; for $h=-s$ reverse the grid.
\end{proof}

The archive checks the binomial identity on $456$ node choices, translated bin identities, and,
for $\chi=1$, $r=2$, $\abs h\le1$, the two-prime lower bound $\sum_{p=2,3}\abs{b_p(h)}^2>\tfrac12$
and, for $\zeta$ and $\chi_{-4},\chi_8,\chi_{-8}$, that the twelve primes $3,\dots,47$ of \cref{rec:lem:frame} already
exceed the conductor bias $\sqrt6\log2$ in every balanced direction:
\begin{center}
\artifact{translated_dirichlet_zero_divisors/checks/translation_check.txt}\\
\artifact{translated_dirichlet_zero_divisors/checks/lattice_check.txt}
\end{center}

\paragraph{The zeros of a differentiated completion.}
For primitive $\chi$ of conductor $q$ and parity $a$ let $A(s)=(q/\pi)^{(s+a)/2}\Gamma((s+a)/2)$, with
the factor $s(s-1)/2$ included for $q=1$, and $\Xi_\chi=AL(s,\chi)$. Let $P(X)=\prod_{j=1}^d(X-i\alpha_j)$,
$\alpha_j\in\R$, be monic ($P=1$ allowed), $F=P(\dd/\dd s)\Xi_\chi$ and, for a prime $p$, $h=\log p$ and
\[
 S_p(T,H)=\frac{\sqrt{2\pi p}}H\sum_\gamma\exp\Bigl(-\frac{(\gamma-T)^2}{2H^2}+ih\gamma\Bigr),\]\[
 z=\tfrac12\log\frac T{2\pi}+\frac{i\pi}4,\quad c=\tfrac12\log q,\quad Q_P(X)=P(X+c),
\]
the sum over the zeros $\tfrac12+i\gamma$ of $F$ with multiplicity.

\begin{lemma}[Zeros after differentiation; \stG]\label{rec:lem:interlace}
$F$ has order at most one, all its zeros lie on $\Re s=\tfrac12$, a unit interval contains
$O_{\chi,P}(\log(2+\abs u))$ of their ordinates, and $F(s)=(-1)^d\epsilon_\chi\overline{F(1-\bar s)}$.
\end{lemma}

\begin{proof}
Dividing $\Xi_\chi(\tfrac12+iw)$ by a square root of $\epsilon_\chi$ gives a real entire $f$ of order at most
one with only real zeros, by \cref{src:thm:dgrh}. Its genus-one product gives
$\Im(f'/f)(w)=-\Im(w)\sum_tm_t/\abs{w-t}^2$, so $(\dd/\dd w+\alpha)f$ has only real zeros for real $\alpha$: a zero of
order $m$ keeps order $m-1$, and $f'/f$ decreasing from $+\infty$ to $-\infty$ on each gap gives one more
simple zero. Interval counts change by at most one per step, and $D-i\alpha$ corresponds to
$-i(\dd/\dd w+\alpha)$; induct. The last identity follows by conjugation.
\end{proof}

\begin{theorem}[Prime-frequency formula; \stG]\label{rec:thm:primefreq}
For fixed $\chi,P,p$ there is $\eta_p>0$ such that uniformly for $1\le H\le T^{2/3}$,
\[
 S_p(T,H)=-h\chi(p)\frac{Q_P(z-h)}{Q_P(z)}+O_{\chi,P,p}\bigl(T^{-1/5}+(\log T)e^{-\eta_pH^2}\bigr),
\]
including $p\mid q$. With $H=\log\log T$ the error is $o((\log T)^{-k})$ for every $k$.
\end{theorem}

\begin{proof}
In the Banach algebra of Dirichlet coefficient sequences with norm $\sum\abs{u_n}n^{-2}$, put $g=z+c$,
$U_g(s)=\sum_n\chi(n)P(g-\log n)n^{-s}$ and $V_g(s)=-\sum_n\chi(n)(\log n)P(g-\log n)n^{-s}$. Since
$U_g/g^d\to L(s,\chi)$ in norm with error $O(1/\abs g)$ and $L^{-1}$ lies in the algebra, a Neumann series
gives $\norm{U_g^{-1}}=O(\abs g^{-d})$, so $B_g=V_g/U_g=\sum b_n(g)n^{-s}$ has bounded norm, $b_1=0$, and
comparison of coefficients at $p$ in $B_gU_g=V_g$ gives $b_p(g)=-(\log p)\chi(p)P(g-\log p)/P(g)$. For
$s=2+it$, $\abs{t-T}\le T^{3/4}$, Stirling gives $A'/A=g+O(T^{-1/4})$ and $(A'/A)^{(j)}=O(T^{-j})$, hence
$A^{(k)}/A=g^k+O(T^{-1/4}(1+\abs g)^k)$; replacing $\dd/\dd s$ acting on $A$ by $g$ in $F/A$ and $F'/A$ gives $U_g$ and
$gU_g+V_g$ up to $O(T^{-1/4}(1+\abs g)^{d+1})$, and division yields $(F'/F)(2+it)=g+B_g(2+it)+O(T^{-1/5})$.
Apply the residue theorem on $\sigma=2,-1$ to $(F'/F)(s)p^s\exp((s-\tfrac12-iT)^2/(2H^2))$, with horizontal
edges at heights $\gg1/\log(2+\abs t)$ from all ordinates (\cref{rec:lem:interlace}), where
$F'/F=O(\log^2(2+\abs t))$; reflection gives $(F'/F)(-1+it)=-\overline{(F'/F)(2+it)}$. A term of frequency $v$ on
the edge $\sigma$ integrates to $e^{ivT}e^{-H^2v^2/2-(\sigma-1/2)v}$; the only resonance is $n=p$ on the right
edge, and all other frequencies have $\abs v\ge\log(1+1/p)=\delta$, costing $O((\log T)e^{-\delta^2H^2/4})$.
\end{proof}

\begin{theorem}[One-prime reconstruction; \stG]\label{rec:thm:operatorinverse}
For a prime $p\nmid q$, the values $S_p(T,\log\log T)$ along any unbounded sequence of $T$ determine
$\chi(p)$ and the monic $P$, and, if $\deg P>0$, the conductor $q$. Writing the normalized expansion as
$1+\sum_jr_jz^{-j}$, $r_1=-dh$, $r_2=\binom d2h^2+ha_1$ with $Q_P=X^d+a_1X^{d-1}+\dots$, and
\[
 q=\exp\Bigl(\frac2{dh}\Re\Bigl[r_2-\binom d2h^2\Bigr]\Bigr).
\]
Terms through $z^{-d-1}$ determine $P$, and for each $d\ge1$ terms through $z^{-d}$ can agree for distinct
operators of degree $d$, even with $q$ and $\chi$ fixed.
\end{theorem}

\begin{proof}
Divide by the limit $-h\chi(p)\ne0$ and read the coefficients successively. The roots of $Q_P$ are
$-c+i\alpha_j$, so $\Re a_1=dc$. If monic $Q_1,Q_2$ of degree $d$ first differ by $\Delta a_j$ in position $j$, then
$\frac{Q_1(z-h)}{Q_1(z)}-\frac{Q_2(z-h)}{Q_2(z)}=jh\Delta a_jz^{-j-1}+O(z^{-j-2})$, by taking logarithms near infinity, which
recovers $a_j$ from $r_{j+1}$ inductively. For sharpness let $f(x)=\prod_{j=1}^d(x-j)$; for small real $t\ne0$,
$f+t$ has $d$ simple real roots, and $P_1(X)=i^df(X/i)$, $P_2(X)=i^d(f(X/i)+t)$ differ by a constant,
so they agree through $z^{-d}$ and differ at $z^{-d-1}$.
\end{proof}

For the quadratic character modulo $5$, $p=2$ and $P=X^2-iX+2$, the limit is $\log2$, the first
correction gives degree two, $r_2$ gives $q=5$, and the third correction gives $a_2$. For $P=X$ and
$q=1$ the direct sums at $p=3$ over $626$ and $1572$ zeros of $\xi'$ agree with $-\log3(1-\log3/z)$ to
about $10^{-4}$:
\begin{center}
\artifact{dirichlet_prime_frequency_inverse/check_algebra.py}\\
\artifact{dirichlet_prime_frequency_inverse/zero_example_output.txt}
\end{center}

\begin{corollary}[Local distinction of operators; \stG]\label{rec:cor:operatorlocal}
For each $k$ there are $A,B$ such that $S_p$ with $H=A\sqrt{\log\log T}$, restricted to
$\abs{\gamma-T}\le B\log\log T$, equals the main term up to $o((\log T)^{-k})$. Two distinct pairs
$(\chi_1,P_1)\ne(\chi_2,P_2)$ have different zero multisets in $[T-C\log\log T,T+C\log\log T]$ for all
large $T$, and $[X,2X]$ contains $\gg X/\log\log X$ unmatched occurrences; constants depend on the pairs.
\end{corollary}

\begin{proof}
The weights $e^{-(\gamma-T)^2/(4H^2)}$ have sum $O(H\log T)$, so the omitted tail is
\[
 O\bigl((\log T)^{1-B^2/(4A^2)}\bigr);
\]
take $\eta_pA^2$ and $B^2/(4A^2)$ larger than $k+2$. Distinct characters
differ at some prime, where the constant terms differ; for equal characters distinct monic
operators differ at a finite order by \cref{rec:thm:operatorinverse}. Pack disjoint windows in
$[X,2X]$.
\end{proof}

If $d$ is known, interpolating the truncated sums at the $d+2$ centres $T^j$, $1\le j\le d+2$, through
degree $d+1$ in $1/z$ (the nodes $\log T/z(T^j)$ tend to $2/j$, so the Vandermonde inverses are bounded)
estimates the coefficient of $z^{-k}$ with error $O((\log T)^{k-d-2})$, and for $d>0$ the conductor
estimate eventually rounds to $q$; no numerical threshold is claimed.

For $P=1$ the window here is much longer than that of \cref{rec:thm:localid}; the content of the
corollary is the operator. Without \cref{src:thm:dgrh} the same holds for the sum
$\frac{\sqrt{2\pi}}H\sum_{F(\rho)=0}p^\rho\exp((\rho-\tfrac12-iT)^2/(2H^2))$ over complex zeros (\stU): the zeros of
$\Xi_\chi$ lie in $0\le\Re s\le1$, for a real entire $f$ of order at most one with zeros in
$\abs{\Im w}\le\tfrac12$ the paired product gives $\Im(f'/f)(w)=-\sum_v(\Im w-\Im v)/\abs{w-v}^2$, of strict sign
outside the strip, a zero-free intermediate function is excluded by the right-edge lower bound
$\abs{P(D)\Xi_\chi(2+it)/A(2+it)}\gg(\log\abs t)^d$, and Jensen's formula on disks of radius $5$ about
$2+iu$ gives the unit-interval count; the rest of the proof is unchanged.

\subsection{Weights and multiplicity vectors from poles of prime series}\label{rec:sec:poles}
\sources{prime_pole_windows; dirichlet_zero_recovery; large_kernels_and_dirichlet_zeros; dirichlet_inverse_problems; local_euler_recovery; translated_dirichlet_zero_divisors}

Fix $q\ge2$, $G=(\Z/q\Z)^\times$, a weight $f:G\to\C$, and put
\[
 F_f(s)=\sum_{p\nmid q}\frac{f(p\bmod q)\log p}{p^s},\qquad G^m=\set{a^m:a\in G},\qquad U_mf(a)=f(a^m).
\]
With $c_\chi(g)=\abs G^{-1}\sum_ag(a)\overline{\chi(a)}$, let $n_\psi(\rho)=\ord_\rho L(s,\psi^*)$ for
$\rho$ in the open strip and $Z_g(\rho)=\sum_\psi c_\psi(g)n_\psi(\rho)$. The question is what the
poles and residues of $F_f$, or of the progression series, determine about $f$ and about the
multiplicity vectors $(n_\psi(\rho))_\psi$, which may have several nonzero entries at a common
zero.

\begin{proposition}[Residues on power lines; \stG]\label{rec:prop:powerline}
$F_f$ continues meromorphically to $\Re s>0$ with at most simple poles. For $\rho=\tfrac12+i\gamma$,
$\gamma\ne0$, and $m\ge1$,
\[
 \Res_{s=\rho/m}F_f(s)=-\frac{\mu(m)}m\,Z_{U_mf}(\rho),
\]
and every nonreal pole in $\Re s>0$ has this form. In particular only $f|_{G^m}$ enters the line
$\Re s=1/(2m)$, and for nonsquarefree $m$ that line carries no nonreal pole.
\end{proposition}

\begin{proof}
For $\Re s>1$ the Euler products and M\"obius inversion give
\[
 F_f(s)=-\sum_{k\ge1}\mu(k)\sum_\chi c_\chi(f)\frac{L'}L(ks,\chi^k)
 =-\sum_{k\ge1}\mu(k)\sum_\psi c_\psi(U_kf)\frac{L'}L(ks,\psi),
\]
since the coefficient of $p^{-js}$, $p\nmid q$, is $f(p^j)\log p\sum_{k\mid j}\mu(k)$. On
$\Re s\ge\sigma_0>0$ the terms with $k\sigma_0\ge2$ converge normally, each logarithmic derivative
being at most $2^{-(k\sigma_0-2)}\sum_n\Lambda(n)n^{-2}$ and $\sum_\psi\abs{c_\psi(U_kf)}\le\sqrt{\abs G}\norm f_\infty$.
Omitted Euler factors have zeros on $\Re s=0$, and by \cref{src:thm:dgrh} a nonreal pole of the
$k$-term lies on $\Re s=1/(2k)$; a simple pole of $L'/L(ms,\psi)$ at $\rho/m$ has residue
$n_\psi(\rho)/m$.
\end{proof}

\paragraph{Triple-logarithmic windows.}
The multiplicity vectors have full rank in very short windows, without any hypothesis on the
zeros. Write $\mathcal Z_q$ for the set of distinct nontrivial zeros of the primitive inducers of the
characters modulo $q$. These inducers are distinct and form a family
closed under conjugation, so \cref{rec:thm:localrank} shows that the vectors $(n_\psi(\rho))_\psi$,
$\rho\in\mathcal Z_q$, $\abs{\Im\rho-T}\le\pi^{-1}\log\log\log\abs T+C_q$, span $\C^{\widehat G}$ for all large $\abs T$. The
local coefficient statistic of \cref{rec:thm:localcoeff} makes this quantitative.

\begin{corollary}[An approximate inverse in each window; \stU]\label{rec:cor:localmap}
Fix $q$ and $\delta>0$. There is $C_{q,\delta}$ such that for all sufficiently large $\abs T$ there is a linear
map $D_T$ from the data $\set{Z_g(\rho)}$, $\rho\in\mathcal Z_q$, $\abs{\Im\rho-T}\le\pi^{-1}\log\log\log\abs T+C_{q,\delta}$, to
$\C^G$ with $\norm{D_TZ_g-g}_\infty\le\delta\norm g_\infty$ for all $g$.
\end{corollary}

\begin{proof}
Choose a prime $p_a$ in each class $a$, fix $m$, and for $T>0$ let $h_a$ be the test of
\cref{rec:thm:localcoeff} for $p_a$, $m$ and $Q=q$, with window radius $H_a=16p_am+D(q,T)$. Put
$(D_TZ)(a)=-\frac{2\pi\sqrt{p_a}}{\log p_a}\sum_{\abs{\Im\rho-T}\le H_a}h_a(z_\rho)Z(\rho)$. For $g=\psi$ one has
$Z_\psi(\rho)=n_\psi(\rho)$, and $(D_TZ_\psi)(a)$ is the zero sum of $\mathcal E_{p_a}$ for the inducer of $\psi$; by the
theorem and the bound $O(\sqrt p/(T\log^2p))$ for the known integral (proof of \cref{rec:thm:localid}) it
differs from $\psi(a)$ by at most $C2^{-m}+o(1)$. By linearity and
$\sum_\psi\abs{c_\psi(g)}\le\sqrt{\abs G}\norm g_\infty$, choosing $m$ and then the height threshold gives the claim, with
$C_{q,\delta}=\max_a16p_am+O(1)$. For $T<0$ use $\overline{Z_g(\rho)}=Z_{\bar g}(\bar\rho)$.
\end{proof}

\begin{theorem}[Local recovery of the weights; \stG]\label{rec:thm:localrecovery}
Fix $b>1/\pi$. There are $C,T_0$ depending only on $q,b$ such that for every $m\ge1$ and
$\abs T\ge T_0$, the poles and residues of $F_f$ on $s=(\tfrac12+i\gamma)/m$,
$\abs{\gamma-T}\le b\log\log\log\abs T+C$, determine exactly $f|_{G^m}$ when $m$ is squarefree: two
weights give the same observations if and only if they agree on $G^m$, and the observation map has
rank $\abs{G^m}$. For nonsquarefree $m$ it has rank zero.
\end{theorem}

\begin{proof}
By \cref{rec:prop:powerline} the normalized residues are $Z_{U_mf}(\rho)$ on the window, which
misses the real axis for large $T_0$, and \cref{rec:thm:localrank} recovers $U_mf$, hence $f$ on
$G^m$. Conversely all residues factor through $f|_{G^m}$, and every function on $G^m$ extends to $G$.
\end{proof}

Exact recovery of complex weights is a finite calibration: by \cref{rec:cor:localmap} with
$\delta=\tfrac12$ the matrix with columns $D_TZ_{\mathbf 1_{\set{a'}}}$, computed from the known family, is within
$1/2$ of the identity, so its inverse has norm at most $2$.

\begin{corollary}[Lines with poles; natural boundary; \stG]\label{rec:cor:boundary}
The line $\Re s=1/(2m)$ contains a nonreal pole of $F_f$ exactly when $m$ is squarefree and
$\supp f\cap G^m\ne\varnothing$, and then it has poles in every sufficiently high window of the
theorem, in both directions. If $f\ne0$, the imaginary axis is a natural boundary of $F_f$.
\end{corollary}

\begin{proof}
Compare $f$ with $0$. For every prime $m$ not dividing the exponent of $G$ the power map is an
automorphism, so such lines carry poles; for real $t\ne0$ take $T=mt$ to get poles $s_m$ with
$\Re s_m=1/(2m)$ and $\abs{\Im s_m-t}\le(b\log\log\log\abs{mt}+C)/m$, accumulating at $it$ (for $t=0$ use one
large $T$).
\end{proof}

The windows in \cref{rec:thm:localrecovery,rec:cor:boundary} are those of the source. The spanning
statement \cref{rec:thm:localrank} used in their proofs now holds at the endpoint radius
$\pi^{-1}\log\log\log\abs T+C$.

Let $P$ be the set of prime divisors of the exponent of $G$, $S(m)=\set{\ell\in P:\ell\mid m}$ and
$B(y)=\set{\ell\in P:y\in G^\ell}$.

\begin{theorem}[Prescribing the lines; \stU]\label{rec:thm:design}
On squarefree indices, the patterns of lines with nonreal poles are exactly the downward closed
families of sets $S(m)$, and the least $\abs{\supp f}$ realizing a family is its number of maximal
members.
\end{theorem}

\begin{proof}
On the Sylow $\ell$-part an exponent prime to $\ell$ acts invertibly, so $G^m=\bigcap_{\ell\in S(m)}G^\ell$
for squarefree $m$. A supported point $y$ makes exactly the subsets of $B(y)$ visible. Every
$B\subseteq P$ is some $B(y)$: take the identity on the $\ell$-part for $\ell\in B$ and a non-$\ell$th
power for $\ell\notin B$. Each maximal member needs its own point.
\end{proof}

Among squarefree indices a family $\mathcal D$ has relative density
$\sum_{S\in\mathcal D}\prod_{\ell\in S}\frac1{\ell+1}\prod_{\ell\in P\setminus S}\frac\ell{\ell+1}$. Modulo
$31$, for $f=\delta_9-\delta_{27}$ (that is, $3^2$ and $3^3$) the lines carry poles exactly when
$5\nmid m$ and $6\nmid m$, a family with maximal sets $\set2,\set3$ and density $55/72$, so two
supported classes are needed; for $f=\delta_3$ the density is $5/12$. Modulo $4$, a simple zero of
the Dirichlet beta function not shared by $\zeta$ and a simple zeta zero not shared by it give, after
the sign change, $\tfrac12(f(1)-f(3))$ and $\tfrac12(f(1)+f(3))$; such zeros occur at heights
$12.98809801\ldots$ and $14.13472514\ldots$ (\artifact{prime_pole_windows/code/window_demo.txt}).

For a pair of characters, every window of radius $\pi^{-1}\log\log\log T+O(1)$ contains zeros with
$n_{\chi_1}>n_{\chi_2}$ and with $n_{\chi_2}>n_{\chi_1}$, for arbitrary conductors and every balanced real
direction of a larger family (\cref{rec:thm:bothsigns}).

\paragraph{Dyadic blocks with corrupted observations.}
Now let $\mathcal Z_T$ be the set of distinct $\gamma\in(T,2T)$ at which some $\nu_\psi(\gamma)=n_\psi(\tfrac12+i\gamma)$
is nonzero, $\nu(\gamma)=(\nu_\psi(\gamma))_\psi$, $d=\abs G$ and $L_T=T\log T$. Observations are
the normalized residues $X_{m,f}(\gamma)=\nu(\gamma)\cdot\widehat g$, $g=U_mf$, at positions matched to
the known zeros with their multiplicities.

\begin{theorem}[A positive proportion of distinct pole ordinates; \stG]\label{rec:thm:distinct}
There are $c_q>0$ and $T_q$ such that for $T\ge T_q$ and every nonzero $b\in\C^{\widehat G}$,
\[
 c_qT\log T\le\#\set{\gamma\in\mathcal Z_T:\textstyle\sum_\psi b_\psi\nu_\psi(\gamma)\ne0}\ll_qT\log T .
\]
\end{theorem}

\begin{proof}
The family of primitive inducers satisfies the hypotheses of the joint short-interval limit theorem
of Bombieri and Perelli \cite[Theorem~2]{BombieriPerelli1998Distinct}: Euler products of degree one,
functional equations, $\sum_{p\le x}\psi^*(p)\overline{\eta^*(p)}/p=\mathbf 1_{\psi=\eta}\log\log x+C+O_q(1/\log x)$,
and, by \cref{src:thm:dgrh}, no zeros to the right of $\sigma=1/2$. (For a single character adjoin a
second one and take a marginal.) With $t$ uniform on $[T,2T]$, $M\to\infty$,
$M\le\log T/\log\log T$, $h=M/\log T$ and $S_\psi(t)=\pi^{-1}\arg L(\tfrac12+it,\psi^*)$, the vector
$(\Delta_hS_\psi(t)/\sqrt{\log M})_\psi$ converges in law to a centred Gaussian vector $Y$ with
independent coordinates of positive variance, and $\Delta_hN_\psi(t)=M/(2\pi)+\Delta_hS_\psi(t)+O_q(1+h)$.
If the lower bound failed along $T\to\infty$, choose unit real vectors $d_T\to d$ (real or imaginary
parts of $b$) with supports of size $K=o(T\log T)$ and
$M=\min\set{\log T/\log\log T,\sqrt{T\log T/(1+K)}}$. Outside probability $(K+1)h/T=o(1)$ the window
$(t,t+h]$ avoids the support, so $\sum_\psi d_{T,\psi}\Delta_hN_\psi(t)=0$ and
$d_T\cdot Y_T=-M\sum_\psi d_{T,\psi}/(2\pi\sqrt{\log M})+o(1)$, a deterministic quantity up to $o(1)$,
contradicting the nondegenerate limit $d\cdot Y$. The upper bound is the zero count.
\end{proof}

\begin{theorem}[Frequent independent multiplicity vectors; \stG]\label{rec:thm:frequent}
There are $\delta_q>0$ and $B_q$ such that every sufficiently high block contains $d$ linearly
independent integer vectors $v_j$ with $n_T(v_j)=\#\set{\gamma\in\mathcal Z_T:\nu(\gamma)=v_j}\ge\delta_qL_T$
and $\norm{v_j}_1\le B_q$. They can be selected from the observed frequencies alone, by sorting the
types by decreasing frequency and keeping those that increase the rank.
\end{theorem}

\begin{proof}
By the zero count, $\sum_{\gamma\in\mathcal Z_T}\norm{\nu(\gamma)}_1\le C_qL_T$. With $c=c_q$ choose $B$ with
$C_q/B<c/4$; at most $cL_T/4$ ordinates have row sum above $B$, and discarding the at most
$(B+1)^d$ remaining types with fewer than $\delta_qL_T=cL_T/(4(B+1)^d)$ occurrences removes at most
$cL_T/4$ more. A vector annihilating the remaining rows would be supported on at most $cL_T/2$
ordinates, contradicting \cref{rec:thm:distinct}. The greedy selection picks rows of frequency at
least $\delta_qL_T$, whose sums are at most $B_q=\ceil{C_q/\delta_q}$.
\end{proof}

With $A=(v_j)$ and $D_q=d!B_q^{d-1}$, the cofactor formula and $\abs{\det A}\ge1$ give
$\norm{A^{-1}}_{\infty\to\infty}\le D_q$. Hence for every $b\in\C^d$ at least $\delta_qL_T$ distinct
ordinates with $\norm{\nu(\gamma)}_1\le B_q$ have $\abs{\nu(\gamma)\cdot b}\ge D_q^{-1}\norm b_\infty$, and
on a squarefree line with $f|_{G^m}\ne0$ the residues there satisfy
$\abs{\Res}\ge\norm{f|_{G^m}}_\infty/(mdD_q)$. After removing any $\delta_qL_T/2$ ordinates,
$c_q'L_T\norm b_2^2\le\sum_{\norm{\nu(\gamma)}_1\le B_q}\abs{\nu(\gamma)\cdot b}^2\le C_q'L_T\norm b_2^2$.

\begin{proposition}[Both signs of a balanced combination; \stG]\label{rec:prop:bothsigns}
For $d\ge2$ there are $a_q,\tau_q>0$ and $B_q'$ such that for large $T$ and every real $b$ with
$\sum_\psi b_\psi=0$, $\norm b_2=1$, at least $a_qL_T$ distinct ordinates in $\mathcal Z_T$ have
$\norm{\nu(\gamma)}_1\le B_q'$ and $\nu(\gamma)\cdot b\ge\tau_q$, and at least $a_qL_T$ have
$\nu(\gamma)\cdot b\le-\tau_q$.
\end{proposition}

\begin{proof}
If the positive support had size $o(L_T)$ along a sequence, the argument of \cref{rec:thm:distinct}
would give $b_T\cdot\Delta_hN\le0$ on almost all windows; as $\sum b_{T,\psi}=0$ the drift $M/(2\pi)$
cancels, so $b_T\cdot Y_T\le o(1)$ with probability tending to one, against a centred nondegenerate
Gaussian limit. Rows with large sums are excluded by the count, and a uniform margin follows by
compactness of the balanced sphere, since finitely many integer types of bounded sum occur.
\end{proof}

At $H=T$ this gives $\gg T\log T$ ordinates in each direction, more than the mass
$\gg T\log\log T$ of \cref{zer:thm:short}, which on the other hand applies in intervals of
triple-logarithmic length.

\begin{theorem}[Stable recovery; \stG]\label{rec:thm:stable}
There are $\alpha_q,K_q>0$ such that for large $T$, every squarefree $m$ and every $f$, if at most
$\alpha_q\#\mathcal Z_T$ of the normalized residues $X_{m,f}(\gamma)$, $\gamma\in\mathcal Z_T$, are erased or
arbitrarily corrupted and the others have error at most $\eps$, then $f|_{G^m}$ is recovered with
supremum error at most $K_q\eps$, exactly if $\eps=0$.
\end{theorem}

\begin{proof}
Select the rows of \cref{rec:thm:frequent} and put $\alpha_q=\delta_q/(4C_q)$; then at least three
quarters of the observations of each selected type are good. Coordinatewise medians of real and
imaginary parts give $y_j$ with $\abs{y_j-v_j\cdot b}\le\sqrt2\eps$, and $Ab^{\mathrm{rec}}=y$ gives
$\norm{b^{\mathrm{rec}}-b}_\infty\le\sqrt2D_q\eps$. Output $\sum_\psi b_\psi^{\mathrm{rec}}\psi(a_u)$ for
$a_u^m=u$, with $K_q=d\sqrt2D_q$.
\end{proof}

For $q=13$ the recoverable classes are all twelve at $m=1$; $1,3,4,9,10,12$ at $m=2$; $1,5,8,12$ at
$m=3$; none at $m=4$; and $1,12$ at $m=6$. A floating-point illustration modulo $4$ with two of five
pairs corrupted and noise $10^{-6}$ recovers $f(1)=2+3i$, $f(3)=-1+i$ to about $7\cdot10^{-7}$
(\artifact{dirichlet_zero_recovery/layers/observed_example.txt}).

\paragraph{Residues of single progressions.}
For $a\in G$ and $h=\abs G$ let
$D_a(s)=\sum_{n\equiv a}\Lambda(n)n^{-s}=-\frac1h\sum_\chi\overline{\chi(a)}\frac{L'}L(s,\chi)$. At any
$\rho$ in the strip, with $m_\chi=\ord_\rho L(s,\chi)$,
\begin{equation}\label{rec:eq:progres}
 -h\Res_{s=\rho}D_a(s)=F_m(a):=\sum_\chi m_\chi\overline{\chi(a)} .
\end{equation}
For a cyclic subgroup $H\le G$ put $d_H=\abs H$, $D_H=\varphi(d_H)$, let $\recGen H$ be its generators
(the identity for $H=1$), fix $a_H\in\recGen H$, and $K_d=\Q(e^{2\pi i/d})$.

\begin{lemma}[Rational blocks; \stU]\label{rec:lem:blocks}
The map $v\mapsto(F_v(a_H))_H$ is a rational isomorphism $\Q^{\widehat G}\to\prod_HK_{d_H}$, and the
values at the generators of $H$ are the conjugates of $F_v(a_H)$. Hence for $A\subseteq G$,
$\dim_\Q\ker(v\mapsto F_v|_A)=\sum_{H:A\cap\recGen H=\varnothing}D_H$, and $A$ recovers every
nonnegative integer vector from exact residues if and only if it meets every $\recGen H$.
\end{lemma}

\begin{proof}
Applying $e^{2\pi i/d_H}\mapsto e^{2\pi ik/d_H}$ to $F_v(a_H)$ gives $F_v(a_H^k)$; if all coordinates
vanish, $F_v=0$ on $G$ and $v=0$. The generator sets partition $G$, so both sides have dimension
$\abs G$. A nonzero rational kernel vector is, after scaling, a difference of two distinct
nonnegative integer vectors.
\end{proof}

For $M\ge1$ let $\mathcal M_M$ be the nonnegative integer vectors of total at most $M$ and
$\delta_A(M)=\min_{m\ne n\in\mathcal M_M}\max_{a\in A}\abs{F_m(a)-F_n(a)}$. For $A$ meeting every generator set,
let $b_H$ count the inverse pairs $\set{a,a^{-1}}$ met by $A$ in $\recGen H$ when $d_H>2$, and
$\kappa(A)=\max(\set0\cup\set{D_H/(2b_H)-1:d_H>2})$.

\begin{theorem}[Precision exponent; \stU]\label{rec:thm:precision}
For such $A$ and every $M\ge1$, $(2M)^{-\kappa(A)}\le\delta_A(M)\le C_{G,A}M^{-\kappa(A)}$. If $A$ misses a
generator set, $\delta_A(M)=0$ for large $M$. Hence the least $\abs A$ with $\delta_A(M)\gg M^{-\theta}$
is $\#\set{H:d_H\le2}+\sum_{d_H>2}\ceil{\varphi(d_H)/(2(1+\theta))}$.
\end{theorem}

\begin{proof}
For $v=m-n\ne0$ some block $\alpha=F_v(a_H)$ is a nonzero algebraic integer with all conjugates at
most $2M$. If $d_H\le2$ it is an observed nonzero integer. Otherwise at least $2b_H$ embeddings are
observed, of modulus at most $\eps=\max_A\abs{F_v}$, so
$1\le\abs{N(\alpha)}\le\eps^{2b_H}(2M)^{D_H-2b_H}$. For the upper bound take a block attaining
$\kappa(A)>0$, integer vectors $v_1,\dots,v_D$ spanning its inverse image, and the $(L+1)^D$ vectors
$\sum n_jv_j$, $0\le n_j\le L$, whose observed values lie in a box of side $O(L)$ in $\R^{2b_H}$;
pigeonhole with $J=\ceil{(L+1)^{D/(2b_H)}}-1$ cubes per side gives two images at distance
$O(L^{1-D/(2b_H)})$, and $L=\floor{M/B}$ concludes. The count follows by choosing $b_H$ in each block.
\end{proof}

Modulo $11$, $A=\set{1,10,4,2}$ meets every generator set, so four classes suffice for exact data;
$\delta_A(M)\asymp M^{-1}$, a fifth class does not change the exponent, and adding both $5$ and $8$
gives $\delta_A(M)=1$; uniform separation needs $(\abs G+\abs{G[2]})/2$ classes. For actual zeros of
height at most $T$, $\sum_\chi m_\chi\ll_q\log(T+3)$, so absolute error below
$\tfrac12(2M)^{-\kappa(A)}$ with $M\asymp_q\log(T+3)$ identifies the vector.

The detection question asks less: a single weighted count should see every zero off the line. For
$q\ge3$ choose a generator $u_H$ of each nontrivial cyclic subgroup, put $w(u_H)=\varphi(\abs H)$ and
$w=0$ elsewhere, so $\sum_aw(a)=n-1$ with $n=\abs G$, and let
$D_w(x)=\sum_{k\le x}\Lambda(k)w(k\bmod q)-\frac{n-1}n\sum_{k\le x,(k,q)=1}\Lambda(k)$. Let
$\recav:K\to\Q$ be the normalized trace of $K=\Q(\zeta_{\exp G})$ and $C_\chi(v)=\sum_av(a)\overline{\chi(a)}$.
Since $\varphi(d)\recav(\overline{\chi(u)})=\sum_{a\in\recGen\langle u\rangle}\overline{\chi(a)}$,
\begin{equation}\label{rec:eq:trace}
 \recav C_\chi(w)=\sum_{a\ne e}\overline{\chi(a)}=\begin{cases}n-1,&\chi=\triv,\\-1,&\chi\ne\triv.\end{cases}
\end{equation}

\begin{theorem}[A sparse weighted count that detects zeros; \stU]\label{rec:thm:sparsecount}
The following are equivalent: \textup{(i)} every nonprincipal $L(s,\chi)$ modulo $q$ is zero-free in
$\Re s>1/2$; \textup{(ii)} $D_w(x)=O_q(\sqrt x\log^2(2x))$; \textup{(iii)} $D_w(x)=O_{q,\eps}(x^{1/2+\eps})$
for every $\eps>0$. Assuming RH, the same holds with $\sum_{(k,q)=1}\Lambda(k)$ replaced by $x$.
\end{theorem}

\begin{proof}
$D_w=\frac1n\sum_{\chi\ne\triv}C_\chi(w)\psi(x,\chi)$, so (i) gives (ii) by the explicit formula
\cite{ErnvallHytonenPalojarvi2022Explicit}, imprimitivity costing $O_q(\log x)$. Under (iii),
$s\int_1^\infty D_w(x)x^{-s-1}\dd x=-\frac1n\sum_{\chi\ne\triv}C_\chi(w)\frac{L'}L(s,\chi)$ is holomorphic on
$\Re s>1/2$. At a zero $\rho$ there with multiplicities $m_\chi$, \eqref{rec:eq:trace} gives
$\recav\Res_{s=\rho}=\frac1n\sum_{\chi\ne\triv}m_\chi>0$, a contradiction.
\end{proof}

Modulo $7$ this is $\psi_6+2\psi_2+2\psi_3$, and for cyclic $G$ of order $n$ the support has
$\tau(n)-1$ classes. Say that a weight $v$ \emph{detects multiplicities} if
$\sum_{\chi\ne\triv}m_\chi C_\chi(v)\ne0$ for every nonzero tuple of nonnegative integers.

\begin{theorem}[Detecting weights; \stU]\label{rec:thm:detecting}
A rational weight detects multiplicities if and only if the rational numbers $\recav C_\chi(v)$,
$\chi\ne\triv$, are all positive or all negative. A set $A\subseteq G\setminus\set e$ supports a rational
detecting weight if and only if it meets $\recGen H$ for every nontrivial cyclic $H$, and if it
misses one, no complex weight supported on $A$ detects multiplicities; the least support is the
number of nontrivial cyclic subgroups.
\end{theorem}

\begin{proof}
A strict sign gives detection after taking traces. Conversely the trace is constant on Galois
orbits of characters and the orbit sums are rational: an orbit with zero trace, or two orbits of
opposite signs with suitable positive integer multiples, gives an undetected pattern. If $A$ misses
$B=\recGen H$, put $m_\chi=\abs B-\sum_{a\in B}\chi(a)$, a nonnegative rational integer with
$m_\triv=0$; orthogonality gives $\sum_\chi m_\chi\overline{\chi(g)}=n(\abs B\mathbf 1_{g=e}-\mathbf 1_{g\in B})$,
which vanishes on $A$, so no weight on $A$ detects it. The generator sets are disjoint.
\end{proof}

The criterion for exact recovery of vectors (\cref{rec:lem:blocks}) includes the identity class,
which detection excludes; allowing $e$, the singleton weight $\mathbf 1_{\set e}$ detects.

\subsection{Quadratic characters from finitely many zeros}\label{rec:sec:quadratic}
\sources{quadratic_spectra}

Let $\chi$ be a primitive nonprincipal quadratic character of conductor $q$ and parity $\kappa$, and
$p_-(\chi)$ the least prime with $\chi(p)=-1$ (the least inert prime for the character of a
quadratic field). Under \cref{src:thm:dgrh} write the zeros as $\tfrac12+i\gamma$ with both signs. For
compact intervals $I,J$ of positive length put $v_I(\gamma)=\int_Ie^{i\gamma x}\dd x$ and
$C_\chi(I,J)=\Re\sum_\gamma\overline{v_I(\gamma)}v_J(\gamma)$, an absolutely convergent pairing since
$\abs{v_I(\gamma)}\le\min(\abs I,2/\abs\gamma)$. With $g_{I,J}(u)=\abs{I\cap(J-u)}$ and
$K_\kappa(u)=e^{-(\kappa+1/2)u}/(1-e^{-2u})$, the Guinand--Weil formula
\cite{ChirrePereiraDeLaat2021Primes} gives, for $J-I\subset(0,\infty)$,
\begin{equation}\label{rec:eq:packetweil}
 C_\chi(I,J)=-\int_0^\infty g_{I,J}(u)K_\kappa(u)\dd u-\sum_{n\ge2}\frac{\Lambda(n)\chi(n)}{\sqrt n}g_{I,J}(\log n).
\end{equation}
Indeed, for smooth $g$ supported in $(0,\infty)$ and $H(t)=\int g(u)\cos(tu)\dd u$ the conductor term
vanishes, the prime terms are as displayed, and the gamma term, from the integral representation of
$\psi$, is $-\int gK_\kappa$; the trapezoid $g_{I,J}$, whose cosine transform is
$\Re(\overline{v_I(t)}v_J(t))$, is reached by mollification.

\begin{theorem}[Sign threshold; \stG]\label{rec:thm:signthreshold}
The supremum $R_\chi$ of the $R>0$ such that $C_\chi(I,J)<0$ whenever $J-I\subset(0,R)$ equals
$\log p_-(\chi)$. For every $R>R_\chi$ there are equal-length intervals, of arbitrarily small length,
with $J-I\subset(0,R)$ and $C_\chi(I,J)>0$.
\end{theorem}

\begin{proof}
If $J-I\subset(0,\log p_-)$, every prime power in \eqref{rec:eq:packetweil} has $\chi\ge0$, and
$K_\kappa>0$, so the pairing is negative. For $r=\log p_-$, $I_h=[0,h]$ and $J=[r,r+h]$, the weight
$(h-\abs{u-r})_+$ sees for small $h$ only the prime $p_-$, contributing $h\log p_-/\sqrt{p_-}$, while
the gamma term is at most $h^2K_\kappa(r-h)$.
\end{proof}

With $n=p^k$, $0<h<\log(1+1/n)$ and $A_\kappa(r,h)=\int(h-\abs{u-r})_+K_\kappa(u)\dd u$, the same formula
gives the exact identity
$\chi(n)=-\frac{\sqrt n}{h\log p}\bigl(C_\chi(\log n,h)+A_\kappa(\log n,h)\bigr)$, where
$C_\chi(r,h)=\sum_\gamma\frac{4\sin^2(\gamma h/2)}{\gamma^2}\cos(\gamma r)$, and by the explicit zero count
of \cite{BennettMartinOBryantRechnitzer2021Counting} the part with $\abs\gamma>T$ is at most
$16(\log(qT)+1)/T$, or
$\frac{4(\log(qT/(2\pi))+1)}{\pi T}+\frac{4(2b(T)+0.149)}{T^2}$ with $b(T)=0.298\log(qT)+4.358$. For
the character modulo $3$ with $r=\log2$, $h=3/8$ and the $46$ positive ordinates up to $100$, the
estimate of $\chi(2)$ is $-0.98283822$, which rounds to $-1$; the cutoff bound is proved and the
evaluations are numerical (\artifact{quadratic_spectra/spectrum_example.txt}). A higher convolution
power recovers all small prime powers from one list.

For $Q\ge q$ and $X\ge2$ let $h=1/(8X)$, choose an even $r\ge2$ with $T=2r/h=16rX$ such that
\begin{equation}\label{rec:eq:choose}
 2^r\ge\frac{32}{\log2}\,r\sqrt X\log(eQT),
\end{equation}
and let $f$ be the $r$-fold convolution of the uniform density on $[-h/r,h/r]$, so
$\widehat f(t)=\recsinc(ht/r)^r$.

\begin{theorem}[Finite recovery with a rounding margin; \stG]\label{rec:thm:quaddecode}
Let $\Gamma$ be the multiset of ordinates $\abs\gamma\le T$ of $L(s,\chi)$, $\chi$ primitive
nonprincipal quadratic of conductor at most $Q$. For each $n=p^a\le X$ the nearest integer to
\[
 B(n;\Gamma)=-\frac{\sqrt n}{(\log p)f(0)}\sum_{\gamma\in\Gamma}\cos(\gamma\log n)\recsinc(h\gamma/r)^r
\]
is $\chi(n)$, including $0$ at ramified primes, simultaneously for all such $n$, and still with
approximate ordinates $\widetilde\gamma$ whenever
$\sum_\Gamma\abs{\widetilde\gamma-\gamma}\le\log2/(16\sqrt Xh(\log X+h))$. One can take
$T\ll X(\log(2X)+\log\log(3Q))$ with an absolute constant.
\end{theorem}

\begin{proof}
$f=v*v$ with $v$ an even density on an interval of length $h$, so $f(0)=\int v^2\ge1/h$. For $n\le X$
the neighbouring integer logarithms are more than $h$ away, so $g(u)=f(u-\log n)$ in the explicit
formula sees only $n$, and the infinite sum is $-\chi(n)\log p\,f(0)/\sqrt n-\int f(u-\log n)K_\kappa(u)\dd u$,
with $K_\kappa<2/\sqrt n$ on the support; normalized, the integral is at most $2h/\log2$. From
$N(t)\le2t\log(Qt)$ for $t\ge2$,
$\sum_{\abs\gamma>T}\abs{\widehat f(\gamma)}\le(r/h)^rr\int_T^\infty N(t)t^{-r-1}\dd t\le4T\log(eQT)2^{-r}$,
which after normalization is at most $1/4$ by \eqref{rec:eq:choose}. The test has Lipschitz constant
at most $\log X+h$ in $\gamma$, so the perturbation costs $1/16$, and
$1/4+1/(8\log2)+1/16<1/2$. Finally $r\asymp\log(2X)+\log\log(3Q)$ satisfies \eqref{rec:eq:choose}.
\end{proof}

With $X=C\log^2(2Q)$, admissible by the bound $O(\log^2(2M))$ of Lamzouri, Li and Soundararajan
\cite[Theorem~1.1]{LamzouriLiSoundararajan2015Conditional} for the least prime outside the kernel of a
nontrivial quadratic character modulo $M\le Q^2$, the theorem determines every primitive quadratic
character of conductor at most $Q$ from its zeros with $\abs\gamma\ll\log^2(2Q)\log\log(3Q)$;
\cref{rec:thm:lowzero} lowers this height to $O(\log^2(2Q)/\log\log(3Q))$ for characters of every order.

\begin{theorem}[Separation at every higher cutoff; \stG]\label{rec:thm:quaddisplace}
With $X,r,h,T$ as above and $H\ge2T$, suppose two primitive nonprincipal quadratic characters of
conductor at most $Q$ differ at a prime $p\le X$ and have the same number of zeros with
$\abs\gamma\le H$. Then every bijection between the two lists satisfies
$\sum\abs{\gamma-\gamma'}>H/(144\pi\sqrt X)$; for distinct characters, with $X=C\log^2(2Q)$, this is $\gg H/\log(2Q)$.
\end{theorem}

\begin{proof}
Use $\phi(u)=f(u)\recsinc^2(Bu)/f(0)$ with $B=H/8$: $\phi(0)=1$, $\int\phi\le\min(h,\pi/B)$, and
$\widehat\phi$ averages $\widehat f(t+V)/f(0)$ over $\abs V\le2B$, with $\abs{t+V}\ge3\abs t/4$ for
$\abs t>H$. The gamma term is again at most $2h/\log2$, the normalized tail is at most
$\frac{4h\sqrt X}{\log2}H\log(eQH)(4r/(3hH))^r\le1/4$ for $H\ge2T$ (since $r\ge4$), and the
normalized test has Lipschitz constant at most $9\pi\sqrt X/H$. A matching of smaller displacement
would force equal values at $p$ by rounding.
\end{proof}

For $q=Q=3$ and $X=5$ the rule gives $r=14$, $h=1/40$, $T=1120$; with $941$ positive ordinates the
finite sums at $n=2,3,4,5$ are $-0.990596$, $0.003337$, $1.003759$, $-0.998733$, rounding to
$-1,0,1,-1$ (\artifact{quadratic_spectra/decoder_example.txt}; the root enumeration is numerical).

\subsection{Characters, products and abelian fields from low zeros}\label{rec:sec:lowzero}
\sources{dirichlet_low_zero_identification; dirichlet_zeros_and_abelian_fields; robust_dirichlet_spectra_2026-11-26; quadratic_spectra}

Let $\mathcal C(Q)$ be the set of primitive characters of conductor at most $Q\ge3$, the
conductor-one character included, $L_Q=\log(2Q)$, and $m_\chi(t)=\ord_{s=1/2+it}L(s,\chi)$. For
$\tau\in\R$, $H>0$ and $F\ge1$ put
\[
 I=I(\tau,H)=[\tau-H,\tau+H],\qquad J(F,\tau)=\log\log\bigl(e^eF(3+\abs\tau)\bigr).
\]
Sums over $t\in I$ run over the distinct ordinates in the union of the spectra in question. The
height needed to identify a primitive character of conductor at most $Q$, of any order, is
$O(L_Q^2/\log\log(3Q))$: this replaces the quadratic cutoff $O(L_Q^2\log\log(3Q))$ that follows from
\cref{rec:thm:quaddecode}. In this subsection $\widehat g(t)=\int g(u)e^{itu}\dd u$. For $g$ supported in a compact
subinterval of $(0,\infty)$, with $a$ the parity and $\delta_\chi=\mathbf 1_{q=1}$, the explicit formula is
\begin{equation}\label{rec:eq:halfweil}
 \sum_\gamma m_\chi(\gamma)\widehat g(\gamma)=-\sum_{n\ge2}\frac{\Lambda(n)\chi(n)}{\sqrt n}g(\log n)
 -\int_0^\infty g(u)\frac{e^{-(a+1/2)u}}{1-e^{-2u}}\dd u+\delta_\chi\int_0^\infty g(u)\bigl(e^{u/2}+e^{-u/2}\bigr)\dd u,
\end{equation}
the positive-support case of the Dirichlet explicit formula \cite{ChirrePereiraDeLaat2021Primes}, of which
\eqref{rec:eq:packetweil} is the quadratic instance: the conductor term and the reflected prime term
vanish because $g(0)=g(-\log n)=0$, and compact $C^2$ tests follow by smoothing. Every unit interval
$[v,v+1]$ contains $O(\log(3q(2+\abs v)))$ ordinates with multiplicity
\cite{BennettMartinOBryantRechnitzer2021Counting}.

\paragraph{Neighbouring prime powers.}
For a prime $p$ and $B>0$ let
\[
 W_p(B)=\sum_{\substack{n\ne p\\ \abs{\log(n/p)}\le1/2}}\frac{\Lambda(n)}{\log p}\sqrt{\frac pn}\,\min\bigl\{1,(B\log(n/p))^{-2}\bigr\},
\]
and $\mathcal P_X$ the set of primes in $[X,2X]$, $P_X=X/\log X$. A single $W_p$ can be large; their mean is
small at bandwidth $B\gg X/\log X$.

\begin{lemma}[Mean interference from neighbours; \stU]\label{rec:lem:neighbours}
Uniformly for all sufficiently large $X$ and all $B>0$,
$\frac1{P_X}\sum_{p\in\mathcal P_X}W_p(B)\ll\frac X{B\log X}$.
\end{lemma}

\begin{proof}
On the support $Xe^{-1/2}\le n\le2Xe^{1/2}$, the prefactor is bounded, and the minimum is at most
$b(h)=\min(1,(CX/(Bh))^2)$ with $h=\abs{n-p}\ge1$. For $h$ even, sieving the classes $0$ and $\mp h$
modulo primes $\ell\le y=X^{1/16}$ (one class when $\ell\mid h$) and applying the arithmetic large sieve
\cite{Montgomery1973Large} with $R=\sqrt X$ gives
\[
 \#\set{p\in\mathcal P_X:\ p\pm h\text{ prime in }[Xe^{-1/2},2Xe^{1/2}]}\ll\frac X{\log^2X}\,s(h),\qquad
 s(h)=\prod_{\substack{\ell\mid h\\\ell>2}}\frac{\ell-1}{\ell-2};
\]
indeed, by Markov's inequality at least half of the Euler product of the sieve weights is carried by
$d\le R$, and Mertens' product formula bounds its reciprocal by $s(h)/\log^2y$. Odd $h$ give no prime
pairs. Since $\sum_{h\le z}s(h)\le z\prod_{\ell>2}(1+1/(\ell(\ell-2)))\ll z$, dyadic summation gives
$\sum_hb(h)s(h)\ll X/B$, so prime neighbours contribute $O(X^2/(B\log^2X))$ to the sum over centres.
The $O(\sqrt X)$ proper prime powers in the range contribute, each summed over all integer centres,
$O(\sum_hb(h))=O(X/B)$, which after division by $P_X$ is absorbed.
\end{proof}

\paragraph{Tests on translated windows.}
For $H\ge800$ take $B=H/4$, $m=\floor{H/200}\ge4$, $w_m(u)=(1-4u^2)_+^m$ and
\[
 g_p(u)=e^{-i\tau(u-\log p)}w_m(u-\log p)\recsinc^2\bigl(B(u-\log p)\bigr),\qquad h_p=\widehat g_p,
\]
so that $g_p(\log p)=1$. Since $\int\recsinc^2(Bu)\dd u=\pi/B$,
\begin{equation}\label{rec:eq:lowlip}
 \abs{h_p(t)}\le\pi/B,\qquad\abs{h_p'(t)}\le(\log p+\tfrac12)\pi/B\qquad(t\in\R).
\end{equation}
Integrating $m$ times by parts and using Cauchy's estimate on circles of radius one for $(1-4z^2)^m$
gives $\abs{\widehat w_m(t)}\le(10m/\abs t)^m$; the factor $\recsinc^2(Bu)$ is an average of $e^{ivu}$ over
$\abs v\le2B$, so $\abs{h_p(t)}\le(20m/\abs{t-\tau})^m$ for $\abs{t-\tau}\ge H$, and dyadic summation with the
unit-interval count gives
\begin{equation}\label{rec:eq:lowtail}
 \sum_{\abs{\gamma-\tau}>H}\abs{h_p(\gamma)}\ll He^{-c_0H}\log\bigl(2q(\abs\tau+H+3)\bigr).
\end{equation}
For primitive $\chi_1,\dots,\chi_r$ of conductors at most $Q$ and $v\in\C^r$ let
$A(n)=\sum_jv_j\chi_j(n)$, $M(t)=\sum_jv_jm_{\chi_j}(t)$ and
$Z_v(p)=-\frac{\sqrt p}{\log p}\sum_{t\in I}h_p(t)M(t)$. Then \eqref{rec:eq:halfweil} gives the identity
\begin{equation}\label{rec:eq:lowobs}
 Z_v(p)=A(p)+\sum_{n\ne p}\frac{\sqrt p}{\log p}\frac{\Lambda(n)}{\sqrt n}g_p(\log n)A(n)+E_v(p),
\end{equation}
\[
 \abs{E_v(p)}\ll\frac{\norm v_1}{B\log X}+\abs{v_{\mathrm{triv}}}\frac X{B\log X}
 +\norm v_1\frac{\sqrt X}{\log X}He^{-c_0H}\log\bigl(2Q(\abs\tau+H+3)\bigr)\qquad(p\in\mathcal P_X),
\]
with $v_{\mathrm{triv}}$ the coefficient of $\zeta$ (zero if absent): on the support of $g_p$ the gamma
kernel is $O(X^{-1/2})$ and the pole kernel $O(\sqrt X)$, and $\int\abs{g_p}\le\pi/B$. Forming $Z_v(p)$ needs no
knowledge of conductors or parities.

\begin{theorem}[Linear independence in a finite window; \stG]\label{rec:thm:family}
Let $\chi_1,\dots,\chi_r$ be distinct primitive characters of conductors at most $Q$, and
$X=C_0r^2L_Q^2$ with $C_0$ a large absolute constant. For every $v\in\C^r$ and $\tau\in\R$,
\[
 H\ge C_1\bigl\{rX/\log X+J(Q,\tau)\bigr\}\ \Longrightarrow\
 \sum_{t\in I}\Bigl\lvert\sum_{j=1}^rv_jm_{\chi_j}(t)\Bigr\rvert\ \ge\ cH\frac{\log X}{\sqrt{rX}}\norm v_2 .
\]
Consequently $(m_{\chi_j}(t))_{t\in I,\,j\le r}$ has column rank $r$, and some $r$ ordinates in $I$ give an
invertible minor.
\end{theorem}

\begin{proof}
By \cref{rec:lem:gram} for a fixed smooth $0\le w\le1$ supported in $(1,2)$, with total error
$O(\sqrt XL_Q\norm v_1^2)$ and $\abs{A(p)}\le\sqrt r\norm v_2$,
\begin{equation}\label{rec:eq:lowrms}
 \sum_{p\in\mathcal P_X}\abs{A(p)}^2\gg P_X\norm v_2^2,\qquad\frac1{P_X}\sum_{p\in\mathcal P_X}\abs{A(p)}\gg\frac{\norm v_2}{\sqrt r}.
\end{equation}
Since $\abs{g_p(\log n)}\le\min\set{1,(B\log(n/p))^{-2}}$, the off-diagonal term of \eqref{rec:eq:lowobs} has mean at
most $\norm v_1P_X^{-1}\sum_pW_p(B)\ll\sqrt r\norm v_2X/(B\log X)$ by \cref{rec:lem:neighbours}, below a quarter of
\eqref{rec:eq:lowrms} for $B\ge CrX/\log X$. The first two error terms are also small, and
$H\gg rX/\log X$ together with $H\ge CJ(Q,\tau)$ absorbs the third, as
$\log(2Q(\abs\tau+H+3))\ll\exp(J(Q,\tau))+\log(H+3)$. Hence some $p$ has $\abs{Z_v(p)}\gg\norm v_2/\sqrt r$, while
\eqref{rec:eq:lowlip} gives $\abs{Z_v(p)}\ll\frac{\sqrt X}{H\log X}\sum_{t\in I}\abs{M(t)}$.
\end{proof}

\begin{corollary}[Stability under partial matching; \stG]\label{rec:cor:lowmatch}
Under the hypotheses of \cref{rec:thm:family} with $r=2$, match submultisets of the two zero lists in $I$
bijectively, let $E$ be the number of unmatched occurrences and $D$ the sum of the matched
displacements. Then $E+(\log(2X)+1)D\gg H\log X/\sqrt X$; in particular every complete matching has
$D\gg H/L_Q$. The same holds for half-open $I$.
\end{corollary}

\begin{proof}
For $v=(1,-1)$ some prime has $\abs{Z_v(p)}\gg1$; bound unmatched entries by the first and matched pairs
by the second inequality of \eqref{rec:eq:lowlip}.
\end{proof}

The same construction reads every prime value at once on a common set of primes.

\begin{proposition}[A common set of readable primes; \stG]\label{rec:prop:read}
Let $X=AL_Q^2$ with $A$ a large absolute constant, and $H\ge B_1\{X/\log X+J(Q,\tau)\}$ with $B_1$ large.
There is a set $\mathcal P\subseteq\mathcal P_X$, independent of the character, such that with
$Z(p;\chi)=-\frac{\sqrt p}{\log p}\sum_{t\in I}h_p(t)m_\chi(t)$,
\[
 \abs{Z(p;\chi)-\chi(p)}\le\tfrac18\qquad(\chi\in\mathcal C(Q),\ p\in\mathcal P),
\]
also with either endpoint of $I$ omitted, and every pair of distinct $\chi,\psi\in\mathcal C(Q)$ has at least
$bX/\log X$ primes $p\in\mathcal P$ with $\abs{\chi(p)-\psi(p)}\ge1$, and, if $\psi\notin\set{\chi,\bar\chi}$, at least
$bX/\log X$ with $\abs{\Re\chi(p)-\Re\psi(p)}\ge\tfrac12$, for an absolute $b>0$.
\end{proposition}

\begin{proof}
On $(\Z/M\Z)^\times$, $M=\lcm(q_\chi,q_\psi)\le Q^2$, the characters remain distinct; with weights
$w_p=(\log p)W(p/X)$, $W$ a fixed smooth weight with $0\le W\le1$ supported in $(1,2)$ and $\int W>0$, \cref{rec:lem:gram} applied to the monomials of $\abs{\chi-\psi}^2$ (mean $2$ on the group)
and of $(\Re\chi-\Re\psi)^2$ (mean at least $1$ for nonconjugate pairs, the real parts being orthogonal),
each of total absolute coefficient at most $4$, shows that the prime-weighted means differ from the
group means by at most $\tfrac14$ for $A$ large, uniformly in $Q$; as the squares are at most $4$ and
$w_p\le\log(2X)$, at least $2bX/\log X$ primes witness each pair. By \eqref{rec:eq:lowobs} with $v=e_\chi$,
$\abs{Z(p;\chi)-\chi(p)}\le W_p(B)+\abs{E(p)}$, and the error is at most $1/16$ for $B_1$ large. By
\cref{rec:lem:neighbours} and Markov's inequality, all but $bX/\log X$ primes of $\mathcal P_X$ have
$W_p(H/4)\le1/16$; let $\mathcal P$ be these. The tail bound covers either endpoint.
\end{proof}

\begin{theorem}[Identification from low zeros; \stG]\label{rec:thm:lowzero}
There are absolute $C,c>0$ such that, with $H_Q=L_Q^2/\log\log(3Q)$:
\begin{enumerate}
\item the multiset of ordinates $0<\gamma\le CH_Q$, with multiplicity, determines $\chi\in\mathcal C(Q)$;
\item for distinct members, even after matching occurrences in these lists at distance at most
$c/L_Q^2$, at least $cL_Q$ occurrences remain unmatched;
\item the multiset of $\abs\gamma$ over all signed ordinates with $\abs\gamma\le CH_Q$ determines the pair
$\set{\chi,\bar\chi}$, and conjugate characters have the same folded data at every height.
\end{enumerate}
\end{theorem}

\begin{proof}
Take $H\asymp H_Q$ with a large lower constant, $\tau=H$, and $I$ with its left endpoint omitted, which
is $(0,2H]$; then $H\gg X/\log X+J(Q,\tau)$ for $X=AL_Q^2$. By \cref{rec:cor:lowmatch},
$E\gg H\log X/\sqrt X\gg L_Q$ once $D$ is absorbed: each list has $O(HL_Q)$ occurrences, so displacements
at most $c/L_Q^2$ give $D\ll cH/L_Q$, and $(\log X)D$ is a small multiple of $H\log X/\sqrt X$ for small $c$,
as $\sqrt X=\sqrt AL_Q$. This proves (i) and (ii). For (iii) take $\tau=0$: the test is then real and even
about $\log p$, so $\Re h_p(\gamma)=\cos(\gamma\log p)\widehat{g_0}(\gamma)$ is even in $\gamma$ and $\Re Z(p;\chi)$ is computed
from the folded list. For nonconjugate pairs a prime of \cref{rec:prop:read} has
$\abs{\Re\chi(p)-\Re\psi(p)}\ge\tfrac12>2\cdot\tfrac18$. Finally $L(\bar s,\bar\chi)=\overline{L(s,\chi)}$ reflects the signed list.
\end{proof}

A list of positive ordinates keeps only the zeros above the real axis, while a folded list
superimposes both half-planes; for a complex character they carry different information. The
discrepancy in (ii) may occur at an ordinate shared with different multiplicities.

\begin{corollary}[A dictionary of $O(\log Q)$ primes; \stG]\label{rec:cor:dict}
There are $k=O(L_Q)$ primes in $\mathcal P$, chosen independently of the unknown character, and an
absolute $d>0$, such that every distinct pair in $\mathcal C(Q)$ has $\abs{\chi(p)-\psi(p)}\ge1$ on at least $dk$ of
them, and every nonconjugate pair has $\abs{\Re\chi(p)-\Re\psi(p)}\ge\tfrac12$ on at least $dk$. A decoder
choosing a candidate within $\tfrac18$ on the most coordinates succeeds with arbitrary corruption of
fewer than $dk/2$ readings if the others have error at most $\tfrac18$. For dictionaries of exact prime
evaluations the order of $k$ cannot be reduced, even for quadratic characters.
\end{corollary}

\begin{proof}
A fraction $\beta>0$ of $\mathcal P$ witnesses each pair. Choose $k=\ceil{C_2L_Q}$ primes uniformly without
replacement; a hypergeometric tail bound gives probability at most $e^{-\beta k/8}$ of fewer than $\beta k/2$
witnesses for a given pair, and a union bound over at most $2Q^4$ pairs (there are at most $Q^2$
characters) gives existence. On a witnessing coordinate the two radius-$\tfrac18$ disks are disjoint.
For the lower bound, the $\gg Q/\log Q$ prime conductors in $(Q/2,Q]$ give distinct quadratic characters
and each prime evaluation takes at most three values, so $3^k\gg Q/\log Q$.
\end{proof}

\begin{example}\label{rec:ex:dict101}
For $q=101$, $X=10^4$, $R=8\cdot10^4$ and the test half-width $\tfrac18$, the primes $11197$, $14107$, $11027$
(indices $60,32,39$ to the base $2$) separate all $100$ characters modulo $101$ by chords of length at
least one, with certified leakage bounds below $0.05$ obtained by enumerating every prime power in an
enlarged window. For the quartic character $\chi$ modulo $5$ with $\chi(2)=i$, interval arithmetic with
Euler--Maclaurin remainders and Rouch\'e's theorem shows unconditionally that the disk of radius
$10^{-4}$ about $\tfrac12+4.13290i$ contains exactly one simple zero of $L(s,\bar\chi)$, which lies on the
critical line by reflection, and no zero of $L(s,\chi)$:
\begin{center}
\artifact{dirichlet_low_zero_identification/checks/prime_dictionary.txt}\
\artifact{dirichlet_low_zero_identification/checks/quartic_certificate.txt}
\end{center}
\end{example}

\begin{corollary}[Products of $L$-functions; \stG]\label{rec:cor:products}
Fix $k\ge1$. A product $\prod_jL(s,\chi_j)^{e_j}$ with at most $k$ distinct factors $\chi_j\in\mathcal C(Q)$ and
arbitrary positive integer exponents is determined, with its factors and exponents, by the pooled
multiset of its positive zero ordinates up to $O_k(L_Q^2/\log\log(3Q))$, without labels.
\end{corollary}

\begin{proof}
Orders add under multiplication at every point of the open strip. Two different products give a
nonzero integer vector $v$ of exponent differences on at most $2k$ distinct characters, the empty
product included; apply \cref{rec:thm:family} with $r\le2k$ on a positive window as in the proof of
\cref{rec:thm:lowzero}.
\end{proof}

\paragraph{Abelian fields.}
For a finite abelian $K/\Q$ let $\mathcal X_K$ be its group of Dirichlet characters, each represented by
its primitive character. Then $\zeta_K(s)=\prod_{\chi\in\mathcal X_K}L(s,\chi)$ and
$\abs{\operatorname{disc}K}=\prod_{\chi\in\mathcal X_K}q_\chi$, and inside a common cyclotomic field $\mathcal X_K$ determines $K$, its
annihilator being the subgroup fixing $K$ \cite{Washington1997Introduction}. Write $m_K(t)$ for the
multiplicity of $\zeta_K$ at $\tfrac12+it$.

\begin{theorem}[Abelian fields of unrestricted degree; \stG]\label{rec:thm:fields}
There are absolute $C,c>0$ such that for distinct finite abelian $K,L\subset\overline\Q$ with absolute
discriminants at most $D\ge3$, $r$ the size of the symmetric difference of $\mathcal X_K$ and $\mathcal X_L$, every
$\tau\in\R$ and $H\ge C\{\log^2(2D)+J(D,\tau)\}$,
\[
 \sum_{t\in I}\abs{m_K(t)-m_L(t)}\ge cH\sqrt r\,\frac{\log\log(3D)}{\log(2D)} .
\]
In particular the zeros with $\abs t\le C'\log^2(2D)$, or the positive ones up to that height,
determine an abelian field of absolute discriminant at most $D$, independently of its degree.
\end{theorem}

\begin{proof}
After cancelling common factors, $M(t)=m_K(t)-m_L(t)=\sum_{j\le r}v_jm_{\chi_j}(t)$ with $v_j=\pm1$, the
trivial character cancelling, $S_0=\sum_j\log(2q_j)\ll\log(2D)$ and $r\ll\log(2D)$, since each $q_j\ge3$ and
$\sum\log q_j\le2\log D$. Take $X=C_2\log^2(2D)$. Summing the errors of \cref{rec:lem:gram} over all pairs
gives $O(\sqrt XrS_0)$, below half the diagonal $rX\int w$; with a second weight majorizing the
indicator of $\mathcal N_X=[e^{-1/2}X,2e^{1/2}X]$,
$\sum_{p\in\mathcal P_X}\abs{A(p)}^2\gg rP_X$ and $\sum_{p\in\mathcal N_X}\abs{A(p)}^2\ll rP_X$. The vector
$b_n=(\Lambda(n)/\log X)\sqrt{X/n}\,A(n)$ on $\mathcal N_X$ has $\norm b_2^2\ll rP_X+r^2\sqrt X/\log X\ll rP_X$, using
$\abs{A(n)}\le r$ at proper prime powers. The off-diagonal matrix
$T_{pn}=\mathbf 1_{p\ne n}\sqrt{p/X}(\log X/\log p)g_p(\log n)$ has $\abs{T_{pn}}\ll(X/B)^2\abs{n-p}^{-2}$, so its row and
column sums are $O((X/B)^2)$ and Schur's test gives $\norm{Tb}_2\ll(X/B)^2\norm b_2$, a quarter of the
central norm for $B\ge K_0X$. With $v_{\mathrm{triv}}=0$, $\norm v_1=r\ll\sqrt X$ and $H\ge C(X+J(D,\tau))$ the
errors are another quarter, so $\sum_{p\in\mathcal P_X}\abs{Z_v(p)}^2\gg rP_X$, some $\abs{Z_v(p)}\gg\sqrt r$, and
\eqref{rec:eq:lowlip} concludes. Positive ordinates suffice by the window $[H,3H]$.
\end{proof}

\begin{example}[Two biquadratic fields]\label{rec:ex:biquad}
$K=\Q(\sqrt5,\sqrt{221})$ and $L=\Q(\sqrt{13},\sqrt{85})$ are totally real with discriminant $1105^2$
and share $\zeta$ and $L(s,\chi_{1105})$, while
$\chi_5(3)+\chi_{221}(3)-\chi_{13}(3)-\chi_{85}(3)=-4$. Unconditionally, $\zeta_K(\tfrac12+it)$ vanishes for some
$1.20<t<1.23$, while $\zeta_L(\tfrac12+it)\ne0$ on $[1.20,1.23]$: outward interval enclosures of the
theta integral $F_d(t)=4\int_0^\infty e^{u/2}\cos(tu)\sum_n\chi_d(n)e^{-\pi n^2e^{2u}/d}\dd u$, which differs from
$L(\tfrac12+it,\chi_d)$ by nonzero factors, give a sign change for $d=221$ and positivity for
$d=13,85,1105$, and Euler--Maclaurin gives $\Im\zeta(\tfrac12+it)<-0.41$ there. Under \cref{src:thm:dgrh},
$\sum_{0<t\le80}\abs{m_K(t)-m_L(t)}\ge7$, from the test $g(u)=\phi((u-\log3)/a)$, $a=2/7$,
$\phi(x)=(x(1-x^2)_+^7)'$, whose support meets only $\log3$, with $\abs{\widehat g(t)}\le C_6\abs t^{-6}$,
$C_6\le2629707255/16$, and the explicit zero count of
\cite[Corollary~1.2]{BennettMartinOBryantRechnitzer2021Counting} for the tail:
\begin{center}
\artifact{dirichlet_zeros_and_abelian_fields/checks/theta_example.txt}\\
\artifact{dirichlet_zeros_and_abelian_fields/checks/field_example.txt}
\end{center}
\end{example}

\paragraph{Explicit constants and a clipped matching metric.}
For finite real multisets $A,B$ let $d_{\mathrm{edit}}(A,B)$ be the least cost of a partial matching, an
unmatched occurrence costing $1$ and a matched pair $\min\set{1,\abs{a-b}}$, and let $\Gamma_\chi(H)$ be the
signed ordinates of $L(s,\chi)$ with $\abs\gamma\le H$, with multiplicity.

\begin{lemma}[The comparison metric; \stU]\label{rec:lem:edit}
$d_{\mathrm{edit}}$ is a metric on finite real multisets, an optimal matching pairs all occurrences of the
smaller one, and for $w$ with $\norm w_\infty\le S$ and Lipschitz constant at most $L$,
$\bigl\lvert\sum_{a\in A}w(a)-\sum_{b\in B}w(b)\bigr\rvert\le\max\set{2S,L}\,d_{\mathrm{edit}}(A,B)$.
\end{lemma}

\begin{proof}
Pairing two unmatched occurrences costs at most $1$ instead of $2$. For the triangle inequality
compose matchings; where an occurrence of the middle multiset is matched on one side only, delete
the surviving endpoint, whose cost $1$ is bounded by the movement plus the deletion already
charged. A matched pair contributes at most $\min\set{2S,L\abs{a-b}}$.
\end{proof}

\begin{lemma}[Explicit coarse counts and an angular prime; \stG]\label{rec:lem:angular}
For primitive $\chi$ of conductor $q>1$ and $T\ge2$, $N_\chi(T)\le T\log(qT)$, and $N_\zeta(T)\le2T\log(2T)$; so
$N_\chi(T)\le2T\log(2qT)$ for all primitive $\chi$. If $M\ge3$, $\psi$ is a nonprincipal character modulo $M$ and
$x=10^6\log^2(2M)$, some prime $p\nmid M$ with $x/e\le p\le ex$ has $\Re\psi(p)\le\tfrac12$, hence
$\abs{1-\psi(p)}\ge1$. Consequently two distinct $\chi_1,\chi_2\in\mathcal C(Q)$ satisfy $\abs{\chi_1(p)-\chi_2(p)}\ge1$ at a
prime $p<12\cdot10^6\log^2(2Q)$.
\end{lemma}

\begin{proof}
By \cref{src:thm:dgrh} and the functional equation,
$\Re(\Lambda_\chi'/\Lambda_\chi)(\tfrac12+u)=\sum_\gamma u/(u^2+\gamma^2)$ for $u>0$, and each zero with $\abs\gamma\le T$ contributes
at least $1/(2T)$ at $u=T$. At $\sigma=T+\tfrac12$, $\abs{L'/L(\sigma,\chi)}<\tfrac12$ and $\psi_\Gamma(y)\le\log y$ give
$\Re(\Lambda_\chi'/\Lambda_\chi)(T+\tfrac12)<\tfrac12\log(qT)$, using $7e<8\pi$; for $\xi$, $\zeta'/\zeta(\sigma)<0$ gives
$\xi'/\xi(T+\tfrac12)\le\log(2T)$. For the prime, apply \eqref{rec:eq:halfweil} to the triangle
$g(v)=(1-\abs{v-\log x})_+$, weighted sums $W_\chi(x)=\sum\Lambda(n)\chi(n)n^{-1/2}g(\log n)$: the gamma term is below
$2/\sqrt x$, the zero sum at most $B(q)=20\log(2q)+21$ by the count, the pole term for $\zeta$ is
$c(\sqrt x+x^{-1/2})$ with $c=(\sinh(1/4)/(1/4))^2>1$, removing ramified Euler factors costs less than
$\tfrac52\log M$, and prime powers of exponent at least two at most $\tfrac12(1+\log x)^2+4$. If
$\Re\psi(p)>\tfrac12$ throughout the support, positivity gives $\Re W_\psi\ge\tfrac12W_{\chi_{0,M}}-\tfrac32E_{\mathrm{pp}}$, hence
$c\sqrt x\le47.5\log(2M)+89+1.5(1+\log x)^2<1000\log(2M)$, a contradiction. For the last claim apply this
to $\chi_1\overline{\chi_2}$ modulo $M=\lcm(q_1,q_2)\le Q^2$.
\end{proof}

\begin{theorem}[Comparison from a separating prime; \stG]\label{rec:thm:interface}
Let $X\ge2$, $A,B\ge1$, suppose $N_\chi(t)\le A(t+1)\log(B(t+2))$ for $t\ge2$ for the characters considered,
and put $h=1/(64X)$, $D=1+\log(B(256X+2))$, $Z=80A\sqrt XD$, with $m\ge8$ even and $m\ge\frac2{\log3}\log Z$,
$H_0=256mX$. If $\abs{\chi_1(p)-\chi_2(p)}\ge1$ at a prime $p\le X$, then for every $H\ge H_0$,
\[
 d_{\mathrm{edit}}\bigl(\Gamma_{\chi_1}(H),\Gamma_{\chi_2}(H)\bigr)\ge\frac{5\log2}{128\pi}\frac H{\sqrt X}.
\]
\end{theorem}

\begin{proof}
Let $f$ be the density of a sum of $m$ independent uniforms on $[-h/m,h/m]$, so
$\widehat f(t)=\recsinc(ht/m)^m$ and $f(0)\ge1/h$ (write $f=g*g$), and put
$\phi_H(u)=\frac{f(u)}{f(0)}\recsinc(Hu/8)^2$, $K_H=\widehat{\phi_H}$, $w_{p,H}(t)=\frac{\sqrt p}{\log p}e^{it\log p}K_H(t)$. Then
$\phi_H(0)=1$, $\int\phi_H\le\min(h,8\pi/H)$, and the support of the frequency test meets the logarithm of
one positive integer only, $p$. By \eqref{rec:eq:halfweil},
$\chi(p)=\delta_\chi P_H(p)-G_{a,H}(p)-\sum_\gamma w_{p,H}(\gamma)$, with $0\le G_{a,H}\le4h/\log2\le\tfrac1{16}$ and
$0\le P_H(p)\le32\pi X/H<\tfrac1{16}$. As $\recsinc(Hu/8)^2$ is the characteristic function of a triangular
density on $[-H/4,H/4]$, $\abs{K_H(t)}\le h(4m/(3h\abs t))^m$ for $\abs t\ge H$, and Stieltjes summation against
the count gives the tail $E(H)$, decreasing in $H$, with
$E(H_0)\le10A\sqrt XDm^23^{-m}\le\tfrac18$, since $m^2\le3^{m/2}$ and $3^{m/2}\ge Z$. So the two zero sums over
$\Gamma_{\chi_j}(H)$ differ by at least $\tfrac58$, while $\max\set{2\norm{w_{p,H}}_\infty,\operatorname{Lip}w_{p,H}}\le16\pi\sqrt X/(H\log2)$;
\cref{rec:lem:edit} concludes.
\end{proof}

\begin{theorem}[Robust identification with explicit constants; \stG]\label{rec:thm:robust}
Put $X_Q=\ceil{12\cdot10^6L_Q^2}$, $D_Q=1+\log(2Q(256X_Q+2))$,
$m_Q=2\ceil{\tfrac12\max\set{8,\frac2{\log3}\log(160\sqrt{X_Q}D_Q)}}$ and $H_Q'=256m_QX_Q=O(L_Q^2\log\log(3Q))$. For
$H\ge H_Q'$ and distinct $\chi_1,\chi_2\in\mathcal C(Q)$,
\[
 d_{\mathrm{edit}}\bigl(\Gamma_{\chi_1}(H),\Gamma_{\chi_2}(H)\bigr)\ge\Delta_Q(H)=\frac{5\log2}{128\pi}\frac H{\sqrt{X_Q}}>\frac H{10^6\log(2Q)} .
\]
Hence a finite multiset within $d_{\mathrm{edit}}$-distance $\Delta_Q(H)/2$ of some spectrum, for example after
$r$ erasures, $s$ insertions and movements $e_j$ with $r+s+\sum_j\min\set{1,\abs{e_j}}<\Delta_Q(H)/2$, identifies
the character, and $\abs{\Gamma_{\chi_1}(H)}+\abs{\Gamma_{\chi_2}(H)}-2\abs{\Gamma_{\chi_1}(H)\cap\Gamma_{\chi_2}(H)}\ge\Delta_Q(H)$.
\end{theorem}

\begin{proof}
\Cref{rec:lem:angular} gives the separating prime $p\le X_Q$ and the count with $A=2$, $B=2Q$; apply
\cref{rec:thm:interface} and $5\log2/(128\pi\sqrt{13\cdot10^6})>10^{-6}$. The rest is the triangle inequality.
\end{proof}

The readings $a_p(\widetilde\Gamma)=-\sum_{t\in\widetilde\Gamma}w_{p,H}(t)$, $p\le X_Q$, of a list $\widetilde\Gamma$ with corruption cost
at most $\eta$ satisfy $\abs{a_p-\chi(p)}\le\tfrac3{16}+16\pi\sqrt{X_Q}\eta/(H\log2)$, so retaining the candidates
compatible at every $p\le X_Q$ is a finite decoder needing neither conductor nor parity. The oriented
phase is necessary, since $\Gamma_{\bar\chi}(H)=-\Gamma_\chi(H)$. For the quartic character modulo $5$ with
$\chi(2)=i$, its conjugate and its square, numerical lists to height $40$ and a damaged list (two
erasures, three insertions, movements below $0.012$) are nearest to the true character with gap
$16.91$ (\artifact{robust_dirichlet_spectra_2026-11-26/demo/demo_results.json}; numerical, not certified
complete). Compared with \cref{rec:thm:lowzero}, this theorem has a larger height and explicit
constants and bounds the clipped cost, which charges a distant movement at most one; for quadratic
characters \cref{rec:thm:quaddisplace} gives the unclipped total displacement with the smaller
constant $1/(144\pi)$ from the somewhat lower height $2T$.

\subsection{Twist multisets from balanced zero statistics}\label{rec:sec:homometry}
\sources{dirichlet_zero_packet_homometry; dirichlet_packet_reconstruction}

A \emph{packet} is a finite nonempty multiset of characters; pooling the zeros of its $L$-functions,
let $Z_A(T)$ be the multiset of ordinates in $(0,T]$. An observer who sees only differences of
heights can measure, for $d\ge2$ and real $t_j\ne0$ with $\sum t_j=0$,
\[
 S_A(t;T)=\frac1T\sum_{\gamma\in Z_A(T)}e^{it\gamma},\qquad \mathcal B_A(\mathbf t;T)=\prod_{j=1}^dS_A(t_j;T),
\]
since a common translation of all heights cancels. Fix a modulus $Q$, $G=(\Z/Q\Z)^\times$, a character
$\psi$ modulo $Q$ and a subgroup $H\le\widehat G$ of exponent $E$; a function $A:H\to\Z_{\ge0}$ denotes the
packet with $A(h)$ copies of $\psi h$, all factors being primitive of conductor $Q$. With
$\zeta=e^{2\pi i/E}$ and $m_A=\sum_hA(h)$, put
$f_{A,a}(u)=\sum_{h:h(a)=\zeta^u}A(h)$ and $F_A(a,k)=\sum_uf_{A,a}(u)\zeta^{ku}$ for $a\in G$. By the
Landau formula of Ford and Zaharescu \cite[Lemma~2.1]{FordZaharescu2015Unnormalized} and
\cref{src:thm:dgrh}, $S_A(t;T)=L_A(t)+O_{A,t}(\log(2T)/T)$, where
\begin{equation}\label{rec:eq:landau}
 L_A(k\log p)=-\frac{\log p}{2\pi p^{\abs k/2}}\psi(p)^kF_A(p,k)\qquad(p\nmid Q,\ k\in\Z\setminus\set0)
\end{equation}
and $L_A(t)=0$ at all other nonzero frequencies.

\begin{theorem}[Local classification; \stG]\label{rec:thm:localclass}
The following are equivalent: \textup{(i)} the limits of all balanced measurements agree for $A,B$;
\textup{(ii)} they agree for the finitely many tuples at a single prime $p$, of orders
$2\le d\le\max(2,E)$, one prime in each reduced class; \textup{(iii)} for each $a\in G$, $f_{B,a}$ is a
cyclic rotation of $f_{A,a}$.
\end{theorem}

\begin{proof}
Given (iii), a tuple with a slot outside the support of \eqref{rec:eq:landau} has zero limit for both;
otherwise $t_j=k_j\log p_j$ and unique factorization gives $\sum_{p_j=p}k_j=0$ at each prime, so the
rotation factors cancel. Conversely, $(E\log p,-E\log p)$ gives $m_A$; at the chosen prime the tests
give every product $\prod_jF(b_j)$ with $\sum b_j\equiv0\pmod E$, $d\le E$. For real $y$ on $\Z/E\Z$ the
power sums $\sum_sz_s^d$ of $z_s=\sum_uy_uf(u-s)$ are combinations of these products, so Newton's
identities give the multiset $\set{z_s}$; with $y_u=(m_A+1)^u$ each $z_s$ encodes one rotation in base
$m_A+1$.
\end{proof}

\begin{corollary}[Coprime packet size; \stG]\label{rec:cor:coprime}
If $\gcd(m_A,E)=1$, all balanced data agree if and only if $B$ is obtained from $A$ by multiplying
every member by one character of $H$.
\end{corollary}

\begin{proof}
With $D_A=\prod_hh^{A(h)}$ the local rotations $v_a\in\mu_E$ satisfy $v_a^{m_A}=(D_B/D_A)(a)$; if
$m_Au\equiv1\pmod E$, then $v_a=(D_B/D_A)^u(a)$ for all $a$, and orthogonality on $G$ gives the
multiplicities.
\end{proof}

For quadratic twists ($E=2$) let $c_A(h)=\sum_uA(u)A(hu)$ count pairwise quotients. Then
$\widehat c_A=\abs{\widehat A}^2$ on $\widehat H$, every element of $\widehat H$ is evaluation at some $a\in G$,
and a two-bin histogram of known mass is determined up to a swap by the difference of its bins;
hence all balanced limits agree if and only if all pair limits agree, if and only if $c_A=c_B$
(homometry). For packets satisfying (iii) and real $w\in C_c^\infty(\R\setminus\set0)$ with $\int w=0$,
the pair statistic $\sum_{\gamma,\gamma'\in Z_A(T)}w(\gamma-\gamma')$ of the two packets differs by
$O(T/(\log T)^{1/3})$, by the cross-character formula of
\cite[Theorem~2]{FordZaharescu2015Unnormalized}, whose main term depends on
$\Lambda(n)^2\abs{\sum_{\chi\in A}\chi(n)}^2$, unchanged by local rotation.

\begin{theorem}[Twenty-eight indistinguishable twist classes; \stG]\label{rec:thm:six}
Let $Q=5\cdot13\cdot17\cdot29=32045$, $\psi_i$ a character of order four modulo $q_i$,
$\psi=\prod\psi_i$ and $h_x=\prod_i\psi_i^{2x_i}$ for $x\in\F_2^4$, so every $\psi h_x$ is primitive of
conductor $Q$. The six-factor packets supported where $x_1x_2+x_3x_4=1$ and where
$x_1x_2+x_3x_4+x_1x_3=1$ have the same balanced limits at every order. Exactly $448$ six-element
subsets of $\F_2^4$ share their data, forming $28$ classes under a common twist, and
\[
 \sum_{0<\gamma\le T}\abs{m_A(\gamma)-m_B(\gamma)}\ge\frac{\sqrt2\log2}\pi T-O(\log(2T)),
\]
the same bound holding with $H$ in place of $T$ on every window $(T,T+H]$, up to $O(\log(2T+2H))$.
\end{theorem}

\begin{proof}
The Walsh transforms of $(-1)^{q_A}$ and $(-1)^{q_B}$ are $4(-1)^{v_1v_2+v_3v_4}$ and the same times
$(-1)^{v_2v_4}$, so both sets have size six and all nontrivial indicator coefficients of modulus two,
and homometry applies; the bilinear parts differ, and twists act by translation. Difference counts
$6$ at zero and $2$ elsewhere mean that the $15$ pair sums are distinct and nonzero; such a six-set
sums to zero with every five points affinely independent, and conversely, so there are
$16\cdot15\cdot14\cdot12\cdot8/6!=448$ of them, in translation orbits of size $16$. At $p\nmid Q$ the
coefficients differ exactly when $(p/q_2)=(p/q_4)=-1$, for one quarter of the classes and for $p=2$,
with difference four; apply \eqref{rec:eq:landau} at $\log2$.
\end{proof}

\begin{theorem}[Uniformity in the frequencies; \stG]\label{rec:thm:homuniform}
If the local histograms of $A$ and $B$ are rotations of each other at every class, then for fixed
$d\ge2$, $a>0$, $0<\delta<2$ and $\eps>0$,
\[
 \sup_{\sum t_j=0,\ a\le\abs{t_j}\le(2-\delta)\log T}\abs{\mathcal B_A(\mathbf t;T)-\mathcal B_B(\mathbf t;T)}
 \ll T^{-\delta/2+\eps}.
\]
\end{theorem}

\begin{proof}
With $X=T^{2-\delta}$, $x=e^{\abs t}$ and $n$ a nearest integer to $x$, the uniform Landau formula
approximates $S_A(t;T)$ by $0$ unless $n=p^k$ is an unramified prime power, and then by
$-\frac{\log p}{2\pi\sqrt x}\psi(p)^{sk}F_A(p,sk)K_T(t-s\log n)$ with $K_T(u)=T^{-1}\int_0^Te^{ivu}\dd v$,
$s=\operatorname{sgn}t$, with error $O(X^{1/2+\eta}\log(2XT)/T)$. If the products $U,V$ of the $n_j$
over positive and negative $t_j$ agree, the leading products agree; otherwise some detuning is at
least $\abs{\log(U/V)}/d\ge1/(d(1+\min(U,V)))$, and $\abs{K_T(u)}\le\min(1,2/(T\abs u))$ bounds both
products by $O((\log2X)^d/T)$.
\end{proof}

\begin{theorem}[One cubic measurement at odd prime exponent; \stG]\label{rec:thm:cubic}
Let $E=\ell$ be an odd prime and choose one prime $p\equiv a\pmod Q$ for each $a\in G$. The limit at
$(\log p,\log p,-2\log p)$ determines $f_{A,a}$ up to rotation; if $\ell\nmid m_A$, these measurements
determine the packet up to a common twist.
\end{theorem}

\begin{proof}
With $\alpha_f=\sum_uf(u)\zeta^u$ and $\sigma(\zeta)=\zeta^2$, the normalized cubic is
$C(f)=\alpha_f^2\overline{\sigma(\alpha_f)}$. If $C(f)=0$ all entries are equal. Otherwise, for
$C(g)=C(f)$ and $b=\alpha_g/\alpha_f$, $\sigma(b)=\bar b^{-2}$, so the radii along the
$\sigma$-orbit satisfy $r_{j+1}=r_j^{-2}$ and are all $1$; then $\sigma(b)=b^2$, $b$ is a root of
unity of odd order in $\Q(\zeta)$, hence $b=\zeta^s$, and $g$ minus the rotated $f$ is constant, hence
zero. The last claim is \cref{rec:cor:coprime}.
\end{proof}

A nonzero difference of two cubics is an algebraic integer with conjugates at most $2m^3$, so it
has modulus at least $(2m^3)^{-(\ell-3)/2}$. The quadratic data cannot replace the cubic: one copy
at $0$ and two at $1$, and its reflection, have equal pair data but are not rotations of each other.
More collisions come from the Maiorana--McFarland bent functions \cite{Dillon1974Elementary}: with
$2r$ primes $q_i\equiv1\pmod4$ and quadratic twists identified with $\F_2^r\times\F_2^r$, the supports
of $x\cdot y+g(y)$, $g(0)=0$, all have the same balanced data, and they fall into exactly
$2^{2^r-1-r}$ common-twist classes; for $r\ge3$ the supports for $g=0$ and $g=y_1y_2y_3$ are not even
affinely equivalent. The finite enumerations are
\artifact{dirichlet_zero_packet_homometry/packet_checks.txt} and
\artifact{dirichlet_zero_packet_homometry/cubic_checks.txt}.

\paragraph{The least number of local measurements.}
At one unramified prime the balanced limits give products $\prod_iF(k_i)$ with $\sum_ik_i\equiv0\pmod E$ of the
Fourier coefficients $F(k)=\sum_xf(x)\zeta_E^{kx}$ of the local histogram $f:C_E=\Z/E\Z\to\Z_{\ge0}$, of known mass
$m=\sum_xf(x)$; rotation multiplies $F(k)$ by $\zeta_E^{kt}$, and the question is how many individual products,
of what orders $s$, determine $f$ up to rotation. Let $\mathcal D=\set{d\mid E:d>1}$, and $\mathcal I$ the set of unordered
pairs $\set{d,e}\subset\mathcal D$ with $\gcd(d,e)>1$.

\begin{theorem}[Exact measurement counts; \stU]\label{rec:thm:probecount}
The least size of a fixed separating list of products, and of a deterministic adaptive procedure,
in the worst case, are
\[
 N_{\mathrm{fixed}}=\abs{\mathcal D}+\abs{\mathcal I}=\tfrac12\bigl(\tau(E)^2+\tau(E)-1-\tau(E^2)\bigr),\qquad N_{\mathrm{adaptive}}=\abs{\mathcal D}=\tau(E)-1 .
\]
A fixed list of minimum size exists with every product of order at most six, and the lower bounds hold
with unrestricted orders, even at the single mass $m=2E^2$.
\end{theorem}

For $E=p^r$ the counts are $r(r+1)/2$ and $r$; for $E=30$ they are $22$ and $7$.

\begin{lemma}[Products on one Galois orbit; \stU]\label{rec:lem:probefactor}
For $d\in\mathcal D$ put $k_d=E/d$, $\alpha_d=F(k_d)\in K_d=\Q(\zeta_d)$, so $F(uk_d)=\sigma_u(\alpha_d)$ for
$u\in(\Z/d\Z)^\times$. If $a>b\ge0$ and $P(\alpha)=\alpha^a\prod_{j\le b}\sigma_{u_j}(\alpha)$, then $P(\alpha)=0$ only for $\alpha=0$, and
$P(\alpha)=P(\beta)\ne0$ forces $\beta/\alpha$ to be a root of unity of $K_d$, lying in $\mu_d$ if $d$ and $a+b$ are odd.
\end{lemma}

\begin{proof}
Applying all $\sigma_v$ to $P(z)=1$, $z=\beta/\alpha$, gives multiplicative relations with exponent matrix
$aI+\sum_jP_j$, $P_j$ permutation matrices, which is nonsingular by the sup norm; the adjugate gives
$z^{\det}=1$. The roots of unity of $K_d$ are $\mu_d$ for even $d$ and $\mu_{2d}$ for odd $d$, and for odd $d$ an
extra root $-\zeta_d^t$ has image in $-\mu_d$ when $a+b$ is odd.
\end{proof}

\begin{proof}[Proof of \cref{rec:thm:probecount}]
Every residue modulo odd $d$, and every even residue modulo even $d$, is a sum of two units, by
the Chinese remainder theorem. \emph{Pure measurements.} For $d=2$ use the units $(1,-1)$, for odd $d$
$(1,1,-2)$, for even $d>2$ with $3\nmid d$ $(1,1,1,-3)$, and for $6\mid d$ $(1,1,1,1,u,v)$ with $u+v=-4$. By
\cref{rec:lem:probefactor} equal values give equal activity of each order and a unique relative phase
$t_d\bmod d$ with $\beta_d=\zeta_d^{t_d}\alpha_d$. \emph{Bridges.} For $\set{d,e}\in\mathcal I$, $g=\gcd(d,e)$, take units of order $d$
summing to $d/g$ and of order $e$ summing to $-e/g$, with $\ell(n,g)$ units on the side of order $n$:
$\ell=1$ if $n=g$, $3$ if $n>g$ is even with $n/g$ odd, and $2$ otherwise. The product enforces
$t_d\equiv t_e\pmod g$ when both orders are active, and the generalized Chinese remainder theorem gives
a common rotation; Fourier inversion recovers $f$. \emph{Lower bounds.} With the Ramanujan sums
$C_d(x)=\sum_{u\in(\Z/d\Z)^\times}\zeta_d^{ux}$ and $M=2E$: distinguishing $M$ from $M+C_d$ needs a product all of
whose indices have order $d$, also along any adaptive path; and distinguishing $M+C_d(x)+C_e(x)$ from
$M+C_d(x)+C_e(x-1)$, for $\gcd(d,e)>1$, needs a product whose set of orders is exactly $\set{d,e}$. All
these histograms have mass $2E^2$. \emph{Adaptive upper bound.} Process the divisors in a fixed order;
with $L=\lcm$ of the processed active orders and $g=\gcd(d,L)$, $h=d/g$, ask for $F(k_d)^hB$, with $B$ a
nonzero product of known coefficients of index sum $-E/g$. This determines $\alpha_d$ up to $\mu_h$, exactly
the ambiguity left by rotations by multiples of $L$.
\end{proof}

\begin{theorem}[Least maximum order; \stU]\label{rec:thm:probedegree}
For $E>1$ the least possible maximum order over all separating lists is
$\max(\set{c_d:d\in\mathcal D}\cup\set{\ell(d,g)+\ell(e,g):\set{d,e}\in\mathcal I})$, with $c_d=2,3,4,6$ the lengths of the pure
rows; the fixed list above attains both minima at once. The optimal adaptive count is attained
with orders at most $6+3\omega(E)$, with order three for odd prime powers and four for powers of two.
\end{theorem}

\begin{proof}
Inputs with arbitrary nonzero entries on one order are realized by the histograms
$f_\alpha(x)=M+\Tr_{K_d/\Q}(A\alpha\zeta_d^{-x})$, for large integers $A$ and $M$. For $d>2$ compare $\alpha=1$ with $b\in K_d$ of
modulus one at every embedding and not a root of unity (for example $(n+\zeta_d)/(n+\zeta_d^{-1})$): all
balanced quadratic products agree, and for even $d$ balanced pure products have even length. For
$6\mid d$ take $b\in K_6$, on which a unit acts as $b^{\pm1}$ according to its class modulo $6$, so that a
balanced row of length below six cannot see $b$. For the pairs shift the $e$-part by $1$ or $g/2$; the
partial index sum of the $d$-part must be a nonzero multiple of $E/g$, which needs $\ell(n,g)$ units on
each side. For the adaptive orders replace $F(k_d)^h$ by $\alpha_d^a\sigma_u(\alpha_d)\sigma_v(\alpha_d)$ with $a\in\set{3,4}$ and
$h-a=u+v$, and decompose $-E/g$ into primary components, each a sum of at most three indices of an
active order.
\end{proof}

\begin{proposition}[When the mass removes the ambiguity; \stU]\label{rec:prop:probemass}
If the convolution matrix $A_f$ has $\gcd(\det A_f,E)=1$, the pure measurements determine $f$ up to rotation
among integral histograms of the same mass; this holds for $E=p^r$ whenever $p\nmid m$. For $E=p^r$, $r\ge2$,
the single pure measurement at the indices of order $E$ separates all rotation classes of mass $m$ exactly
when $m\le p$.
\end{proposition}

\begin{proof}
If $b$ has the same pure measurements, the ratio $G/F$ takes values in $\mu_E$; its inverse transform $v$
satisfies $A_fv=b$, so $(\det A_f)v$ and $Ev$ are integral, hence $v\in\Z^E$, and Parseval with $\sum v=1$
makes $v$ a point mass. For $E=p^r$, $\mathbb F_p[X]/((X-1)^E)$ has $f$ as a unit when $p\nmid m$. For the
last claim the aligned difference is divisible by $\Phi_E(X)=\sum_{i<p}X^{iE/p}$, hence constant on fibres
modulo $E/p$, and a nonzero difference costs at least $p$ units of mass on each side; for $m>p$ the
histograms $(m-p)\mathbf 1_{\set0}+\mathbf 1_U$ and $(m-p)\mathbf 1_{\set0}+\mathbf 1_V$ for distinct fibres $U,V$ collide.
\end{proof}

\begin{corollary}[Logarithmic windows of zeros; \stG]\label{rec:cor:probewindow}
In the setting of \cref{rec:thm:localclass}, fix $Q,p\nmid Q,E$ and the mass $m$. There is $C$ such that for
all large $T$ the fixed list of \cref{rec:thm:probecount}, evaluated on the relative ordinates of the pooled
positive zeros in every window $(T,T+C\log(2QT))$, determines the local histogram at $p$ up to rotation.
\end{corollary}

\begin{proof}
With $S(t)=Y^{-1}\sum_{T<\gamma<T+Y}e^{it\gamma}$, the Landau formula of Banks \cite{Banks2023Comparing} under
\cref{src:thm:dgrh} gives $S(k\log p)=-\frac{\log p}{2\pi p^{\abs k/2}}\psi(p)^kF(k)+O(\log(2QT)/Y)$ for fixed $k\ne0$. Lift each
row to nonzero integers with sum exactly zero; the product of the $S(k_j\log p)$ depends only on
differences of ordinates and the factors $\psi(p)^{k_j}$ cancel. Distinct values of a product of length
$s\le6$ differ by at least $(2m^s)^{-(\varphi(E)-2)/2}$ for $E>2$, as a nonzero algebraic integer of $K_E$ with
conjugates at most $2m^s$, and by at least one for $E\le2$; choose $C$ to make the errors smaller than
half this and round.
\end{proof}

All lower-bound examples occur arithmetically: for a prime $Q\equiv1\pmod{2E}$, a generator $\psi$ of the
characters modulo $Q$, $H$ of order $E$ and $p$ a primitive root modulo $Q$, every histogram is realized by
primitive twists $\psi h$. For $E=30$, $Q=61$, $p=2$, the histograms $60+C_6(x)+C_{10}(x)$ and
$60+C_6(x)+C_{10}(x-1)$ agree in every product of order at most five, while the sextic with indices
$(5,5,5,3,3,-21)$ gives $\pm729{,}000{,}000$:
\begin{center}
\artifact{dirichlet_packet_reconstruction/zero-packets/check_probes.txt}\\
\artifact{dirichlet_packet_reconstruction/zero-packets/rigidity_output.txt}
\end{center}

\subsection{Matching Euler factors of products of quadratic $L$-functions}\label{rec:sec:qmatch}
\sources{dirichlet_packet_reconstruction}

Let $X$ be a group of characters modulo $Q$ of order dividing two, $r=\dim_{\F_2}X\ge1$, and $V=\widehat X$,
written additively. A packet $A$ is a multiset of $k\ge1$ members of $X$, and at $p\nmid Q$ its local
polynomial $\prod_{\chi\in A}(1-\chi(p)T)$ is determined by the trace $\sum_{\chi\in A}\chi(p)$. For distinct packets $A,C$ let
$\delta(A,C)$ be the density of primes with different local polynomials. By the prime number theorem in
progressions,
\[
 \delta(A,C)=2^{-r}\abs{\supp H},\qquad H(v)=\sum_{\chi\in X}c_\chi\chi(v),\quad c_\chi=\operatorname{mult}_A(\chi)-\operatorname{mult}_C(\chi),
\]
with $H\ne0$, $H(0)=0$, $B=\sum\abs{c_\chi}\le2k$, and after cancellation each packet has size $B/2$.

\begin{lemma}[Dyadic valuation; \stU]\label{rec:lem:walsh}
For a nonzero integral combination $H$ of Walsh characters of $\F_2^r$ and $t=\min_{H(x)\ne0}v_2(H(x))$, the map
$x\mapsto2^{-t}H(x)\bmod2$ is a nonzero Boolean polynomial of degree at most $t$. A nonzero Boolean
polynomial of degree at most $s\le r$ has relative weight at least $2^{-s}$, with equality only for the
indicator of an affine flat.
\end{lemma}

\begin{proof}
In Boolean coordinates each Walsh character is a product of factors $1-2x_i$, so a monomial of degree
$j$ has coefficient divisible by $2^j$, and by $2^t$ by M\"obius inversion; divide and reduce. For the
weight write $f_0+x_rf_1$ and induct. At equality, translates of the support coincide or are
disjoint, since a nonzero Boolean difference of degree at most $s-1$ has weight at least $2^{1-s}$.
\end{proof}

\begin{theorem}[Exact minimum and the high-agreement spectrum; \stU]\label{rec:thm:qmatch}
Put $D=\floor{\log_2(2k)}$ and $d=\min(r,D)$. The least positive disagreement density at rank $r$ and size
$k$ is $2^{-d}$. If $r>D$, equality holds exactly when $H=\pm2^d\eta\mathbf 1_{a+W}$ with $\eta\in X$ and an affine
subspace $a+W\not\ni0$ of codimension $d$; if $r\le D$, exactly when $H=m2^r\mathbf 1_{\set a}$, $a\ne0$, $\abs m2^r\le2k$.
If $r>D$, the densities strictly below $2^{1-d}$ are exactly
\[
 \delta_j=2^{1-d}(1-2^{-j}),\qquad1\le j\le\min(d,r-d),\quad k\ge2^d-2^{d-j},
\]
and every pair realizing $\delta_j$ reduces after cancellation to two disjoint sets of exactly $2^d-2^{d-j}$
distinct characters; if $r\le D$ only the minimum occurs below twice the minimum.
\end{theorem}

\begin{proof}
Since $\abs H\le2k$, $t\le D$, and \cref{rec:lem:walsh} gives $\delta\ge2^{-t}\ge2^{-d}$. At equality all nonzero values
are $\pm2^d$ on an affine flat $a+W$, $H(a+w)=2^ds(w)$, and Fourier inversion makes every normalized
coefficient of $s$ integral; by Parseval exactly one is $\pm1$, so $s$ is a signed character. The
converse uses the two halves of the Fourier support, padded by common entries. Below twice the
minimum $t=d$ and the Boolean polynomial is the disagreement indicator, whose weights are given by
the Kasami--Tokura classification of Reed--Muller weights below twice the minimum
\cite{KasamiTokura1970Weight}. Then $\abs{c_\chi}\le2^d\delta<2$, so $B=\sum c_\chi^2=4^d\delta$ and $B/2=2^d-2^{d-j}\le k$, and
$\abs{\supp H}$ is even, forcing $j\le r-d$. Two $d$-dimensional subspaces $U,U'\subset X$ meeting in dimension $d-j$
give $H=2^d(\mathbf 1_{U^\perp}-\mathbf 1_{(U')^\perp})$ with density $\delta_j$. For $r\le D$ densities are multiples of $2^{-r}$.
\end{proof}

For example, for $r\ge6$ two products of seven factors can agree at $25/32$ of the primes, but not at
any proportion strictly between $25/32$ and $13/16$. The classification is checked on $723{,}057$ pairs
and $95$ constructions through rank ten:
\begin{center}
\artifact{dirichlet_packet_reconstruction/prime-traces/verification.txt}
\end{center}

\begin{theorem}[Finite detection and density recovery; \stG]\label{rec:thm:qmatchfinite}
For a smooth $w\ge0$ supported in $(\tfrac12,1)$ with integral one and $f:V\to\C$, let
$M_x(f)=\sum_{p\nmid Q}(\log p)w(p/x)f([p])$. Then $\abs{M_x(f)-x\mathbb E f}\ll_w\norm f_A\sqrt x\log(2Q)$, where
$\norm f_A=\sum_\chi\abs{\widehat f(\chi)}$. Every distinct pair has a distinguishing prime in $(x/2,x)$ for
$x\ge C_wk^2\log^2(2Q)$; if $r>D$ and $0<\delta<2^{1-d}$, the disagreement set $S$ has
$M_x(\mathbf 1_S)=x\delta+O_w(\sqrt x\log(2Q))$ uniformly in $k,r,Q$, and rounding $4^dM_x(\mathbf 1_S)/(2x)$ recovers $B/2$, hence
$\delta$, for $x\ge C_w'16^d\log^2(2Q)$.
\end{theorem}

\begin{proof}
The first estimate is \cref{rec:lem:gram} after Fourier expansion. Apply it to $H^2$, with
$\norm{H^2}_A\le B^2$ and $\mathbb EH^2\ge B$. In the high-agreement range $\mathbf 1_S=4^{-d}H^2$, so $\norm{\mathbf 1_S}_A<4$, and
$4^d\delta/2=B/2$ by \cref{rec:thm:qmatch}.
\end{proof}

\subsection{Shortest positive averages that cancel prescribed zeros}\label{rec:sec:dilation}
\sources{quadratic_spectra; logarithmic_dirichlet_zero_filters_2026-11-26}

A probability measure $\mu$ on the logarithmic scale acts on a power $x^s$ by the multiplier
$F_\mu(s)=\int e^{su}\dd\mu(u)$, and a positive average preserves inequalities between the averaged
prime sums. How short can it be if it must annihilate prescribed zero contributions? Fix $\tau>0$,
distinct reals $\beta_1,\dots,\beta_r$, positive integers $m_j$, $s_j=\beta_j+i\tau$, $m=\sum m_j$,
$h=\pi/\tau$, $b=mh$, and require of a compactly supported probability measure $\mu$
\begin{equation}\label{rec:eq:fjets}
 F_\mu^{(k)}(s_j)=\int u^ke^{s_ju}\dd\mu(u)=0\qquad(1\le j\le r,\ 0\le k<m_j).
\end{equation}

\begin{theorem}[Minimum width; \stU]\label{rec:thm:fwidth}
Every solution of \eqref{rec:eq:fjets} has support diameter at least $m\pi/\tau$, and up to
translation exactly one has that diameter: $\mu_*=\sum_{\ell=0}^mw_\ell\delta_{\ell h}$ with
$\sum_\ell w_\ell z^\ell=\prod_j\bigl((z+e^{h\beta_j})/(1+e^{h\beta_j})\bigr)^{m_j}$.
\end{theorem}

\begin{proof}
The space $V=\operatorname{span}_\R\set{u^ke^{\beta_ju}:k<m_j}$ has dimension $m$, and a nonzero element
has at most $m-1$ real zeros with multiplicity (divide by $e^{\beta_1u}$, differentiate, use Rolle).
So there is a unique-up-to-scalar $q\in V$ vanishing at $h,\dots,(m-1)h$, with no other real zeros,
and with a sign choice $g(u)=q(u)\sin(\tau u)\ge0$ on $[0,b]$, vanishing exactly on
$\set{0,h,\dots,b}$, simply at the ends and doubly inside. If a solution has diameter at most $b$,
translate it to $[0,b]$; the imaginary parts of \eqref{rec:eq:fjets} give $\int g\dd\mu=0$, so $\mu$ lives
on the grid, $F_\mu(s)=P(e^{hs})$, and $P$ vanishes to order $m_j$ at $-e^{h\beta_j}$; with $P(1)=1$ this
is the stated polynomial, whose coefficients are positive, including both endpoint masses.
\end{proof}

With the $2m$ real coordinates $Y(u)=(\Re,\Im)(u^ke^{s_ju})_{j,k}$, a norm $\norm\cdot$ with dual
$\norm\cdot_*$ and $\eps(\mu)=\norm{\int Y\dd\mu}$, the extremizer is stable and its interval cannot be
shortened cheaply.

\begin{proposition}[Sharp stability; \stU]\label{rec:prop:fstable}
For probability measures on $I=[0,b]$ with $\eps(\mu)\le1$, $W_1(\mu,\mu_*)\le C\eps(\mu)^\alpha$ with
$\alpha=1$ for $m=1$ and $\alpha=1/2$ for $m\ge2$, and these exponents are best possible.
\end{proposition}

\begin{proof}
With $g=v\cdot Y$ as above, $g(u)\ge c\operatorname{dist}(u,G)^2$ on $I$ (or $\ge c\operatorname{dist}(u,G)$ if
$m=1$), and $\int g\dd\mu\le\norm v_*\eps$, so rounding to the grid costs $O(\eps^\alpha)$ (Cauchy--Schwarz
in the quadratic case) and moves the moments by $O(\eps^\alpha)$. Total mass and the real parts of the
jets determine signed grid weights, by the root argument, so the grid weights are within
$O(\eps^\alpha)$ of $w_\ell$. Splitting an interior atom by $\pm\eta$ gives $\eps=O(\eta^2)$ and transport
cost $\asymp\eta$; for $m=1$ moving an endpoint atom gives both of order $\eta$.
\end{proof}

\begin{theorem}[First-order cost of shortening; \stU]\label{rec:thm:fleak}
For $0<d<b$ let $D(d)=\min\set{\eps(\mu):\supp\mu\subset[d/2,b-d/2]}$. Then
\[
 \lim_{d\downarrow0}\frac{D(d)}d=\frac{w_0g'(0)-w_mg'(b)}{2\norm v_*}>0 .
\]
\end{theorem}

\begin{proof}
By convex duality, $D(d)=\max_{\norm z_*\le1}\min_{u\in I_d}z\cdot Y(u)$. A direction with $z\cdot Y\ge0$
on $I$ integrates to zero against $\mu_*$, so it vanishes at the nodes with zero derivative inside;
writing $z\cdot Y=p\cos\tau u+r\sin\tau u$ with $p,r\in V$ forces $p=0$ and $r\in\R_{\ge0}q$, so the
supporting directions form the ray of $v$. With $a_0=g'(0)/2$, $a_m=-g'(b)/2$, $a_\ell=0$ otherwise and
$K=\sum w_\ell a_\ell>0$, evaluating near-maximizers on the grid with the two endpoints moved inward
gives $\limsup D(d)/d\le K$. Conversely, grid evaluation $V\to\R^{m+1}$, $p\mapsto((-1)^\ell p(\ell h))$, has
image the hyperplane orthogonal to $(w_\ell)$, so one can prescribe $(-1)^\ell p(\ell h)=K-a_\ell$; then
$g+dp\cos\tau u$ is at least $dK-O(d^2)$ on $I_d$ and has dual norm $1+O(d)$.
\end{proof}

For $\tau=\pi$, $(\beta_1,\beta_2)=(0,\log2)$ and simple conditions the weights are $(1/3,1/2,1/6)$, the
certificate is $(2-2^u)\sin\pi u$, and for the maximum of the two complex moduli $D(d)=(\pi/9)d+o(d)$.
The width bound alone also follows from Obrechkoff's sector inequality
\cite{EremenkoFryntov2016Remarks}.

\begin{theorem}[A block of consecutive harmonics; \stU]\label{rec:thm:hblock}
Fix integers $p,H\ge1$, $\tau>0$, distinct $\beta_j$ and $m_j\ge1$, and put $B=2p+H-1$,
$h=2\pi/(B\tau)$, $n=H\sum m_j$, $a_j=e^{h\beta_j}$, $\omega=e^{2\pi i/B}$. If a compactly supported
probability measure satisfies $F_\mu^{(k)}(\beta_j+i\ell\tau)=0$ for $p\le\ell\le p+H-1$, $k<m_j$ and all
$j$, its support has diameter at least $nh$, and the unique minimizer up to translation has masses
$w_t$ at $th$, $0\le t\le n$, with
\[
 \sum_tw_tz^t=\frac{P(z)}{P(1)},\qquad P(z)=\prod_j\Bigl[\prod_{\ell=p}^{p+H-1}(z-a_j\omega^\ell)\Bigr]^{m_j}.
\]
For $p=H=1$ this is \cref{rec:thm:fwidth}.
\end{theorem}

\begin{proof}
Each inner factor has positive coefficients: with $\theta_s=(2s-H+1)\pi/B$, it is $a_j^HR(z/a_j)$ for
$R(y)=\prod_s(1+ye^{i\theta_s})=\sum c_ty^t$, and $(1+\alpha y)R(\omega y)=(1+\alpha\omega^Hy)R(y)$ gives
$c_t/c_{t-1}=\sin((H-t+1)\pi/B)/\sin(t\pi/B)>0$. For the converse, with $u=x/h$ and
$Q(z)=\prod_j\prod_s(z-a_je^{i\theta_s})^{m_j}$, let $E(u)$ be the sum of the residues of $z^{u-1}/Q(z)$ and
$g=(-1)^{n+1}\sin(\pi u)E(u)$, a real combination of the test functions. A keyhole contour gives, for
$0<u<n$, $E(u)=\frac{(-1)^{n+1}\sin\pi u}\pi\int_0^\infty t^{u-1}\abs{Q(-t)}^{-1}\dd t$, so $g\ge0$ with double
zeros at interior nodes and simple ones at $0,n$. Hence $\mu$ lives on the grid and its mass
polynomial is a multiple of $P$.
\end{proof}

Stability with exponents $1/2$ (for $n\ge2$) and $1$, and the shortening law, hold as before; for the
residue certificate $g'(0)=\pi/(h\abs{Q(0)})$, $g'(b)=-\pi/h$ and $w_0/w_n=\abs{Q(0)}$, so the slope is
$\pi w_n/(h\norm v_*)$. For $p=1$, $H=2$, $h=1$, $e^{\beta_1}=1/2$, $e^{\beta_2}=2$ with multiplicities $1,2$
the optimal width is $6$, and the masses are
\[
 (w_0,\dots,w_6)=\tfrac1{343}(16,48,108,92,57,18,4),
\]
which are the coefficients of $(4z^2+2z+1)(z^2+2z+4)^2/343$. The exact check is recorded in
\artifact{quadratic_spectra/filters_example.txt}.

For prime sums, let $W$ be smooth with compact support in the interval $(1,2)$, and put
$S_W(x)=\sum_n\Lambda(n)W(n/x)$. A zeta zero contributes $-x^\rho\widetilde W(\rho)$ to $S_W$, with
$\widetilde W$ the Mellin transform, and a positive average of $S_W(xe^u)$ multiplies this by
$F_\mu(\rho)$. For $\rho=\tfrac12+i\gamma$, $\gamma>0$, and $h=\pi/\gamma$, the average
\[
 \frac{e^{h/2}S_W(x)+S_W(e^hx)}{1+e^{h/2}}
\]
cancels the conjugate pair; its $m$-fold iterate has a zero of
order $m$ at $\rho$, logarithmic span $mh$, minimal by \cref{rec:thm:fwidth}, and multiplier
$O(\abs\eta^m)$ at a zero displaced by $\eta$.

\paragraph{Many frequencies on the critical line: logarithmic support.}
The theorems above prescribe one frequency $\tau$ (or a block of harmonics) and several real parts. The
complementary problem prescribes all zero ordinates up to a height $B$ on the critical line. Here
$\widehat\mu(z)=\int e^{izt}\dd\mu(t)$, and a zero of order at least $r$ at $\lambda$ means
$\widehat\mu^{(j)}(\lambda)=0$ for $j<r$. For a finite symmetric multiset $\Lambda\subset\R\setminus\set0$ write $m_\Lambda$ for its
multiplicity function.

\begin{lemma}[A positive interpolating density; \stU]\label{rec:lem:sepfilter}
Let $\Gamma\subset\R\setminus\set0$ be finite and symmetric with the distinct points of $\Gamma\cup\set0$ at least $\delta>0$
apart, $A=4\pi/\delta$, $g_A(t)=A^{-1}(1-\abs t/A)_+$ and $f_A(z)=(\sin(Az/2)/(Az/2))^2$. There is an even
probability density $w_\Gamma$ on $[-A,A]$ with $\tfrac56g_A\le w_\Gamma\le\tfrac65g_A$ and $\widehat w_\Gamma=0$ on $\Gamma$.
\end{lemma}

\begin{proof}
With $M_{\lambda\kappa}=f_A(\lambda-\kappa)$, $b_\lambda=f_A(\lambda)$ and $0\le f_A(x)\le4/(A^2x^2)$, separation gives
$\norm{M-I}_1\le\frac8{A^2\delta^2}\sum_j j^{-2}=\tfrac1{12}$ and $\norm b_1\le\tfrac1{12}$, so $a=-M^{-1}b$ has $\norm a_1\le\tfrac1{11}$, is real,
and satisfies $a_{-\lambda}=a_\lambda$ by uniqueness. Then $Q(t)=1+\sum_\lambda a_\lambda e^{-i\lambda t}\in[\tfrac{10}{11},\tfrac{12}{11}]$, its
$g_A$-weighted integral $Z$ lies in the same range, and $w_\Gamma=g_AQ/Z$ has
$\widehat w_\Gamma(\kappa)=Z^{-1}(f_A(\kappa)+\sum_\lambda a_\lambda f_A(\kappa-\lambda))=0$.
\end{proof}

\begin{theorem}[A positive filter with explicit support; \stU]\label{rec:thm:finitefilter}
Let $\Lambda\subset[-B,B]\setminus\set0$ be a finite symmetric multiset whose positive elements are at least $\delta>0$,
with local occupancy $K=\sup_{x\ge0}\sum_{x<\lambda\le x+\delta,\lambda>0}m_\Lambda(\lambda)$. For every integer $N\ge0$ there is an even
probability measure $\mu$ on $[-H,H]$, $H=4\pi K/\delta+N/B$, with $\ord_\lambda\widehat\mu\ge m_\Lambda(\lambda)$ on $\Lambda$ and
$\abs{\widehat\mu(\omega)}\le(\sin1)^N$ for real $\abs\omega\ge B$.
\end{theorem}

\begin{proof}
Colour the positive occurrences in nondecreasing order greedily with $K$ colours, reusing a colour only
at distance at least $\delta$; equal occurrences get different colours. Reflect each colour class to a
symmetric set, apply \cref{rec:lem:sepfilter}, and convolve, which adds the vanishing orders. Convolving
further with $N$ copies of the uniform measure on $[-1/B,1/B]$ multiplies the transform by
$(\sin(\omega/B)/(\omega/B))^N$, and $\sup_{\abs x\ge1}\abs{\sin x/x}=\sin1$.
\end{proof}

With $c_0=-\log\sin1$ and $N=\ceil{c_0^{-1}\log(1/\eps)}$ the tail is at most $\eps$. If $\widehat\mu$ has a zero of order
$r$ at real $\lambda$ and $\mu$ is a probability measure on $[-H,H]$, then $\int P(T+t)e^{i\lambda(T+t)}\dd\mu(t)=0$ for every
polynomial $P$ of degree below $r$ and real $T$, and $\abs{\widehat\mu(\lambda+\eta)}\le\min(1,(H\abs\eta)^r/r!)$.

\begin{theorem}[Logarithmic support for Dirichlet zero filters; \stU]\label{rec:thm:logfilter}
For a fixed finite family $\mathcal F$ of primitive characters closed under conjugation, let $\Lambda_B$ be the
multiset of nonzero ordinates $\abs\gamma\le B$ of their nontrivial zeros, multiplicities added over the family.
There are even probability densities $\mu_B$, $B\ge2$, supported in $[-H_B,H_B]$ with $H_B=O_{\mathcal F}(\log(B+3))$,
vanishing on $\Lambda_B$ to at least the aggregate multiplicity, with
$\sup_{\abs\omega\ge B}\abs{\widehat\mu_B(\omega)}\le e^{-c_0B\log B}$.
\end{theorem}

\begin{proof}
The local count gives occupancy $K_B=O_{\mathcal F}(\log(B+3))$ for a fixed $\delta\le1$ below the least nonzero
absolute ordinate; take $N=\ceil{B\log B}$ in \cref{rec:thm:finitefilter}.
\end{proof}

No simplicity, distinctness across characters or independence of ordinates is used, and the
ordinates are real by definition; \cref{src:thm:dgrh} enters only when they are used as Fourier
frequencies of prime errors. The order of the support is optimal.

\begin{lemma}[Jensen support bound; \stU]\label{rec:lem:jensensupport}
If a probability measure on $[-H,H]$ has Fourier transform vanishing on a finite symmetric multiset
$\Lambda\subset[-B,B]\setminus\set0$ to the given multiplicities, then
$H\ge\frac\pi{2B}\sum_{\lambda\in\Lambda}\log\frac B{\abs\lambda}=\frac\pi B\int_0^B\frac{N_\Lambda(t)}t\dd t$, with $N_\Lambda$ the positive counting
function.
\end{lemma}

\begin{proof}
$F=\widehat\mu$ has $F(0)=1$ and $\abs{F(Re^{i\theta})}\le e^{HR\abs{\sin\theta}}$, so
Jensen's formula gives
\[
 \sum_{\abs z<R}\log(R/\abs z)\le2HR/\pi;
\]
keep the prescribed zeros and let $R\downarrow B$. Translation handles an
uncentred interval.
\end{proof}

With the Riemann--von Mangoldt formula this gives, for multiplicity-preserving cancellation of all
zeta ordinates up to $B$, $H\ge\tfrac12\log B-\tfrac12\log(2\pi)-1-O(\log^2B/B)$ (\stU), and for a family of $d$
characters $H\ge\frac d2\log B+\tfrac12\sum_\chi\log(f_\chi/2\pi)-d-O_{\mathcal F}(\log^2B/B)$. Under RH (\stG), Montgomery's
proved pair-sum formula gives $\sum_{0<\gamma\le T}m_\gamma^2\le(\tfrac43+o(1))N(T)$ over distinct ordinates
\cite{Montgomery1973Pair,ConreyGhoshGonek1998Simple}, so at least $(\tfrac56-o(1))N(T)$ ordinates are distinct, by
$m-1\le(m^2-m)/2$, and even simple cancellation at the distinct zeta ordinates needs
$H\ge(\tfrac5{12}-o(1))\log B$.

\begin{theorem}[A uniform positive averaging rule; \stU]\label{rec:thm:uniformavg}
Let $f(t)=b+\sum_{\gamma\in\Gamma}c_\gamma e^{i\gamma t}$ with $\Gamma\subset\R\setminus\set0$ locally finite and symmetric, coefficients in a
finite-dimensional normed space, $\mathcal A=\sum\norm{c_\gamma}<\infty$, and let $\mu_B$ be a probability measure with
$\widehat\mu_B=0$ on $\Gamma\cap[-B,B]$ and $\abs{\widehat\mu_B}\le e^{-c_0B\log B}$ on $\abs\omega\ge B$. Then
$\sup_t\norm{\int f(t+u)\dd\mu_B(u)-b}\le\mathcal Ae^{-c_0B\log B}$. If $H_B\le C\log B$, one rule of half-width
$O(\log\log(\mathcal A/\eps))$ has error at most $\eps$ at every centre.
\end{theorem}

\begin{proof}
The difference from $b$ is $\sum_{\abs\gamma>B}c_\gamma e^{i\gamma t}\widehat\mu_B(\gamma)$. For the second claim take
$B=2L/(c_0\log L)$, $L=\log(\mathcal A/\eps)$.
\end{proof}

If each low frequency is known only to within $\eta$ in a factor of half-width at most $A$, the bound
becomes $\mathcal A(A\eta+e^{-c_0B\log B})$, by $\abs{\phi(\omega)-\phi(\lambda)}\le A\abs{\omega-\lambda}$ for a probability measure on $[-A,A]$;
a numerical interpolation vector $\widetilde a$ with residual $r=M\widetilde a+b$ satisfies
$\norm{\widetilde a-a}_1\le\tfrac{12}{11}\norm r_1$, which certifies finite kernels from interval data.

\begin{theorem}[Simultaneous averaging on a polylogarithmic range; \stG]\label{rec:thm:polylogavg}
Fix $q$, and in the notation of \cref{zer:lem:ramp} let $R_a(x)=S_a(x)-\log x/n_q-c_a$ and $\beta_a=4v_a(0)$. There
is $C_q$ such that for large $X$ one probability measure $\mu$ on $[-C_q\log\log X,C_q\log\log X]$, independent of
$a$, satisfies for every reduced class $a$
\[
 \int e^{u/2}R_a(Xe^u)\dd\mu(u)=\frac{\beta_a}{\sqrt X}+O_q\Bigl(\frac{(\log X)^{1+C_q/2}}X\Bigr),
\]
using only scales in $[X(\log X)^{-C_q},X(\log X)^{C_q}]$. The average can be realized at each $X$ by at most
$\varphi(q)+1$ scales with nonnegative weights.
\end{theorem}

\begin{proof}
Apply \cref{rec:thm:logfilter} to the signed noncentral ordinates of the inducers modulo $q$, with $B=T=\log X$
and $H=H_T\le C_q\log T$. By \eqref{zer:eq:Zexp} and \cref{rec:thm:uniformavg}, the average of $e^{t/2}R_a(e^t)$ at
$t=T+u$ differs from $\beta_a$ by $\mathcal A_qe^{-c_0T\log T}+O_q((1+T+H)e^{-(T-H)/2})$, with $\mathcal A_q$ the absolute sum of the coefficients in \eqref{zer:eq:Zexp},, and multiplying by $X^{-1/2}$ gives
the claim. For the atoms, the integral of the continuous vector $(e^{u/2}R_a(Xe^u))_a$ lies in the convex
hull of its values, and Carath\'eodory's theorem gives at most $\varphi(q)+1$ points.
\end{proof}

This improves the common compact kernel of \cref{zer:thm:kernel}, whose support is $O_q(\log^2B)$, with scales
$X\exp(\pm O_q((\log\log X)^2))$ and remainder $O(X^{-1+\eps})$; it is a translation-independent rule and does not
compete with centre-dependent exact balancing in bounded windows. An exact rational certificate for
the clustered multiset $\set{3/2,3/2,7/4,11/4,17/5,19/5,47/10}$ and numerical zeta filters for $B\le200$ are
\begin{center}
\artifact{logarithmic_dirichlet_zero_filters_2026-11-26/evidence/exact_certificate.json}\\
\artifact{logarithmic_dirichlet_zero_filters_2026-11-26/evidence/numerical_results.json}
\end{center}
At these small cutoffs the construction is not shorter than one uniform factor per ordinate, and
the improvement is in asymptotic order.

\subsection{Open questions}\label{rec:sec:open}
\sources{dirichlet_inverse_problems; dirichlet_li_rigidity; dirichlet_zero_packet_homometry; dirichlet_zero_recovery; large_kernels_and_dirichlet_zeros; one_gap_dirichlet_zeros; prime_pole_windows; prime_power_rigidity; prime_zero_reconstruction; quadratic_spectra; three_scale_reconstruction; grh_zero_count_certificate; zeta_zero_moment_bounds; grh_periods_sampling_and_zeros; dirichlet-grh-research; grh-torsion-and-zero-data; dirichlet_low_zero_identification; dirichlet_zeros_and_abelian_fields; local_euler_recovery; robust_dirichlet_spectra_2026-11-26; dirichlet_packet_reconstruction; dirichlet_prime_frequency_inverse; prime_observations_and_progression_free_sets; translated_dirichlet_zero_divisors; logarithmic_dirichlet_zero_filters_2026-11-26}

\begin{enumerate}
\item \emph{Effective versions.} The limits in \cref{rec:thm:universal,rec:thm:factorial,rec:thm:threescale,rec:thm:primepower,rec:thm:liseparation,rec:thm:dyadic,rec:thm:mellin,rec:thm:operatorinverse}
are exact statements about infinite data. For characters of bounded conductor (and, for Li data,
bounded strides and filter degrees), what sample size, centre or separation time suffices, given
explicit zero counts and controlled measurement error, without any assumption on the spacing of
ordinates? For \cref{rec:thm:dyadic}, can an analytic weight recover all labels on the grid $2^n$ with only
exponential loss of precision in $\abs\gamma$, and how short can the prime window be? For
\cref{rec:thm:mellin}, can one fixed weight on $O(\log X)$ explicit radii retain the zero divisor of every
Galois field, and what finite heights are possible at a prescribed precision? For \cref{rec:thm:rigid}:
how accurately do finitely many noisy distances on a bounded range of lags distinguish two moduli?
For \cref{rec:thm:census} applied to $\Xi$ or to a Dirichlet completion, can explicit derivative and
zero-separation estimates bound the first working order for a disk with rational centre and radius?
\item \emph{The height needed for recovery.} How low can the height needed to recover $(q,d)$ in the
lattice class be? The bound $O(d^2\log^3(2d))$ of \cref{rec:thm:recovery} isolates one atom at a time;
probes that distinguish whole arithmetic signatures could use more of the prime data. What is the
optimal height in \cref{rec:thm:lowzero}? Its proof uses a positive proportion of distinguishing
primes at the scale $\log^2Q$ and bounds the neighbouring Euler coefficients by absolute values; can
cancellation among them give enough readings at a smaller bandwidth? How fast must the constant
grow with the number of factors in \cref{rec:cor:products}, and can the radius
$O(d^2\log^2(2Q)\log(2d))$ of \cref{rec:thm:localid} be reduced using the integrality of the multiplicity
differences? Can the field cutoff $\log^2(2D)$ of \cref{rec:thm:fields} be lowered by another $\log\log D$,
and the term $J(D,\tau)$ reduced uniformly in the discriminant? Which absolute constants make these
cutoffs practical, and how many \emph{distinct} ordinates, not merely multiplicity discrepancies, must
separate two characters or two fields at a comparable height?
\item \emph{Four central values.} Is $B_d>-2$ in \cref{rec:cor:chowla} for every prime
$d\equiv7\pmod8$? This is a subfamily of Chowla's conjecture; the one-sided estimate of the source
does not decide it.
\item \emph{Which multiplicity vectors occur?} \Cref{rec:lem:blocks,rec:thm:precision,rec:thm:detecting}
treat all abstract nonnegative integer vectors. Which vectors $(\ord_\rho L(s,\chi))_\chi$ actually occur
at common zeros, and do the worst-case near-collisions of \cref{rec:thm:precision} occur?
\item \emph{Windows and constants.} \Cref{rec:thm:localrank,rec:thm:localcoeff,rec:thm:bothsigns} reach the
radius $\pi^{-1}\log\log\log T+O(1)$. Is the coefficient $1/\pi$ optimal, and does a window of fixed radius
suffice at a fixed modulus, for two fixed characters, for two operators on the same completion
(\cref{rec:cor:operatorlocal}, at present $C\log\log T$), or for the full range of steps in
\cref{rec:thm:coalesce}? Is \cref{rec:thm:micro} sharp for the actual zeros? Can one force an ordinate at
which exactly one of several given functions vanishes, at a comparable radius? How small can the
constants be made in terms of $q$ and the family, and how much error in locations and residues is
admissible? Does positive-fraction recovery (\cref{rec:thm:stable}) hold in every window of length
$T^\theta$ for a fixed $\theta<1$, what is an explicit admissible corruption fraction already for $q=4$, and
what precision in positions suffices when close ordinates may be merged? For a nonabelian class
function, what replaces the power-image rank when Artin factors share constituents or have poles?
\item \emph{Pooled zero statistics.} Can \cref{rec:thm:homuniform} be extended to frequencies with
$e^{\abs t}\asymp T^2$? What is the smallest packet size giving a nontrivial ambiguity for each
exponent? For a fixed small mass, what are the exact measurement counts of
\cref{rec:thm:probecount} as functions of $(E,m)$, what is the least maximum order at the optimal
adaptive count for general $E$, and how small can the height constant of \cref{rec:cor:probewindow} be
made? Do products of $L$-functions of characters of odd prime order admit an equally explicit range of
local matching densities as in \cref{rec:thm:qmatch}, and can one density measurement in
\cref{rec:thm:qmatchfinite} be made at a smaller scale?
\item \emph{Omitted primes.} For primes bounded by $B$, distinct sets of at least three primes have
distinct short-gap polynomials (\cref{rec:thm:gaps}); how does the minimal separation of their
coefficient vectors depend on $B$, and what finite height and precision distinguish them? For two
omitted primes, recovery would follow from the four exponentials conjecture.
\item \emph{Local Taylor data.} For tied tail entries \cref{rec:thm:alternating} holds while the
extremal statement of \cref{rec:thm:sharpendpoint} can fail: which multiplicity patterns determine the
sharp endpoint? How fast may the degree grow with the stride while the root selection stays uniform,
and which stride and degree give the narrowest certified interval from moments with given errors?
For \cref{rec:thm:oneinterval}, how should the markers and the interval be chosen under prescribed
observation errors, and are there finite bounds like \cref{rec:thm:onegapbounds} for a complex
character, whose two signs of ordinates carry different masses?
\item \emph{Fewer filters and weights.} Can the three scales of \cref{rec:thm:threescale} be replaced by
two different nonnegative filters with an injective spectral map? Can one real nonnegative prime
weight, perhaps with higher-order correlations, encode the labelled spectrum of
\cref{rec:thm:universal}, and can its prefix precision be improved? Does the finite-family converse
\cref{rec:prop:finiteconverse} extend to the countable aggregate?
\item \emph{Other $L$-functions.} For automorphic $L$-functions of unknown degree the leading term of
the zero count measures degree times stride in \cref{rec:prop:lidecode}; what further spectral
information separates the two? For higher-degree $L$-functions, which gamma parameters can be
separated at a single prime frequency as in \cref{rec:thm:primefreq}, and which need accuracy finer
than every inverse power of $\log T$? Do the universal encoding and the positive-residue
reconstructions extend to families for which the corresponding Riemann hypothesis is known?
\Cref{src:thm:dgrh} does not supply it.
\item \emph{Shortest positive averages.} For which sets of nonconsecutive prescribed frequencies does the
shortest positive dilation average of \cref{rec:thm:fwidth} still have explicit support and shortening
constant? What is the optimal constant in the support $O(\log B)$ of \cref{rec:thm:logfilter}, between the
Jensen lower bound $\tfrac12\log B$ for $\zeta$ and the construction?
\end{enumerate}
